\chapter{BV-BRST formalism and Gaiotto's perspective}
\label{ch:v1-bv-brst}
\index{Batalin--Vilkovisky|see{BV algebra}}

The BRST complex and the chiral bar complex compute the same
derived quotient only on their common genus-$0$ surface. Costello's
BV formalism supplies the derived moduli problem of gauge fields;
the chiral bar construction supplies the collision coalgebra of the
boundary algebra. A comparison theorem must specify the surface on
which these two presentations meet.

On $X=\mathbb{P}^1$, after gauge fixing and on the uncurved relative
complex, the classical BV equation is transported to the square-zero
bar equation. The quantum master equation uses the convention
\[
\hbar \Delta_{\mathrm{BV}}S+\tfrac{1}{2}\{S,S\}_{\mathrm{BV}}=0,
\]
with the factor $\tfrac{1}{2}$ fixed as in
Appendix~\ref{app:signs}. Its genus-$0$ classical part
corresponds to $d_{\mathrm{bar}}^2=0$ under the comparison; the
Laplacian and higher-genus sewing terms require additional data.
For the bosonic string, $c_{\mathrm{matter}}=26$ is the BRST
nilpotence condition for matter plus the $bc$ ghost system, equivalently
the vanishing of the BRST conductor
$K_{\mathrm{BRST,tot}}=(c_{\mathrm{matter}}-26)/2$. This conductor is
the Virasoro half-central-charge seen by the reparametrisation ghosts;
it is not the trace-form collision characteristic
$\kappa(\mathcal H^{\oplus 26}_1)=26$. The Virasoro Koszul self-dual
point is $c = 13$, a different condition.

A natural obstruction sits in the way. At higher genus, the BV
Laplacian receives contributions from handle-gluing amplitudes that
involve the full moduli space $\overline{\mathcal{M}}_{g,n}$, while
the bar differential collects OPE residues along FM boundary strata.
What is proved in this chapter is the genus-$0$ BV/bar comparison
and the all-genera scalar Heisenberg identity, together with the
conditional all-genera coderived comparison of
Theorem~\ref{thm:bv-bar-coderived-vol1}. Beyond genus~$0$, the
obstruction analysis isolates why classes~$\mathsf{G}$,
$\mathsf{L}$, and~$\mathsf{C}$ admit stronger chain-level statements
than class~$\mathsf{M}$. The higher-genus comparison becomes an
isomorphism in $D^{\mathrm{co}}$ under the harmonic-factorization
package: the harmonic discrepancy then factors through the curvature
and the resulting cone is $m_0$-torsion, hence coacyclic. The stronger
chain-level statement is unconditional for
classes~$\mathsf{G}$ and~$\mathsf{L}$, conditional for
class~$\mathsf{C}$ on harmonic decoupling, and fails for
class~$\mathsf{M}$. The Heisenberg case is resolved at the
scalar level at all genera:
$F_g^{\mathrm{BV}}(\cH_\kappa) = F_g^{\mathrm{bar}}(\cH_\kappa)
= \kappa \cdot \lambda_g^{\mathrm{FP}}$
\textup{(}Theorem~\ref{thm:heisenberg-bv-bar-all-genera-vol1}\textup{)}.
For general algebras what survives at every genus is the scalar
anomaly package: the universal MC class $\Theta_\cA$ encodes the
one-loop anomaly coefficient $\kappa(\cA)$, and its projection
onto $H^2(\overline{\mathcal{M}}_{g,n})$ gives $F_g = \kappa(\cA) \cdot \lambda_g$
on the uniform-weight lane \textup{(}Theorem~D\textup{)}. The
multi-weight correction $\delta F_g^{\mathrm{cross}}$ is where the
chain-level identification breaks: the cross-channel term in the
sewing expansion has no counterpart in the BV side until
harmonic decoupling is assumed, and class~$\mathsf{M}$ is precisely
where that assumption fails.

\begin{remark}[BRST anomaly brackets and transferred SC operations]
\label{rem:brst-anomaly-gkw}
\index{Gaiotto--Kulp--Wu!BRST anomaly brackets}
\index{BRST operator!GKW higher operations}
The BRST-anomaly brackets of Gaiotto--Kulp--Wu~\cite{GKW2025}
and the transferred $\mathrm{SC}^{\mathrm{ch,top}}$ operations of
the bar complex describe the same deformation problem on the
genus-$0$ overlap: GKW compute higher operations from Wess--Zumino
consistency on a disk chart; the bar complex computes them from
iterated residues on FM boundary strata
(Remark~\ref{rem:iterated-residues-ainfty}).
The MC equation $D \cdot \Theta + \tfrac{1}{2}[\Theta, \Theta] = 0$
(Theorem~\ref{thm:mc2-bar-intrinsic}) is the all-genera master
equation of which GKW's consistency conditions are the genus-$0$
projections.
\end{remark}

\begin{remark}[Universal-conductor identity: BV/BRST instance;
\ClaimStatusConditional]
\label{rem:bv-brst-universal-conductor}
\index{universal conductor!BV/BRST instance|textbf}
\index{ghost charge!conductor formula}
\index{KSDual!ghost-trace alignment}
\emph{Type signature.} (Open quadrant; chain-level presentation;
Beilinson level $4$ on the chosen ghost resolution; hypothesis
package: conformal chiral algebra $\cA$ with chosen stress tensor and
chosen BRST ghost system; KSDual scope for the equality with the
trace-form conductor.)

For a conformal chiral algebra with chosen stress tensor set
\[
 K_{\mathrm{BRST}}(\cA):=\frac{1}{2}c(\cA).
\]
This is the Virasoro conductor measured by the BRST ghost current. It
agrees with the modular characteristic only in lineages where the
standard-landscape ratio gives $\kappa=c/2$; it differs from the
Heisenberg collision invariant, where
$\kappa(\mathcal H^{\oplus d}_1)=d$ while
$K_{\mathrm{BRST}}(\mathcal H^{\oplus d}_1)=d/2$. The
matter-plus-ghost cancellation
$K_{\mathrm{BRST,tot}}=0\Leftrightarrow Q_{\mathrm{BRST}}^2=0$
(Lemma~\ref{lem:brst-nilpotence}, Remark~\ref{rem:brst-nilpotence-periodicity})
is the BV/BRST instance of the ghost-resolution conductor identity
\[
 K_{\mathrm{univ}}^{\mathrm{BRST}}(\cA)
 \;=\; -\,\frac{1}{2}\,
 c_{\mathrm{ghost}}\bigl(\BRST(\cA)\bigr).
\]
The conductor $K_{\mathrm{univ}}^{\mathrm{BRST}}$ is attached to a
chosen BRST ghost resolution. It is not the ordered-collision invariant
$\kappa(\cA)$, and it is not the Theorem~C complementarity conductor
$K^\kappa(\cA,\cA^!)=\kappa(\cA)+\kappa(\cA^!)$ without an additional
comparison theorem. The ghost-trace conductor is scope-aligned to the
$\mathbb{Z}/2$-fixed sublocus $\mathrm{KSDual}\subset\mathrm{Kosz}(X)$
under the chiral/antichiral involution $\cA\mapsto\cA^!$: on KSDual
the half-normalisation $K_{\mathrm{BRST}}=\tfrac12 c$ pairs with the
Theorem~C complementarity conductor as the equivariant signature, and
the ghost-trace pairs with the bulk
$Z^{\mathrm{der}}_{\mathrm{ch}}(\cA)$ under the
ghost-resolution-to-derived-centre comparison morphism. Off KSDual the
conductor is a chosen-resolution invariant only. Chapter~\ref{chap:universal-conductor} records the
full central-charge normalisation; the present BRST comparison uses the
half-normalisation matching the Virasoro convention
$\kappa(\mathrm{Vir}_c)=c/2$. The identity records the conformal anomaly
cost that a BRST gauging must absorb.
For the bosonic string the reparametrisation ghost system has
$c_{bc}=-26$ and $K_{\mathrm{BRST}}(bc)=-13$; the free matter sector
has $c_{\mathrm{matter}}=26$ and
$K_{\mathrm{BRST}}(\mathcal H^{\oplus 26}_1)=13$, so the total BRST
conductor vanishes. The separate collision invariant remains
$\kappa(\mathcal H^{\oplus 26}_1)=26$.

For the WZW example $\cA=V_k(\fg)$ the affine ghost system contributes
$c_{\mathrm{ghost}}^{\mathrm{WZW}}=-2\dim(\fg)$
(Remark~\ref{rem:wzw-brst-bar-generic-level}), hence
$-\tfrac12c_{\mathrm{ghost}}^{\mathrm{WZW}}=\dim(\fg)$. This cancels
the Virasoro conductor of the Feigin--Frenkel pair:
\[
 \frac{1}{2}\Bigl(c(V_k(\fg))+c(V_{-k-2h^\vee}(\fg))\Bigr)
 = \dim(\fg)
 = -\tfrac{1}{2}\,c_{\mathrm{ghost}}^{\mathrm{WZW}}.
\]
The trace-form modular characteristic obeys the different identity
\[
 \kappa(V_k(\fg))+\kappa(V_{-k-2h^\vee}(\fg))=0,\qquad
 \kappa(V_k(\fg))=\frac{\dim(\fg)(k+h^\vee)}{2h^\vee},
\]
as recorded in the landscape census. Thus the ghost conductor cancels
central charge, while Verdier complementarity cancels the trace-form
curvature coefficient.
Chapter~\ref{chap:universal-conductor} is the reference for the
ghost-resolution conductor and for any theorem comparing it with the
central-charge, Euler, or ghost-number carriers under their stated
hypotheses. The cross-volume bridge identifies the genus-$0$ shadow of
the boundary-bulk ghost-cancellation formula obtained from Vol~II's
universal holographic master theorem
(Vol~II Theorem~\ref{V2-thm:universal-holography-master}, cited
across the volume boundary; cf.\ also the typed KZ--Arnold
bar-superconnection theorem of Chapter~\ref{ch:climax-theorem}).
\end{remark}

\begin{remark}[Structural cross-references]
\label{rem:bv-brst-anchors}
\index{BV/BRST!structural cross-references}
The BV/BRST chapter sits at three structural anchor
points:
\begin{enumerate}[label=\textup{(W\arabic*)}]
\item \emph{Universal conductor.} The identity
 $K_{\mathrm{univ}}^{\mathrm{BRST}}(\cA)
 =-\tfrac12c_{\mathrm{ghost}}(\BRST(\cA))$ is the ghost-resolution
 conductor functor of Chapter~\ref{chap:universal-conductor};
 Lemma~\ref{lem:brst-nilpotence} is its genus-$0$ instance. Identifying
 this carrier with the ordered-collision invariant $\kappa(\cA)$ is a
 separate comparison statement and is not automatic.
\item \emph{Bar--cobar reconstruction.} The bar/cobar counit
 $\Omega(\barB(\cA)) \to \cA$ belongs to reconstruction, not to
 BRST and not to Verdier/Koszul duality. Theorem~\ref{thm:bv-bar-geometric}
 supplies a genus-$0$ comparison between the gauge-fixed BV complex
 and the geometric bar complex; Theorem~A supplies the separate
 reconstruction map from the bar coalgebra back to~$\cA$ on the
 Koszul locus.
\item \emph{Climax identity.} The genus-$0$ BV/bar identification
transports $Q_{\mathrm{BRST}}$ to the ordered-bar differential on the
relative nilpotent locus $c_{\mathrm{matter}}=26$. This is one
face of the typed KZ--Arnold bar-superconnection theorem, together
with the corresponding BRST ghost-conductor comparison in
Chapter~\ref{ch:climax-theorem}: the BV comparison lands in the
bar-complex operator, while the KZ--Arnold superconnection produces
that same operator only after Fulton--MacPherson boundary-residue
realization.
\end{enumerate}
The cross-volume bridge to Vol~III is via the two-stage
CY-to-chiral functor
$\Phi_d^{(\Sigma_{d-1}, C)}
=\mathrm{Sp}^{\mathrm{ch}}_{\Sigma_{d-1},C}\circ\Phi_d^{\mathrm{FA}}$
at $d=3$, with specialisation data $(\Sigma_{d-1}, C)$ named
explicitly: on the five verified Borcherds families
$N \in \{1,2,3,4,6\}$, the conjectural comparison is between the
present BV/BRST conductor and the proved Borcherds weight
$\kappa_{\BKM}(X)=c_N(0)/2$ in the BPS-counting category. The
level-$0$ vertical companion is the factorisation dg-cat on
$(X, D, \tau)$ identified with Vol~III's CY-cat via 6d holomorphic
Chern--Simons under a \emph{quartic} anomaly polynomial
$\int_X\Tr_{\mathrm{ad}}\bigl(A(F_A)^3\bigr)$ on the CY$_3$ trivialised
by $\omega_X\cong\mathcal{O}_X$. The dimensional shift to
six dimensions raises the obstruction to degree~$4$ in the gauge field
$A$, distinguishing it from the cubic-Casimir anomaly of $5$d hCS; the
identification ``$6$d hCS $=$ $3$d CS in disguise'' has no type.
Companion remark: \ref{rem:bvbrst-6d-hcs-quartic}.
\end{remark}

\begin{remark}[BV/bar dictionary; \ClaimStatusHeuristic]
\label{rem:bv-bar-bridge}
\index{BV algebra!bar complex bridge|textbf}
\emph{Type signature.} (Open quadrant; chain-level on item~(i),
$(\infty, 1)$-categorical heuristic on items~(ii)--(iii); Beilinson
levels $1$--$2$ on the gauge-fixed BV/bar identification, level $4$
on the modular characteristic; hypothesis package: chiral algebra
$\cA$ on a smooth curve $X$ with chosen stress tensor and chosen
gauge fixing $\mathcal L$.)

Let $\cA$ be a chiral algebra on a smooth curve~$X$. The BV and
bar-cobar formalisms are related by the following dictionary:
\begin{enumerate}[label=\textup{(\roman*)}]
\item the geometric BV complex and the geometric bar complex are
 identified at genus~$0$
 \textup{(}Theorem~\textup{\ref{thm:bv-bar-geometric}}, proved
 elsewhere\textup{)};
\item the classical master equation corresponds to the
 genus-$0$ square-zero condition for the bar differential after
 transporting through that external comparison;
\item the scalar genus-$1$ obstruction is the modular characteristic
 $\kappa(\cA)$, the one-loop anomaly
 coefficient.
\end{enumerate}
The chain-level identifications (antifields versus
cogenerators, the BV Laplacian, the full quantum master equation,
BRST versus bar cohomology) remain heuristic in this manuscript.

Item~\textup{(i)} is Theorem~\ref{thm:bv-bar-geometric}. The scalar
genus-$1$ coefficient in item~\textup{(iii)} is the content of
Theorem~\ref{thm:genus-universality} and, for Virasoro, of
Computation~\ref{comp:virasoro-curvature}. The remaining entries are
the intended BV reading of these proved algebraic statements, not
chain-level equalities.
\end{remark}

\begin{remark}[Heuristic BV reading of the modular MC hierarchy]
\label{rem:modular-qme-bv}
\index{quantum master equation!modular}
\index{modular quantum master equation|textbf}

The proved object is the universal Maurer--Cartan class~$\Theta_\cA$
\textup{(}Theorem~\ref{thm:mc2-bar-intrinsic}\textup{)}.
Write $\Theta_\cA = \Theta_0 + \hbar\,\Theta_1 + \hbar^2\,\Theta_2
+ \cdots$ for the genus expansion. The BV-shaped modular equation
\begin{equation}\label{eq:modular-qme-bv}
d\,\Theta_\cA + \tfrac{1}{2}[\Theta_\cA, \Theta_\cA]
+ \hbar\,\Delta\,\Theta_\cA = 0
\end{equation}
in the completed cyclic deformation Lie algebra
$\Defcyc^{\mathrm{mod}}(\cA)$ breaks into a hierarchy:
\begin{alignat*}{2}
&\text{genus } 0\colon\quad
 &&d\,\Theta_0 + \tfrac{1}{2}[\Theta_0,\Theta_0] = 0
 \qquad\text{(classical ME $=$ $d_{\mathrm{bar}}^2 = 0$)}, \\
&\text{genus } 1\colon\quad
 &&d\,\Theta_1 + [\Theta_0,\Theta_1] + \Delta\,\Theta_0 = 0
 \qquad\text{(one-loop anomaly $=$ $\kappa(\cA)\cdot\omega_1$)}, \\
&\text{genus } 2\colon\quad
 &&d\,\Theta_2 + [\Theta_0,\Theta_2]
 + \tfrac{1}{2}[\Theta_1,\Theta_1]
 + \Delta\,\Theta_1 = 0
 \qquad\text{(two-loop interference)}.
\end{alignat*}
The displayed hierarchy is a heuristic BV translation of the proved
modular MC hierarchy, not a proved BV identity: the
genus-$0$ line is a classical master equation, the genus-$1$
line a one-loop anomaly equation, and higher lines encode
nonlinear modular interference. At the scalar level, the trace
of~\eqref{eq:modular-qme-bv} recovers the genus tower
$F_g = \kappa(\cA)\lambda_g$ on the uniform-weight lane by
Theorem~D. At the
full level, no BV identity is asserted here; equation~\eqref{eq:modular-qme-bv}
is the conjectural target of the modular cumulant transform
(\S\ref{subsec:completion-kinematics-theory} of the concordance).

In the Feynman-diagram language of
Chapter~\ref{ch:v1-feynman}: $\Theta_0$ collects tree
diagrams, $\Theta_1$ collects one-loop diagrams, and $\Theta_g$
collects connected $g$-loop diagrams. In physics language, the modular MC hierarchy is an all-loop Ward identity. The connected
stable-graph exponential
$\Theta_\cA = \sum_{\Gamma\,\text{connected}}
\frac{\hbar^{g(\Gamma)}}{|\operatorname{Aut}(\Gamma)|}\,\Phi_\Gamma$
makes this precise: the vertices are cumulant primitives, internal
edges are factorization propagators, and the symmetry factor is the
standard graph automorphism weight.
\end{remark}

\begin{remark}[String field theory interpretation; \ClaimStatusHeuristic]
\label{rem:sft-bar-identification}
\index{string field theory!bar complex identification}
\index{bar complex!string field theory reading}
In the heuristic string-field-theory reading, the genus-$0$
bar--BRST comparison sends the bar differential
$d_{\barB}$ to the BRST charge~$Q$ of open string field theory. The
bar curvature~$m_0$ then plays the role of the tadpole, the shadow
amplitude $\mathrm{Sh}_{g,n}(\Theta_\cA)$ plays the role of the
closed SFT vertex at genus~$g$ with $n$~punctures, and the MC
equation is the algebraic analogue of the classical closed SFT
master equation.
\end{remark}

\section{BV formalism for chiral algebras}

\begin{definition}[BV data for a chiral algebra]
\label{def:bv-data-chiral}
\index{antibracket}
\index{BV algebra!data}
Let $\cA$ be a chiral algebra on a smooth curve~$X$.
The BV formalism associates to~$\cA$:
\begin{enumerate}[label=\textup{(\roman*)}]
\item \emph{Fields}: $\phi \in \mathcal{A}$, sections of the chiral algebra sheaf.
\item \emph{Antifields}: $\phi^+ \in \mathcal{A}^*[1]$, the Serre dual shifted by~$1$.
\item \emph{BV bracket}: the degree-$+1$ odd Poisson bracket
 $\{-,-\}_{\mathrm{BV}}$ generated by the shifted Serre pairing,
 with graded skew sign
 $\{F,G\}_{\mathrm{BV}}
 = -(-1)^{(|F|+1)(|G|+1)}\{G,F\}_{\mathrm{BV}}$.
\item \emph{Action}: $S[\phi, \phi^+]$ satisfying the classical master
 equation $\{S,S\}_{\mathrm{BV}} = 0$.
\end{enumerate}
\end{definition}

\begin{theorem}[Genus-$0$ BV complex and geometric bar complex;
\ClaimStatusProvedElsewhere{} \cite{CG17}]
\label{thm:bv-bar-geometric}
\index{BV algebra!bar complex identification}
\emph{Type signature.} (Open quadrant; chain-level presentation;
Beilinson levels $1\to 2$, the bar-coalgebra direction; hypothesis
package: chiral algebra on $X=\mathbb{P}^1$, BV presentation with
local observables recovering $\cA$, genus-$0$ uncurved surface.)

Let $\cA$ be a chiral algebra on $X = \mathbb{P}^1$, equipped with a
BV presentation whose local observables recover~$\cA$. On the
genus-$0$ uncurved surface there is a quasi-isomorphism
\[
C_{\mathrm{BV}}(\mathcal{A}) \xrightarrow{\;\sim\;}
\barB^{\mathrm{ch}}(\mathcal{A}).
\]
Under this comparison, the classical BV differential
$Q_{\mathrm{BV}}=\{S,-\}_{\mathrm{BV}}$ is transported to the
genus-$0$ bar differential $d_{\mathrm{bar}}$. The comparison is the
Costello--Gwilliam BV/factorization-observable formalism specialized
to the chiral bar model; it is not a statement about the full QME or
about bar-cobar reconstruction.
\end{theorem}

\begin{proof}[Geometric construction supporting \textup{Theorem~\ref{thm:bv-bar-geometric}}]
\emph{Step~1: Bar generators.}

The degree-$n$ bar complex is:
\[\bar{B}^n(\mathcal{A}) = \Omega^*(\overline{C}_{n+1}(X), \mathcal{A}^{\boxtimes (n+1)})\]

The generators are operator insertions $\phi_i \in \mathcal{A}$ at marked points, paired with logarithmic $1$-forms $\eta_{ij} = d\log(z_i - z_j)$ dual to the collision divisors.

\emph{Step~2: BV bracket and signs.}

The BV bracket is a residue operation on the log-de~Rham complex
of~$\overline{C}_{n+1}(X)$, with the cohomological sign convention
of Appendix~\ref{app:signs}: for homogeneous functionals,
\[
\{F,G\}_{\mathrm{BV}}
=-(-1)^{(|F|+1)(|G|+1)}\{G,F\}_{\mathrm{BV}}.
\]
On desuspended bar generators the Koszul sign is the sign obtained by
moving $s^{-1}\phi$ past the intervening logarithmic form factors.
In local coordinates near the divisor
$D_{ij} = \{z_i = z_j\}$, with $\epsilon = z_i - z_j$, the
degree-zero residue component is:
\[\{\phi(z_i), \eta_{jk}\}
= \delta_{ij}\,\Res_{\epsilon = 0}
\Bigl[\frac{\partial\phi}{\partial z_i}\,
\frac{d\epsilon}{\epsilon(z_i - z_k)}\Bigr]
+ \delta_{ik}\,\Res_{\epsilon = 0}
\Bigl[\frac{\partial\phi}{\partial z_i}\,
\frac{d\epsilon}{\epsilon(z_i - z_j)}\Bigr].\]
Both terms are Poincar\'e residues of logarithmic forms along
normal-crossing divisors of the smooth
variety~$\overline{C}_{n+1}(X)$: the residue map
$\Res_{D_{ij}}\colon \Omega^k_{\log}(\overline{C}_{n+1})
\to \Omega^{k-1}_{D_{ij}}$ is an algebraic morphism of sheaves
on a smooth variety. The factor $1/(z_i - z_k)$ is a pole along
the transverse divisor $D_{ik}$ (Theorem~\ref{thm:FM}(3):
$D_{ij} \cap D_{ik}$ has normal crossings), so on~$D_{ij}$
it restricts to a regular meromorphic function with poles only
along~$D_{ij} \cap D_{ik}$, a codimension-$2$ stratum.
The FM compactification resolves all coincidence singularities
into normal-crossing divisors, and the log-de~Rham complex on
this resolution is the correct algebraic framework: it replaces
distributional regularization with algebraic residues.

The logarithmic $1$-forms $\eta_{ij} = d\log(z_i - z_j)$ generate
$H^1(C_n(X); \mathbb{C})$ by the Arnold--Brieskorn theorem;
for $X = \mathbb{C}$, all $\binom{n}{2}$ classes are linearly independent in degree~$1$
(the Arnold relations constrain products in degree~$\geq 2$), giving
$\dim H^1(C_n(\mathbb{C})) = \binom{n}{2}$. (For a curve of genus~$g > 0$, additional generators from $H^1(X)$ contribute $2gn$ classes.)

\emph{Step~3: Master equation.}

The classical master equation $\{S,S\}_{\mathrm{BV}}=0$ is transported
to $\dzero^2=0$ for the genus-$0$ bar differential:
\[d = d_{\text{strat}} + d_{\text{int}} + d_{\text{res}}\]
\end{proof}

\begin{remark}[Literature confirmation: Costello--Gwilliam]
The genus-$0$ BV/bar identification used in this chapter is consistent
with the Costello--Gwilliam BV formalism. In the CG framework, a
perturbative quantum field theory determines a factorization algebra of
observables, and for holomorphic theories this is the source of the
vertex-algebra structure treated in Chapter~5 of~\cite{CG17}.
Theorem~\ref{thm:bv-bar-geometric} is the chiral-bar model of that same
local observable package on $X=\mathbb{P}^1$: the BV differential is
the genus-$0$ bar differential, and collision residues are the
factorization products written in logarithmic-form language.

The overlap with~\cite{CG17} is the genus-$0$ local comparison and the
general BV/factorization dictionary. The bar-side genus expansion,
harmonic comparison, and coderived package are separate inputs; the
all-genera coderived theorem
\textup{(}Theorem~\ref{thm:bv-bar-coderived-vol1}\textup{)} belongs to
this chapter's construction, not to the theorem package attributed to CG.
\end{remark}

\subsection{Quantum master equation}

\begin{definition}[BV Laplacian]
\label{def:bv-laplacian}
\index{BV Laplacian|textbf}
The BV Laplacian $\Delta_{\mathrm{BV}}$ is the second-order differential operator
\[
\Delta_{\mathrm{BV}} = \sum_i \frac{\partial^2}{\partial \phi_i \partial \phi_i^+},
\]
contracting each field-antifield pair $(\phi_i, \phi_i^+)$.
In the geometric realization on $\barB^{\mathrm{ch}}(\cA)$, the heuristic formula
\[
\Delta_{\mathrm{BV}} \approx \sum_{i<j} \int \delta(z_i - z_j)\,
\frac{\partial}{\partial \eta_{ij}}
\]
inserts a diagonal residue, contracting the logarithmic form $\eta_{ij}$ via a
delta-function along $D_{ij}$. This is the loop-insertion operator in the bar
complex, distinct from the cobar functor $\Omega$ (which recovers $\cA$ from
$B(\cA)$ by bar-cobar inversion, Theorem~\ref{thm:bar-cobar-isomorphism-main}).
The precise algebraic version of this operator is conditional
\textup{(}Theorem~\ref{thm:config-space-bv}\textup{)}.
\end{definition}

\begin{remark}[Heuristic QME/MC comparison; \ClaimStatusHeuristic]
\label{rem:qme-bar-cobar}
\index{quantum master equation|textbf}
At genus~$0$, the classical part of the BV master equation matches the
square-zero bar differential after transporting through the external
comparison of Theorem~\ref{thm:bv-bar-geometric}. Extending this to a
full identification of the quantum master equation
\[
\hbar\,\Delta_{\mathrm{BV}}S + \tfrac{1}{2}\{S,S\}=0
\]
with a Maurer--Cartan equation in a BV-enhanced bar deformation
complex would require the conditional BV package recorded in
Theorems~\ref{thm:config-space-bv} and~\ref{thm:bv-functor}. That
full comparison is not proved here.

The classical genus-$0$ comparison combines the external BV/bar
identification with the square-zero bar differential on the
uncurved locus. The quantum, Laplacian, and functorial parts are
isolated as conditional content later in this chapter.
\end{remark}

\begin{proposition}[Finite-type PVA BV/QME gate;
\ClaimStatusConditional]
\label{prop:bvbrst-finite-type-pva-qme-gate}
\index{Poisson vertex algebra!BV action gate}
\index{quantum master equation!finite-type PVA gate}
\textup{[}Type signature:
$(\mathrm{Open/HT},\mathrm{BV\ local\ functional},
\mathrm{levels}=1\to3,
\mathrm{hypothesis\ package}=\mathrm{finite\text{-}jet\ PVA}
\ +\ \mathrm{Khan\text{--}Zeng\ half\text{-}space\ action}
\ +\ \mathrm{analytic\ SDR/QME\ package})$\textup{].}
Let $\mathcal P$ be a freely generated finite-jet Poisson vertex
algebra with Hamiltonian PVA bivector $\Pi$.  The Khan--Zeng
half-space construction supplies a classical BV action
\begin{equation}
\label{eq:bvbrst-pva-classical-action}
S_{\mathcal P}^{\mathrm{cl}}
=
\int_{\RR_t\times\CC_z}
\left\langle \phi^+,(dt+\bar\partial)\phi\right\rangle
\;+\;
\int_{\RR_t\times\CC_z}\Pi(\phi,\phi^+,\partial_z)
\end{equation}
in the local functional complex determined by the PVA generators,
their antifields, and the field-parity signs.  In this finite-type
classical setting,
\begin{equation}
\label{eq:bvbrst-cme-pva-hamiltonian}
\{S_{\mathcal P}^{\mathrm{cl}},
S_{\mathcal P}^{\mathrm{cl}}\}_{\mathrm{BV}}=0
\quad\Longleftrightarrow\quad
[\Pi,\Pi]_{\mathrm{PVA}}=0,
\end{equation}
so the classical master equation is equivalent to the PVA Hamiltonian
condition \textup{(}the $\lambda$-Jacobi identity with
sesquilinearity and Leibniz rules\textup{)}.

The all-loop quantum master equation is not a consequence of
\eqref{eq:bvbrst-cme-pva-hamiltonian}.  To state
\[
\hbar\,\Delta_{\mathrm{BV}}S+\tfrac12\{S,S\}_{\mathrm{BV}}=0
\]
for the quantized theory one must also provide:
\begin{enumerate}[label=\textup{(\roman*)}]
\item an analytic strong deformation retract for the propagator
complex;
\item a Stokes image choice for boundary strata and self-image
renormalisation for coincident fields;
\item field-parity signs compatible with the BV Laplacian and the
PVA parity;
\item a Costello-style renormalisation package proving that the
counterterms are finite and local in the chosen completion.
\end{enumerate}
Only after these data are supplied may one assert an all-loop QME,
boundary vertex algebra, boundary OPE formula, or reflected-weight
identity.  The reflected weight identity is therefore
\begin{equation}
\label{eq:bvbrst-reflected-weight-gate}
K_{\mathrm{bdry}}^{!}=K_{\mathrm{slab}}-K_{\mathrm{bdry}}
\end{equation}
only in a slab or boundary model for which the slab reduction,
renormalised boundary OPEs, and OCA comparison have been constructed.
\end{proposition}

\begin{proof}
Khan--Zeng construct the half-space local functional from the finite
PVA bracket by the standard BV pairing between fields and antifields.
Under this construction the BV bracket of the interaction functional
with itself is the variational Schouten bracket of the PVA bivector.
Thus the classical master equation is precisely the condition
$[\Pi,\Pi]_{\mathrm{PVA}}=0$, i.e. the PVA Hamiltonian condition
including the $\lambda$-Jacobi identity, sesquilinearity, and Leibniz
rules.  The remaining assertions are exclusions: the BV Laplacian,
Stokes image, self-contraction counterterms, and field-parity choices
enter only in the quantum term
$\hbar\Delta_{\mathrm{BV}}S$.  They are not present in the classical
Schouten calculation, so no all-loop QME or boundary OPE theorem
follows until those analytic data are added.
\end{proof}

\label{subsec:bv-bar-identification}
\index{BV algebra!bar complex identification}

\begin{remark}[Heuristic genus-\texorpdfstring{$0$}{0} amplitude/bar comparison;
\ClaimStatusHeuristic]\label{rem:genus0-amplitude-bar}
\index{string amplitude!bar complex}
Let $\cA$ be a chiral algebra on $\mathbb{P}^1$, and let
$\omega_1, \ldots, \omega_n$ be cocycles in
$\bar{B}^1(\cA)$. The heuristic genus-$0$ string-amplitude pairing
\[
\mathcal{A}_0(\omega_1, \ldots, \omega_n)
\;=\;
\int_{\overline{\mathcal{M}}_{0,n}}
 \omega_1 \wedge \cdots \wedge \omega_n
\]
parallels the bar-complex Euler characteristic
$\chi(\bar{B}^n_0(\cA))$, where $\bar{B}^n_0$ denotes the
genus-$0$ component of the degree-$n$ bar complex. The exact
identification between these two quantities is not proved in the
manuscript.
The bar complex is realized on
$\overline{C}_n(\mathbb{P}^1) = \overline{\mathcal{M}}_{0,n+1}$
(Theorem~\ref{thm:bar-cobar-isomorphism-main}),
with the bar differential given by residues of the FM propagator
$\eta_{ij} = d\log(z_i - z_j)$ along collision divisors.
The displayed comparison is the expected physics reading of these
configuration-space forms, but the chapter does not prove the required
identification between the physical integration pairing and the Euler
characteristic of the genus-$0$ bar complex.
\end{remark}

\begin{remark}[Genus-$1$ partition function]\label{rem:genus1-bv}
At genus~$1$, the $\hat{A}$-genus gives $F_1 = \kappa(\cA)/24$
\textup{(}Theorem~D\textup{)}. For $\cA = \mathrm{Vir}_c$ with
$\kappa(\mathrm{Vir}_c) = c/2$: $F_1 = c/48$.
\end{remark}

\begin{remark}[Anomaly as curvature]\label{rem:anomaly-curvature-bv}
\index{anomaly!curvature identification}
By Theorem~\ref{thm:anomaly-koszul}, the genus-$0$ uncurved bar
model of
$\cA_{\mathrm{tot}}=\cA_{\mathrm{matter}}\otimes\cA_{\mathrm{ghost}}$
is square-zero on the BRST Virasoro lane if and only if
$K_{\mathrm{BRST,tot}}=0$; for the bosonic string this is the
condition $c_{\mathrm{matter}}=26$ after adding the $bc$ ghost system
of central charge $-26$. Outside this Virasoro lane the curvature
coefficient is the lineage-specific $\kappa(\cA)$ of the landscape
census. When that curvature coefficient is nonzero, the scalar
projection of the universal MC class is
$\Theta_\cA^{\min}=\kappa(\cA)\cdot\eta\otimes\Lambda$
\textup{(}Theorem~\ref{thm:explicit-theta}\textup{)}, and the
higher-genus fiberwise bar model is curved by this class.
\end{remark}

\section{Gauge fixing and BRST}

\subsection{BRST from BV}

\begin{definition}[BRST operator]
\label{def:brst-operator}
\index{BRST cohomology|textbf}
\index{BRST operator|textbf}
The BRST operator $Q_{\mathrm{BRST}}$ arises from gauge fixing the BV action.
Choose a Lagrangian submanifold $\mathcal{L}$ of the total space of fields and
antifields:
\[
Q_{\mathrm{BRST}} = Q_{\mathrm{BV}}\big|_{\mathcal{L}}.
\]
In the chiral algebra context, the BRST charge decomposes by antighost number
(not ghost number):
\[
Q_{\mathrm{BRST}} = Q_0 + Q_1 + Q_2 + \cdots,
\]
where $Q_k$ has antighost number~$k$ and operator dimension~$k-1$.
The total BRST charge has definite ghost number~$+1$.
\end{definition}

\begin{theorem}[BRST cohomology on a nilpotent gauge-fixed complex;
\ClaimStatusProvedElsewhere{} \cite{CG17,Polchinski1998}]
\label{thm:brst-physical-states}
Let a BV theory be equipped with a gauge-fixing Lagrangian
$\mathcal L$ such that the restricted BRST operator
$Q_{\mathrm{BRST}}$ is square-zero on the relative state complex and
all residual gauge symmetries are generated by this operator. Then
the physical on-shell state space is the BRST cohomology
\[
H^*(Q_{\mathrm{BRST}})\cong \mathcal A_{\mathrm{phys}}.
\]
\end{theorem}

\begin{example}[Free \texorpdfstring{$bc$}{bc} ghost system]
The $bc$ system has fields $b(z)$ of weight $\lambda$ and $c(z)$ of weight $1-\lambda$, OPE $b(z)c(w) \sim (z-w)^{-1}$, central charge $c_{bc} = 1 - 3(2\lambda - 1)^2$, and BRST operator $Q = \oint c(z) \bigl(T_{\text{matter}}(z) + \tfrac{1}{2}T_{\text{ghost}}(z)\bigr) dz$.

For the bosonic string, $\lambda = 2$ (so $b$ has weight~$2$ and $c$ has weight~$-1$), giving $c_{\text{ghost}} = 1 - 3(3)^2 = -26$ and ghost stress tensor $T_{\text{ghost}} = -2b\,\partial c - (\partial b)\,c$.

On the relative complex $\ker(b_0) \cap \ker(L_0^{\mathrm{tot}})$,
with the standard intercept normalisation, the BRST nilpotence
condition is:
\[
Q_{\mathrm{BRST}}^2 = 0
\iff c_{\mathrm{matter}} + c_{\mathrm{ghost}} = 0,
\quad\text{i.e., } c_{\mathrm{matter}} = 26
\]
(on the full complex, the anomaly is the quadratic ghost operator
of Lemma~\ref{lem:brst-nilpotence}, not a $c_0$ term).

In the bar complex, the residue component $d_{\mathrm{res}}$
of the bar differential heuristically corresponds to $Q_{\mathrm{BRST}}$
under the identification of Theorem~\ref{thm:bv-bar-geometric}.
\end{example}

\subsection{Coupling to topological gravity}

\begin{theorem}[Coordinate cocycle for collision logarithmic forms;
\ClaimStatusProvedHere]
\label{thm:log-form-ghost-law}
\index{ghost!transformation law}
\index{logarithmic form!ghost transformation}
The logarithmic forms $\eta_{ij} = d\log(z_i - z_j)$
on the configuration space $C_n(X)$ transform under a coordinate
change $z\mapsto w(z)$ by the exact logarithmic cocycle
\[
\eta_{ij} \;\mapsto\; \eta_{ij}
+ d\log\!\left(\frac{w(z_i) - w(z_j)}{z_i - z_j}\right).
\]
Near the diagonal $z_i \to z_j$, the correction term
approaches $d\log w'(z_i)$. Thus the FM propagator carries the
same coordinate-change cocycle that the Faddeev--Popov
diffeomorphism ghost detects after linearisation, but the
logarithmic form is an even de~Rham class, not the odd BRST ghost
field itself. The Arnold relations
$\eta_{ij} \wedge \eta_{jk} + \eta_{jk} \wedge \eta_{ki}
+ \eta_{ki} \wedge \eta_{ij} = 0$
in the Orlik--Solomon algebra provide the quadratic relation on
collision residues used by the genus-$0$ bar differential.
\end{theorem}

\begin{proof}
Direct computation: $\eta_{ij} = d\log(z_i - z_j)$
maps to $d\log(w(z_i) - w(z_j))$ under $z \to w(z)$.
Factoring:
\[
d\log(w(z_i) - w(z_j)) = d\log(z_i - z_j)
+ d\log\!\left(\frac{w(z_i) - w(z_j)}{z_i - z_j}\right)
\]
and near the diagonal, L'H\^opital gives
$(w(z_i) - w(z_j))/(z_i - z_j) \to w'(z_i)$, so
the correction approaches $d\log w'(z_i)$. Infinitesimally, for
$w(z)=z+\epsilon v(z)$,
\[
\delta_v\eta_{ij}
=d\!\left(\frac{v(z_i)-v(z_j)}{z_i-z_j}\right),
\]
whose diagonal limit is $d(v'(z_i))$. This is the coordinate
cocycle that enters the Faddeev--Popov ghost transformation, while
the parity and complex are different. The Arnold relations are
verified in Theorem~\ref{thm:arnold-relations}.
\end{proof}

\begin{conjecture}[Topological BRST comparison for collision ghosts;
\ClaimStatusConjectured]
\label{conj:v1-bar-topological-brst}
Let $\mathrm{Vect}(X)$ act on a chiral algebra~$\mathcal A$ by
conformal vector fields, and choose a Faddeev--Popov resolution of
this action. The coordinate cocycle of
Theorem~\ref{thm:log-form-ghost-law} is conjectured to extend to a filtered
genus-$0$ comparison map
\[
\Phi_{\mathrm{top}}\colon
C^*_{\mathrm{top}\text{-}\mathrm{BRST}}
(\mathcal A\otimes \mathrm{Vect}(X))
\longrightarrow
\bar{B}^{\mathrm{ch}}_0(\mathcal A),
\]
whose associated graded sends a ghost for the collision vector field
$z_i-z_j$ to the logarithmic form~$\eta_{ij}$. The conjectured
quasi-isomorphism is restricted to the nilpotent genus-$0$ relative
locus and requires the PBW lifting hypotheses of
Theorem~\ref{thm:brst-bar-genus0}; no all-genera dg-algebra
isomorphism is asserted.
\end{conjecture}

\section{The genus-\texorpdfstring{$0$}{0} BRST-bar chain map}
\label{sec:brst-bar-chain-map}
\index{BRST!bar complex chain map|textbf}

\subsection{Setup: the BRST complex}

Let $\cA$ be a conformal vertex algebra of central charge~$c$ on
$X = \mathbb{P}^1$, and let $\cA_{\mathrm{gh}} = bc$ denote the
$bc$-ghost system of weights $(2, -1)$ with $c_{\mathrm{gh}} = -26$.
The \emph{BRST complex} of~$\cA$ is the vertex algebra
$\cA_{\mathrm{tot}} = \cA \otimes \cA_{\mathrm{gh}}$ equipped with
the BRST differential:
\begin{equation}\label{eq:brst-differential}
Q_{\mathrm{BRST}}
= \oint \! c(z)\Bigl(T_{\cA}(z) + \tfrac{1}{2}\,T_{\mathrm{gh}}(z)\Bigr)\,dz
= \sum_{n \in \mathbb{Z}} \Bigl(
 \sum_m L^{\cA}_m\, c_{n-m}
 + \tfrac{1}{2}\sum_{p+q=n}(p-q)\, {:}b_p\, c_q\, c_{n-p-q}{:}
\Bigr)
\end{equation}
where $L^{\cA}_n$ are the Virasoro modes of~$\cA$ and $b_n, c_n$ are
the ghost modes with $\{b_m, c_n\} = \delta_{m+n,0}$. The
normal-ordering intercept is fixed to $a=1$; equivalently, in the
presentation
\[
Q_{\mathrm{BRST}}
=\sum_n c_{-n}(L_n^\cA+\tfrac12L_n^{\mathrm{gh}}-a\delta_{n,0})
\]
one takes
$a=1$. This is the convention used in
Lemma~\ref{lem:brst-nilpotence}.

\begin{lemma}[BRST nilpotence; \ClaimStatusProvedElsewhere{} \cite{FGZ86}]
\label{lem:brst-nilpotence}
Let $Q_{\mathrm{BRST}}$ be the bosonic-string BRST charge with
normal-ordering intercept $a=1$. As operators on the full
matter-plus-ghost Fock complex,
\[
 Q_{\mathrm{BRST}}^2
 =
 \frac{c_{\mathrm{matter}}-26}{12}
 \sum_{n\in\mathbb Z}(n^3-n)\,{:}c_{-n}c_n{:}.
\]
In particular, $Q_{\mathrm{BRST}}^2 = 0$ on the relative complex
$\cA_{\mathrm{tot}}^{\mathrm{rel}} = \ker(b_0) \cap \ker(L_0^{\mathrm{tot}})$
as an operator if and only if $c_{\mathrm{matter}}=26$.
\end{lemma}

\begin{remark}[BRST nilpotence and the bar construction]
\label{rem:brst-nilpotence-periodicity}
\index{nilpotence-periodicity correspondence!BRST instance}
The condition $Q_{\mathrm{BRST}}^2=0$ is the BRST instance of the
nilpotence-periodicity correspondence
(Remark~\ref{rem:nilpotence-periodicity}). The numerical invariant in
this sentence is the BRST conductor
\[
 K_{\mathrm{BRST,tot}}
 =\frac{c_{\mathrm{matter}}+c_{\mathrm{ghost}}}{2}
 =\frac{c-26}{2},
\]
not the bare trace-form characteristic of the underlying chiral
algebra. Vanishing of $K_{\mathrm{BRST,tot}}$ makes the relative
genus-$0$ bar model square-zero on the chosen stress-tensor lane. When
$K_{\mathrm{BRST,tot}}\neq0$, the quadratic ghost anomaly
$\frac{c_{\mathrm{matter}}-26}{12}
\sum_n(n^3-n){:}c_{-n}c_n{:}$ is the BRST curvature
obstructing nilpotence; the $n=0$ term vanishes because
$n^3-n=0$.
\end{remark}

For the bosonic string at $c_{\mathrm{matter}}=26$, the relative BRST
cohomology of $(\cA_{\mathrm{tot}},Q_{\mathrm{BRST}})$ is the
semi-infinite cohomology of the Virasoro action on the matter module,
graded by ghost number. This is the special Virasoro/string
realisation of the general semi-infinite construction defined in
\S\ref{sec:semi-infinite-bar}; it is not the definition of
semi-infinite cohomology for an arbitrary chiral algebra.

\subsection{The chain map}

\begin{theorem}[Genus-\texorpdfstring{$0$}{0} BRST-bar quasi-isomorphism;
\ClaimStatusConditional]
\label{thm:brst-bar-genus0}
\index{BRST!bar complex quasi-isomorphism}
Let $\cA$ be a PBW conformal vertex algebra on $\mathbb{P}^1$ with
$c(\cA) = 26$ \textup{(}so that
$\cA_{\mathrm{tot}} = \cA \otimes bc$ satisfies
$K_{\mathrm{BRST,tot}} = 0$\textup{)}. Assume the conformal-weight
filtrations on the relative BRST complex and on the genus-$0$ chiral
bar complex are complete and exhaustive, that their associated graded
complexes are identified by the classical bar--CE comparison, and that
the filtered lifting obstruction groups vanish. Then there exists a
chain map
\[
\Phi \colon
 \bigl(\cA_{\mathrm{tot}}^{\mathrm{rel}},\, Q_{\mathrm{BRST}}\bigr)
 \longrightarrow
 \bigl(\barB^{\mathrm{ch}}_0(\cA_{\mathrm{tot}}),\, d_{\mathrm{bar}}\bigr)
\]
from the relative BRST complex to the genus-$0$ bar complex, which is a
quasi-isomorphism of cochain complexes.
\end{theorem}

\begin{proof}
The hypotheses give a PBW filtration on both complexes, reducing the
comparison to the classical finite-dimensional setting at the
associated-graded level.

\emph{Step~1: Conformal weight filtration.}
Both complexes carry a $\mathbb{Z}_{\geq 0}$-filtration by
\emph{total conformal weight}:
\begin{align}
F^{\leq N}\cA_{\mathrm{tot}}^{\mathrm{rel}}
 &= \bigoplus_{\substack{\text{states of total}\\\text{weight } \leq N}}
 \cA_{\mathrm{tot}}^{\mathrm{rel}}, \\
F^{\leq N}\barB^{\mathrm{ch}}_0(\cA_{\mathrm{tot}})
 &= \bigoplus_{\substack{n \geq 0}}
 \Omega^{\leq N}\!\bigl(
 \overline{C}_{n+1}(\mathbb{P}^1),\,
 \cA_{\mathrm{tot}}^{\boxtimes(n+1)}
 \bigr),
\end{align}
where the filtration on bar generators counts the sum of conformal
weights of inserted operators plus the form degree on the
configuration space. Both differentials respect this filtration:
$Q_{\mathrm{BRST}}$ preserves conformal weight (the BRST current
has weight~$1$ and the contour integral lowers by~$1$), and the bar
differential $d_{\mathrm{bar}}$ is computed by OPE residues at
collision divisors, which preserve the total weight.

\emph{Step~2: Associated graded.}
On the associated graded $\mathrm{gr}\,F^{\bullet}$, both complexes
degenerate to classical constructions. For
$\mathrm{gr}\,\cA_{\mathrm{tot}}^{\mathrm{rel}}$: only the
weight-$0$ part of $Q_{\mathrm{BRST}}$ survives, which is the
Chevalley--Eilenberg differential of the finite-dimensional Lie algebra
$\mathfrak{g}_0 = \cA_{(1)}$ (the weight-$1$ subspace, viewed as a Lie
algebra via the zero-mode bracket $[a, b] = a_{(0)} b$).
For $\mathrm{gr}\,\barB^{\mathrm{ch}}_0(\cA_{\mathrm{tot}})$: only
the leading OPE pole contributes to the bar differential, giving
the classical bar complex $\barB(\mathfrak{g}_0)$ of the
finite-dimensional Lie algebra~$\mathfrak{g}_0$.

In both cases, the associated graded complex is the
Chevalley--Eilenberg / bar complex of a finite-dimensional Lie algebra.
By Arkhipov~\cite{Arkhipov97}, there is a natural quasi-isomorphism
\begin{equation}\label{eq:arkhipov-classical}
\Phi_0 \colon C^{\bullet}_{\mathrm{CE}}(\mathfrak{g}_0)
 \xrightarrow{\;\sim\;}
 \barB(\mathfrak{g}_0)
\end{equation}
given by the bar-CE comparison map: on generators,
$\Phi_0(x_1 \wedge \cdots \wedge x_n)
= [x_1 | \cdots | x_n]_{\mathrm{bar}}$
(the corresponding bar element), extended by the shuffle product
to all degrees. This is a classical result in homological algebra
(e.g., Loday--Vallette~\cite{LV12}, \S2.2): the bar complex and
the Chevalley--Eilenberg complex of a Lie algebra are
quasi-isomorphic via the antisymmetrization map.

\emph{Step~3: Spectral sequence comparison.}
The conformal weight filtration induces convergent spectral sequences:
\begin{align}
{}^{\mathrm{BRST}}E_1^{p,q}
 &= H^{p+q}\bigl(\mathrm{gr}^p\, \cA_{\mathrm{tot}}^{\mathrm{rel}}\bigr)
 \;\Longrightarrow\;
 H^{p+q}(Q_{\mathrm{BRST}}), \\
{}^{\mathrm{bar}}E_1^{p,q}
 &= H^{p+q}\bigl(\mathrm{gr}^p\, \barB^{\mathrm{ch}}_0(\cA_{\mathrm{tot}})\bigr)
 \;\Longrightarrow\;
 H^{p+q}(d_{\mathrm{bar}}).
\end{align}
Both spectral sequences converge because the filtrations are bounded
below (conformal weight $\geq 0$) and exhaustive (every element has
finite weight), and the complexes are complete with respect to the
filtration (the vertex algebra is $\mathbb{Z}_{\geq 0}$-graded by
conformal weight with finite-dimensional graded pieces).

The map $\Phi_0$ of Step~2 defines an isomorphism
${}^{\mathrm{BRST}}E_1^{p,q} \xrightarrow{\sim}
{}^{\mathrm{bar}}E_1^{p,q}$ at the $E_1$ page: at each conformal
weight~$p$, the associated graded complexes are identified by~$\Phi_0$,
which is a quasi-isomorphism.

\emph{Step~4: Lifting to a chain map.}
The assumed filtered lifting data promotes the $E_1$ comparison to a
map of filtered complexes. Once this filtered chain map is constructed,
the standard spectral sequence comparison theorem
(Eilenberg--Moore~\cite{EM66}; see also
Weibel~\cite[Theorem~5.2.12]{Weibel94}) implies that a filtered chain
map inducing an isomorphism on the $E_1$ page is a quasi-isomorphism.

Concretely, $\Phi$ is constructed inductively on filtration degree.
At level $N = 0$, $\Phi$ is the classical map~$\Phi_0$. At level~$N$,
given $\Phi|_{F^{\leq N-1}}$, one extends to $F^{\leq N}$ by solving
a lifting problem: the obstruction to extending lies in
$H^{n+1}(\mathrm{gr}^N\,\barB^{\mathrm{ch}}_0)$, which is the
$(n+1)$-st cohomology of the classical bar complex at weight~$N$.
The theorem assumes these obstruction classes vanish; in the standard
free-field and affine PBW cases this is the Koszul-resolution
acyclicity used in the preceding paragraph. The induction then proceeds.

The map $\Phi$ so constructed is a filtered chain map inducing an
isomorphism on $E_1$ pages. By the Eilenberg--Moore comparison
theorem, $\Phi$ is a quasi-isomorphism.
\end{proof}

\begin{remark}[BV/BRST reading of the master square]
\label{rem:bv-convergence}
\index{BV complex!convergence with bar complex}
Theorem~\ref{thm:brst-bar-genus0} compares the relative BRST complex
with the genus-$0$ chiral bar complex under the PBW lifting
hypotheses. It does not identify $\barB_X(\cA^!)$ with the ghost
complex. The ghost complex is the external BRST resolution determined
by the gauge symmetry and its chosen stress tensor; $\barB_X(\cA)$ is
the collision coalgebra; $\cA^!$ is obtained from the bar coalgebra by
Verdier duality on the Koszul locus; and
$Z^{\mathrm{der}}_{\mathrm{ch}}(\cA)
=C^\bullet_{\mathrm{ch}}(\cA,\cA)$ is the chiral derived centre, the
closed-sector object. The master square therefore commutes on the relative
BRST locus $K_{\mathrm{BRST,tot}}=0$ after the filtered comparison
map~$\Phi$ is constructed, while the curvature square
$\dfib^{\,2}=\kappa(\cA)\omega_1$ belongs to the separate
lineage-specific bar curvature package.
\end{remark}

\begin{remark}[Low-degree shape of \texorpdfstring{$\Phi$}{Phi}]
\label{rem:brst-bar-explicit}
For the bosonic string $\cA = \mathcal{H}^{26}$, the chain map
$\Phi$ has the following low-degree associated-graded form. The
displayed ghost representatives live first in the absolute ghost Fock
model; the relative complex uses their $b_0$-horizontal projections.
\begin{itemize}
\item \emph{Degree~$0$}: $\Phi(|0\rangle) = 1 \in \barB^0$.
 A BRST-closed state of ghost number~$0$ maps to the unit in
 the bar complex.
\item \emph{Degree~$1$}: $\Phi(c_1 \alpha_{-1}^i |0\rangle)
 = [\alpha^i] \in \barB^1$, where $\alpha^i$ are the Heisenberg
 generators. The physical vertex operators (momentum eigenstates)
 map to bar generators.
\item \emph{Degree~$2$}:
$\Phi(\pi_{\mathrm{rel}}(c_0 c_1 |0\rangle)) = \eta_{12}$,
 where $\eta_{12} = d\log(z_1 - z_2)$ is the FM propagator.
 In the relative complex the state $c_0c_1|0\rangle$ is replaced
 by its $b_0$-horizontal representative; the formula records the
 associated-graded ghost--propagator identification, not a literal
 relative-chain equality.
\end{itemize}
In general, ghost insertions $c_{n_1} \cdots c_{n_k}$ map to
products of logarithmic forms $\eta_{i_1 j_1} \wedge \cdots \wedge
\eta_{i_r j_r}$ on the configuration space, with the precise
correspondence determined by the residue map of the FM propagator.
\end{remark}

\begin{corollary}[Physical anomaly cancellation at genus \texorpdfstring{$0$}{0};
\ClaimStatusConditional]
\label{cor:anomaly-physical-genus0-vol1}
\index{anomaly cancellation!BRST proof}
The genus-$0$ BRST anomaly cancellation is
$K_{\mathrm{BRST,tot}}=0$, equivalently
$c_{\mathrm{matter}}+c_{\mathrm{ghost}}=0$. On the Virasoro
stress-tensor lane this is the bar curvature-vanishing condition of
Theorem~\textup{\ref{thm:anomaly-koszul}}. Specifically:
\begin{enumerate}[label=\textup{(\roman*)}]
\item $Q_{\mathrm{BRST}}^2 = 0$ on $\cA_{\mathrm{tot}}^{\mathrm{rel}}$
 if and only if $c(\cA) = 26$ \textup{(}Lemma~\textup{\ref{lem:brst-nilpotence})}.
\item $d_{\mathrm{bar}}^2 = 0$ on $\barB^{\mathrm{ch}}_0(\cA_{\mathrm{tot}})$
 on this stress-tensor lane if and only if
 $K_{\mathrm{BRST,tot}}=0$, i.e., $c(\cA)=26$
 \textup{(}Theorem~\textup{\ref{thm:anomaly-koszul})}.
\item The chain map $\Phi$ of Theorem~\textup{\ref{thm:brst-bar-genus0}}
 identifies both anomaly cancellation conditions under the PBW lifting
 hypotheses of that theorem: both sides are well-defined cochain
 complexes precisely when $c = 26$, and $\Phi$ is a quasi-isomorphism
 relating their cohomologies.
\end{enumerate}
\end{corollary}

\begin{proof}
Parts (i) and (ii) are the cited results. Part (iii):
Theorem~\ref{thm:brst-bar-genus0} constructs $\Phi$ under
the hypotheses $c = 26$ and filtered PBW lifting, giving a quasi-isomorphism that
identifies $H^*(Q_{\mathrm{BRST}}) \cong H^*(d_{\mathrm{bar}})$.
If $c \neq 26$, the BRST side has
\[
Q_{\mathrm{BRST}}^2
=\frac{c-26}{12}\sum_{n\in\mathbb Z}(n^3-n){:}c_{-n}c_n{:}\neq0
\]
as an operator on the full complex,
and the bar side has
$\dfib^{\,2}=K_{\mathrm{BRST,tot}}\cdot\omega_1
\cdot \mathrm{id} \neq 0$: neither is a cochain complex,
and both obstructions are controlled by the same numerical invariant
$c - 26$. The dictionary
\[
K_{\mathrm{BRST,tot}}
= K_{\mathrm{BRST}}(\cA_{\mathrm{matter}})
 + K_{\mathrm{BRST}}(\cA_{\mathrm{ghost}})
= \tfrac{c}{2} + \tfrac{-26}{2}
= \tfrac{c-26}{2}
\]
(Computation~\ref{comp:virasoro-curvature} for $\kappa(\mathrm{Vir}_c) = c/2$;
$K_{\mathrm{BRST}}(\text{$bc$-ghost}, \lambda{=}2)
=c_{\mathrm{gh}}/2=-13$)
converts the BRST obstruction $Q^2 \propto (c-26)$ into
the bar obstruction
$\dfib^{\,2}\propto K_{\mathrm{BRST,tot}}$ on this Virasoro lane.
\end{proof}

\begin{proposition}[Matched Koszul transport preserves BRST anomaly cancellation; \ClaimStatusConditional]
\label{prop:koszul-brst-anomaly-preservation-vol1}
\index{Koszul duality!BRST anomaly preservation}
\index{anomaly cancellation!Koszul duality}
Let $\cA_{\mathrm{tot}} = \cA_{\mathrm{matter}} \otimes \cA_{\mathrm{ghost}}$
be a BRST-quantized gauge system with total BRST conductor
$K_{\mathrm{BRST,tot}}
=K_{\mathrm{BRST}}(\cA_{\mathrm{matter}})
+K_{\mathrm{BRST}}(\cA_{\mathrm{ghost}})=0$.
Assume that the Koszul partner $\cA_{\mathrm{matter}}^!$ is equipped
with a compatible BRST stress tensor satisfying
\[
K_{\mathrm{BRST}}(\cA_{\mathrm{matter}}^!)
=-K_{\mathrm{BRST}}(\cA_{\mathrm{matter}}).
\]
Let $\widetilde{\cA}_{\mathrm{ghost}}$ be a transported ghost sector
whose BRST conductor satisfies
\[
K_{\mathrm{BRST}}(\widetilde{\cA}_{\mathrm{ghost}})
=-K_{\mathrm{BRST}}(\cA_{\mathrm{ghost}}).
\]
Then the dual total system
\[
 \cA_{\mathrm{tot}}^{\vee}
 \;=\; \cA_{\mathrm{matter}}^! \otimes \widetilde{\cA}_{\mathrm{ghost}}
\]
has vanishing total BRST conductor
\[
 K_{\mathrm{BRST}}(\cA_{\mathrm{tot}}^{\vee})
 \;=\; K_{\mathrm{BRST}}(\cA_{\mathrm{matter}}^!)
 + K_{\mathrm{BRST}}(\widetilde{\cA}_{\mathrm{ghost}})
 \;=\; -K_{\mathrm{BRST,tot}}
 \;=\; 0,
\]
and the BRST anomaly cancellation of
Corollary~\textup{\ref{cor:anomaly-physical-genus0-vol1}} holds for
$\cA_{\mathrm{tot}}^{\vee}$ at genus~$0$.
\end{proposition}

\begin{proof}
The BRST conductor is additive under tensor product because central
charges add. The two displayed hypotheses give
$K_{\mathrm{BRST}}(\cA_{\mathrm{matter}}^!)
=-K_{\mathrm{BRST}}(\cA_{\mathrm{matter}})$ and
$K_{\mathrm{BRST}}(\widetilde{\cA}_{\mathrm{ghost}})
=-K_{\mathrm{BRST}}(\cA_{\mathrm{ghost}})$, hence the total conductor
of $\cA_{\mathrm{tot}}^\vee$ is $-K_{\mathrm{BRST,tot}}=0$. The
genus-$0$ BRST cancellation follows from
Corollary~\ref{cor:anomaly-physical-genus0-vol1} applied to
$\cA_{\mathrm{tot}}^\vee$. No assertion is made that bar-cobar
reconstruction alone supplies the transported ghost sector or the
dual stress tensor.
\end{proof}

\begin{remark}[Scoping: why the ghost involution matters]
\label{rem:koszul-brst-scoping}
The hypothesis on $\widetilde{\cA}_{\mathrm{ghost}}$ is essential and
cannot be dropped. Consider the bosonic string:
$K_{\mathrm{BRST}}(\cA_{\mathrm{matter}})=13$ because
$c_{\mathrm{matter}}=26$, while the collision characteristic of the
same rank-$26$ Heisenberg algebra is
$\kappa(\mathcal H^{\oplus 26}_1)=26$. The $bc$ system at
$\lambda=2$ has
$c_{bc}=1-3(2{\cdot}2{-}1)^2=-26$ and
$K_{\mathrm{BRST}}(bc)=-13$, so
$K_{\mathrm{BRST,tot}}=13+(-13)=0$. A Koszul or Verdier operation on
the matter bar coalgebra does not automatically transform the
reparametrisation ghosts. If the ghost sector is kept fixed while the
matter stress tensor is replaced by a partner of conductor $-13$, the
new total conductor is $-13+(-13)=-26\neq0$, and BRST anomaly
cancellation fails. The proposition therefore requires an additional
ghost transport of conductor $+13$. This is an extra hypothesis on the
BRST package, not an internal self-duality of the $bc$ system and not
the Virasoro involution $c\mapsto26-c$.
\end{remark}

\begin{computation}[Superstring ghost/superghost Koszul dual;
\ClaimStatusConditional]
\label{comp:v1-superstring-ghost-koszul}
\index{superstring!ghost Koszul dual}
\index{$\beta\gamma$ system!superghost Koszul dual}
The bosonic string ghost system is the $bc$ pair at $\lambda = 2$
with $c_{bc} = 1 - 3(2{\cdot}2 - 1)^2 = -26$ and
$K_{\mathrm{BRST}}(bc)=c_{bc}/2=-13$. It is not Virasoro self-dual:
Virasoro self-duality is the same-family fixed point
$\mathrm{Vir}_c^!\simeq\mathrm{Vir}_{26-c}$, hence $c = 13$, whereas
the bosonic ghost sits at $c=-26$. The ghost transport required in
Remark~\ref{rem:koszul-brst-scoping} is a separate parity/Koszul
operation on the ghost package, not a Virasoro self-duality statement.

For the superstring, the ghost sector enlarges to
$(bc) \otimes (\beta\gamma)$, where the superghost
$\beta\gamma$ system sits at $\lambda = 3/2$. By the
formula $c_{\beta\gamma}(\lambda) = 2(6\lambda^2 - 6\lambda + 1)$,
\[
 c_{\beta\gamma}(\tfrac{3}{2})
 \;=\; 2\bigl(6 \cdot \tfrac{9}{4}
 - 6 \cdot \tfrac{3}{2} + 1\bigr)
 \;=\; 2\bigl(\tfrac{27}{2} - 9 + 1\bigr)
 \;=\; 2 \cdot \tfrac{11}{2}
 \;=\; 11,
\]
hence $K_{\mathrm{BRST}}(\beta\gamma)=c_{\beta\gamma}/2=11/2$, and
the transported dual satisfies
$K_{\mathrm{BRST}}((\beta\gamma)')=-11/2$.

The full superghost characteristic is
\[
 K_{\mathrm{BRST,ghost}}^{\mathrm{super}}
 \;=\; K_{\mathrm{BRST}}(bc) + K_{\mathrm{BRST}}(\beta\gamma)
 \;=\; -13 + \tfrac{11}{2}
 \;=\; -\tfrac{15}{2}.
\]
Superstring BRST anomaly cancellation requires
$K_{\mathrm{BRST,tot}} = K_{\mathrm{BRST,matter}}
+ K_{\mathrm{BRST,ghost}}^{\mathrm{super}} = 0$,
hence $K_{\mathrm{BRST,matter}} = 15/2$, equivalently
$c_{\mathrm{matter}} = 15$. For the $N = 1$ superconformal
matter of $d$ flat spacetime dimensions
($d$ free scalars contributing $d$ to $c$, plus $d$ free
Majorana fermions contributing $d/2$), this gives
$c_{\mathrm{matter}} = d + d/2 = 3d/2 = 15$, hence
$d = 10$: the critical dimension of the superstring. The
conductor sum for the transported superghost is
\[
 K_{\mathrm{BRST,ghost}}^{\mathrm{super}}
 + (K_{\mathrm{BRST,ghost}}^{\mathrm{super}})'
 \;=\; -\tfrac{15}{2} + \tfrac{15}{2}
 \;=\; 0,
\]
confirming that the superstring BRST anomaly cancellation
\textup{(}at $d = 10$\textup{)} is preserved by the matching
Koszul involution of Proposition~\ref{prop:koszul-brst-anomaly-preservation-vol1}.
\end{computation}

\begin{remark}[The ghost--superghost Koszul involution;
\ClaimStatusProvedElsewhere]
\label{rem:ghost-superghost-koszul}
\index{ghost!Koszul involution}
\index{superghost!Koszul involution}
The preceding computation identifies the ghost characteristic in
each of the two string backgrounds:
\begin{itemize}[nosep]
\item bosonic string: ghost $= bc$ at $\lambda = 2$, with
 $\kappa = -13$;
\item superstring: ghost $= bc$ at $\lambda = 2$ plus $\beta\gamma$
 at $\lambda = 3/2$, with total
 $K_{\mathrm{BRST}}=-15/2$.
\end{itemize}
In each case, anomaly preservation after matter dualization requires
a ghost transport that reverses the sign of the total ghost
characteristic while preserving the pairing used by the BRST
operator. For the $bc$ system this is not an internal same-family
self-duality: Friedan--Martinec--Shenker~\cite{FMS86} relate the
$bc$ and $\beta\gamma$ systems at bosonised level, and a parity-changed
ghost realization is the relevant opposite-characteristic model. For
the superstring combined ghost sector, the transport acts on the
combined $bc\otimes\beta\gamma$ package. The statement
$K_{\mathrm{BRST,ghost}}^{\mathrm{bos}} = -13$ and
$K_{\mathrm{BRST,ghost}}^{\mathrm{super}} = -15/2$ are the
numerical invariants before this transport, matching the standard
critical dimensions $d = 26$ and $d = 10$ of
Polchinski~\cite{Polchinski1998}. The involution is
\emph{not} the Virasoro self-duality at $c = 13$
\textup{(}which would require $\kappa + \kappa' = 13$\textup{)}:
the ghost transport is native to the $bc$\slash$\beta\gamma$ BRST
package and operates at a different locus.
\end{remark}

\begin{remark}[Higher genus]
\label{rem:brst-bar-higher-genus}
At genus $g \geq 1$, the BRST complex requires additional data:
the moduli-space density, gauge-fixing data, and
Costello renormalization framework~\cite{costello-renormalization} for
handling ultraviolet divergences. The bar complex side is well-defined
at all genera via configuration spaces on $\Sigma_g$
(Chapter~\ref{chap:higher-genus}). No all-genera chain map is
asserted here. Such a map would have to be genus-graded, compatible
with Costello's counterterm expansion, and square-zero only after the
curvature correction recorded in Theorem~\ref{thm:bv-bar-coderived-vol1};
it is not supplied by an unconstructed field-space measure.
The corresponding construction is recorded in
\S\ref{sec:modular-koszul-theory}.
\end{remark}

\section{Semi-infinite cohomology and the bar complex}
\label{sec:semi-infinite-bar}
\index{semi-infinite cohomology|textbf}

\subsection{Semi-infinite cohomology}

\begin{definition}[Semi-infinite cohomology {\cite{FGZ86}}]
\label{def:semi-infinite-cohomology}
\index{semi-infinite cohomology!definition}
Let $\widehat{L} = \bigoplus_{n \in \mathbb{Z}} L_n$ be a
$\mathbb{Z}$-graded Lie algebra with $\dim L_n < \infty$, and let
$M$ be a $\widehat{L}$-module. Write
$\widehat{L} = L_- \oplus L_0 \oplus L_+$ for the decomposition
into negative, zero, and positive modes. The \emph{semi-infinite
complex} is:
\[
C^{\infty/2+n}(\widehat{L}, M)
= M \otimes \textstyle\bigwedge\nolimits^{\infty/2+n}\! \widehat{L}^*
\]
where $\bigwedge^{\infty/2+n}$ denotes the semi-infinite wedge
power: the ``Dirac sea'' vacuum is the semi-infinite form
$\xi_0^* \wedge \xi_{-1}^* \wedge \xi_{-2}^* \wedge \cdots$
(wedging all negative-mode duals), and degree $n$ is measured
relative to this vacuum. The semi-infinite differential
$Q_{\mathrm{si}}$ combines the Chevalley--Eilenberg differential
with a ``charge anomaly'' correction term that ensures
$Q_{\mathrm{si}}^2 = 0$ without any restriction on the central charge.
The \emph{semi-infinite cohomology} is
$H^{\infty/2+\bullet}(\widehat{L}, M) =
H^*(C^{\infty/2+\bullet}(\widehat{L}, M), Q_{\mathrm{si}})$.
\end{definition}

\begin{remark}[Semi-infinite vs.\ string BRST]
\label{rem:semi-infinite-vs-brst}
The string BRST complex of \S\ref{sec:brst-bar-chain-map} couples
$\cA$ to $bc$-ghosts and requires $c = 26$ for $Q^2 = 0$
(Lemma~\ref{lem:brst-nilpotence}). Semi-infinite cohomology is
different: it is defined for any $\mathbb{Z}$-graded Lie algebra
$\widehat{L}$ and any module $M$, with $Q_{\mathrm{si}}^2 = 0$
unconditionally. The ghosts are ``built in'' to the semi-infinite
wedge product rather than introduced as an external $bc$-system.
For the string ($\widehat{L} = \mathrm{Vir}$, $M = \cA$),
semi-infinite cohomology of the Virasoro algebra \emph{does}
reduce to string BRST cohomology when $c = 26$~\cite{FGZ86}.
\end{remark}

\subsection{The bar--semi-infinite identification for Kac--Moody algebras}

\begin{theorem}[Bar complex = semi-infinite complex for KM;
\ClaimStatusConditional]
\label{thm:bar-semi-infinite-km}
\index{semi-infinite cohomology!bar complex identification}
\index{BRST!semi-infinite identification}
Let $\widehat{\fg}_k$ be the universal affine vertex algebra at level
$k \neq -h^\vee$. There is a quasi-isomorphism of filtered cochain
complexes:
\begin{equation}\label{eq:bar-semi-infinite}
\Psi \colon
 \bigl(\barB^{\mathrm{ch}}(\widehat{\fg}_k),\, d_{\mathrm{bar}}\bigr)
 \xrightarrow{\;\sim\;}
 \bigl(C^{\infty/2+\bullet}(\widehat{\fg}_k, V_k),\, Q_{\mathrm{si}}\bigr)
\end{equation}
identifying the chiral bar complex with the semi-infinite complex
\textup{(}with coefficients in the vacuum module $V_k$\textup{)}.
In particular:
\begin{equation}\label{eq:bar-equals-semi-infinite}
H^*(\barB^{\mathrm{ch}}(\widehat{\fg}_k))
\;\cong\;
H^{\infty/2+*}(\widehat{\fg}_k,\, V_k).
\end{equation}
\end{theorem}

\begin{proof}
The proof parallels that of
Theorem~\ref{thm:brst-bar-genus0}, replacing the string BRST
complex with the semi-infinite complex and removing the
$c = 26$ hypothesis.

\emph{Step~1: Conformal weight filtration.}
Both complexes carry bounded-below, exhaustive
$\mathbb{Z}_{\geq 0}$-filtrations by total conformal weight.
On the bar side, this is the PBW filtration of
Theorem~\ref{thm:brst-bar-genus0}, Step~1. On the semi-infinite
side, the conformal weight filtration arises from the
$L_0$-grading on $V_k$ and the mode grading on
$\bigwedge^{\infty/2+\bullet}\!\widehat{\fg}^*_k$. Both
differentials ($d_{\mathrm{bar}}$ and $Q_{\mathrm{si}}$) respect
these filtrations: $d_{\mathrm{bar}}$ preserves total conformal
weight (OPE residues preserve weight), and $Q_{\mathrm{si}}$
preserves it by construction~\cite{FGZ86}.

\emph{Step~2: Associated graded.}
On the associated graded $\mathrm{gr}\,F^\bullet$,
$\mathrm{gr}\,\barB^{\mathrm{ch}}(\widehat{\fg}_k)$ reduces
to the classical bar complex $\barB(\fg)$ of the
finite-dimensional Lie algebra $\fg$
(only the leading OPE pole survives, giving the Lie bracket),
while $\mathrm{gr}\,C^{\infty/2+\bullet}(\widehat{\fg}_k, V_k)$
reduces to the Chevalley--Eilenberg complex
$C^*_{\mathrm{CE}}(\fg, \mathbb{C})$
(the semi-infinite vacuum collapses to $\bigwedge^*\!\fg^*$
when mode corrections are dropped;
see~\cite{FGZ86}, Proposition~2.1).
Both associated graded complexes compute the Lie algebra cohomology
$H^*(\fg, \mathbb{C})$. The classical bar--CE comparison map
$\Phi_0$ of Arkhipov~\cite{Arkhipov97}
(cf.\ \eqref{eq:arkhipov-classical}) provides a natural
quasi-isomorphism $\barB(\fg) \xrightarrow{\sim}
C^*_{\mathrm{CE}}(\fg, \mathbb{C})$.

\emph{Step~3: Spectral sequence comparison.}
The conformal weight filtrations induce convergent spectral sequences
on both sides (bounded below, exhaustive, finite-dimensional
graded pieces). The map $\Phi_0$ induces an isomorphism at the
$E_1$ page. By the Eilenberg--Moore comparison theorem
(\cite{EM66}; cf.\ \cite[Theorem~5.2.12]{Weibel94}), $\Phi_0$ lifts
to a filtered chain map
$\Psi \colon \barB^{\mathrm{ch}}(\widehat{\fg}_k) \to
C^{\infty/2+\bullet}(\widehat{\fg}_k, V_k)$
that is a quasi-isomorphism.

\emph{Step~4: No anomaly obstruction.}
Unlike Theorem~\ref{thm:brst-bar-genus0}, no central charge
condition is needed. On the bar side, $d_{\mathrm{bar}}^2 = 0$
by the Arnold relations
(Theorem~\ref{thm:arnold-three}); the curvature
$m_0 = \kappa \cdot \mathbf{1}$ contributes to the $A_\infty$
structure but does not obstruct the bar differential. On the
semi-infinite side, $Q_{\mathrm{si}}^2 = 0$ by
construction~\cite{FGZ86}: the semi-infinite wedge product
absorbs the charge anomaly that would otherwise require
$c = 26$. Both differentials square to zero at every level
$k \neq -h^\vee$, and the spectral sequence comparison
applies without restriction.

The condition $k \neq -h^\vee$ ensures that the Sugawara
construction is well-defined, providing the $L_0$-grading
that drives the filtration. At critical level
$k = -h^\vee$, the Sugawara stress tensor is undefined, but the
semi-infinite complex is still defined using the loop rotation
grading; the identification at critical level follows from the
Feigin--Frenkel theorem~\cite{Feigin-Frenkel}
(cf.\ Theorem~\ref{thm:critical-level-structure}).
\end{proof}

\begin{remark}[String BRST as a special case of semi-infinite cohomology]
\label{rem:string-brst-special-case}
The string BRST comparison is the Virasoro analogue of
Theorem~\ref{thm:bar-semi-infinite-km}, not a literal specialization
of the affine Kac--Moody theorem. By~\cite{FGZ86}, the string BRST
complex of $\mathcal{H}^{26}\otimes bc$ with $c_{\mathrm{tot}}=0$ is
the semi-infinite complex of the Virasoro algebra acting on
$\mathcal{H}^{26}$. Under the PBW lifting hypotheses of
Theorem~\ref{thm:brst-bar-genus0}, the resulting Virasoro
semi-infinite comparison gives the map~$\Phi$. At
$c_{\mathrm{matter}}=26$, both the bosonic-string BRST anomaly
\[
\frac{c_{\mathrm{matter}}-26}{12}
\sum_{n\in\mathbb Z}(n^3-n){:}c_{-n}c_n{:}
\]
and $Q_{\mathrm{si}}^2$ vanish on the relative complex; away from
$c_{\mathrm{matter}}=26$, this quadratic ghost anomaly measures the
mismatch.
\end{remark}

\subsection{Anomaly duality from complementarity}

\begin{corollary}[Anomaly duality for Kac--Moody pairs;
\ClaimStatusConditional]
\label{cor:anomaly-duality-km-vol1}
\index{anomaly cancellation!Kac--Moody duality}
\index{Feigin--Frenkel duality!anomaly interpretation}
Let $\widehat{\fg}_k$ be an affine Kac--Moody algebra at generic
level $k$, with Feigin--Frenkel dual
$\widehat{\fg}_{k'} = \widehat{\fg}_{-k-2h^\vee}$. Then:
\begin{enumerate}[label=\textup{(\alph*)}]
\item \emph{Bar cohomology = semi-infinite cohomology.}
The bar cohomology computations of
Part~\ref{part:characteristic-datum} \textup{(}Riordan numbers for
$\widehat{\mathfrak{sl}}_2$, the Hilbert series of
Theorem~\textup{\ref{thm:bar-cohomology-level-independence})}
are semi-infinite cohomology computations:
\[
\dim H^n(\barB^{\mathrm{ch}}(\widehat{\fg}_k))
= \dim H^{\infty/2+n}(\widehat{\fg}_k, V_k).
\]

\item \emph{Curvature = conformal anomaly.}
The genus-$1$ curvature coefficient
$\kappa(\widehat{\fg}_k) = (k+h^\vee)\dim\fg/(2h^\vee)$
\textup{(}Theorem~\textup{\ref{thm:genus-universality}(i))} equals
the semi-infinite anomaly coefficient. For affine Kac--Moody pairs,
the Feigin--Frenkel anti-symmetry
$\kappa(\widehat{\fg}_k) + \kappa(\widehat{\fg}_{-k-2h^\vee}) = 0$
of Theorem~\textup{\ref{thm:genus-universality}(ii)} is the
anomaly cancellation for the Feigin--Frenkel pair: the
conformal anomalies of $\widehat{\fg}_k$ and
$\widehat{\fg}_{-k-2h^\vee}$ cancel.
\textup{(}This anti-symmetry holds for KM/free-field families where the
duality is level-negating; it fails for Virasoro, where
$\kappa(\mathrm{Vir}_c) + \kappa(\mathrm{Vir}_{26-c}) = 13 \neq 0$.\textup{)}

\item \emph{Complementarity = anomaly duality.}
The quantum complementarity theorem
\textup{(}Theorem~\textup{\ref{thm:quantum-complementarity-main})}
applied to the Koszul pair
$(\widehat{\fg}_k, \widehat{\fg}_{-k-2h^\vee})$ gives:
\[
Q_g(\widehat{\fg}_k) \oplus Q_g(\widehat{\fg}_{-k-2h^\vee})
= H^*(\overline{\mathcal{M}}_g,\, Z(\widehat{\fg}_k)).
\]
Via the bar--semi-infinite identification, the left side is a
direct sum of semi-infinite cohomology spaces on
$\overline{\mathcal{M}}_g$: the genus-$g$ quantum corrections
of the WZW model at level~$k$ and its Feigin--Frenkel dual sum to
the full moduli space cohomology. This is the mathematical
content of anomaly cancellation for Kac--Moody pairs at all genera.

\item \emph{Universal MC duality.}
The universal Maurer--Cartan classes satisfy
$\Theta_{\widehat{\fg}_k} + \Theta_{\widehat{\fg}_{-k-2h^\vee}}
= 0$ (Theorem~\ref{thm:explicit-theta}), since
$\kappa(\widehat{\fg}_k) + \kappa(\widehat{\fg}_{-k-2h^\vee}) = 0$
by Theorem~\ref{thm:genus-universality}(ii).
The chain-level mechanism underlying~(c) is:
the Verdier involution exchanges
$\Theta_{\widehat{\fg}_k}$ and $-\Theta_{\widehat{\fg}_{k'}}$,
producing the Lagrangian splitting of
Theorem~\ref{thm:quantum-complementarity-main}.
\end{enumerate}
\end{corollary}

\begin{proof}
Part~(a) is the cohomology-level consequence of
Theorem~\ref{thm:bar-semi-infinite-km}. Part~(b): the
curvature $m_0 = \kappa \cdot \mathbf{1}$ in the bar complex
corresponds, under the quasi-isomorphism~$\Psi$, to the
semi-infinite anomaly (the failure of the naive CE differential
to square to zero, corrected by the semi-infinite charge).
The genus universality theorem identifies this with the
conformal anomaly coefficient. The KM anti-symmetry
$\kappa(\widehat{\fg}_k) + \kappa(\widehat{\fg}_{-k-2h^\vee}) = 0$
follows from the Feigin--Frenkel level shift and the formula
$\kappa = (k+h^\vee)\dim\fg/(2h^\vee)$: replacing $k$ by
$-k - 2h^\vee$ negates~$\kappa$.
Part~(c): combine the bar--semi-infinite identification with
Theorem~\ref{thm:quantum-complementarity-main}. The genus-$g$
obstruction space $Q_g(\widehat{\fg}_k)$ is the associated
graded of $H^*(\barB^{(g)}(\widehat{\fg}_k))$, which by
Theorem~\ref{thm:bar-semi-infinite-km} is the genus-$g$
semi-infinite cohomology. The complementarity sum equals
$H^*(\overline{\mathcal{M}}_g, Z(\widehat{\fg}_k))$ by
Theorem~\ref{thm:quantum-complementarity-main}.
\end{proof}

\begin{remark}[The Master Table, revisited]
\label{rem:master-table-brst}
\index{master table of invariants!BRST interpretation}
Via Corollary~\ref{cor:anomaly-duality-km-vol1}, every entry in the Master Table (Table~\ref{tab:master-invariants}) is a semi-infinite anomaly coefficient: $\kappa$ is the semi-infinite anomaly, $c + c' = 2\dim\fg$ the total Feigin--Frenkel anomaly, and $F_g$ the genus-$g$ WZW free energy.
\end{remark}

\subsection{Extension to $\mathcal{W}$-algebras via Drinfeld--Sokolov reduction}

\begin{theorem}[DS transport of bar = semi-infinite for principal \texorpdfstring{$\mathcal{W}$}{W}-algebras;
\ClaimStatusConditional]
\label{thm:bar-semi-infinite-w}
\index{semi-infinite cohomology!W-algebra identification}
\index{Drinfeld--Sokolov reduction!semi-infinite cohomology}
Let $\mathcal{W}^k(\fg, f_{\mathrm{prin}})$ be the principal
W-algebra at level $k \neq -h^\vee$. Assume that the affine
bar--semi-infinite quasi-isomorphism of
Theorem~\ref{thm:bar-semi-infinite-km} is compatible with principal
quantum Drinfeld--Sokolov reduction on the relevant filtered
complexes: the bar construction intertwines the DS BRST
differential up to filtered quasi-isomorphism, the semi-infinite
complex intertwines with the DS differential, and passage to DS
cohomology preserves the affine comparison map. Then there is a
quasi-isomorphism:
\begin{equation}\label{eq:bar-semi-infinite-w}
\barB^{\mathrm{ch}}(\mathcal{W}^k(\fg))
\;\xrightarrow{\;\sim\;}\;
C^{\infty/2+\bullet}(\mathcal{W}^k(\fg),\, M_k)
\end{equation}
where $M_k$ is the vacuum module for $\mathcal{W}^k(\fg)$ and the
right side is the semi-infinite complex defined with respect to the
conformal weight grading. In particular:
\[
H^*(\barB^{\mathrm{ch}}(\mathcal{W}^k(\fg)))
\;\cong\;
H^{\infty/2+*}(\mathcal{W}^k(\fg),\, M_k).
\]
\end{theorem}

\begin{proof}
Under the stated transport hypotheses, begin with the affine
quasi-isomorphism of Theorem~\ref{thm:bar-semi-infinite-km}. The
first compatibility assumption identifies the DS reduction of the
affine bar complex with
$\barB^{\mathrm{ch}}(\mathcal{W}^k(\fg))$, while the second
identifies the DS reduction of the affine semi-infinite complex with
$C^{\infty/2+\bullet}(\mathcal{W}^k(\fg), M_k)$. The third
assumption ensures that the affine comparison map survives passage to
DS cohomology. The resulting reduced map is exactly the displayed
quasi-isomorphism.
\end{proof}

\begin{corollary}[Virasoro bar complex = semi-infinite complex;
\ClaimStatusConditional]
\label{cor:virasoro-semi-infinite}
\index{semi-infinite cohomology!Virasoro}
\index{Virasoro algebra!semi-infinite cohomology}
For the Virasoro algebra $\mathrm{Vir}_c$ at central charge
$c = 1 - 6(k+1)^2/(k+2)$ \textup{(}obtained as
$\mathcal{W}^k(\mathfrak{sl}_2)$\textup{)}:
\[
H^*(\barB^{\mathrm{ch}}(\mathrm{Vir}_c))
\;\cong\;
H^{\infty/2+*}(\mathrm{Vir}_c,\, M_c)
\]
where $M_c$ is the vacuum Virasoro module. The Motzkin-number
bar cohomology dimensions of
Theorem~\textup{\ref{thm:virasoro-chiral-koszul}} are therefore
semi-infinite cohomology dimensions.
\end{corollary}

\begin{proof}
Apply Theorem~\ref{thm:bar-semi-infinite-w} with
$\fg = \mathfrak{sl}_2$, $f = f_{\mathrm{prin}}$.
\end{proof}

\begin{corollary}[Curvature complementarity for principal \texorpdfstring{$\mathcal{W}$}{W}-algebra pairs;
\ClaimStatusConditional]
\label{cor:anomaly-duality-w}
\index{anomaly cancellation!W-algebra duality}
Let $\mathcal{W}^k(\fg)$ be the principal W-algebra, with
Feigin--Frenkel same-family partner
$\mathcal{W}^{k'}(\fg) = \mathcal{W}^{-k-2h^\vee}(\fg)$.
Then:
\begin{enumerate}[label=\textup{(\alph*)}]
\item The curvature coefficients satisfy the same
level-independent complementarity law as the central-charge
package. For principal finite-type $\mathcal{W}$-algebras,
\[
\kappa(\mathcal{W}^k(\fg)) + \kappa(\mathcal{W}^{k'}(\fg))
= \bigl(c(\mathcal{W}^k(\fg)) + c(\mathcal{W}^{k'}(\fg))\bigr)
\varrho(\fg),
\]
where $\varrho(\fg) = \sum_{i=1}^r \frac{1}{m_i+1}$ is the anomaly
ratio from Theorem~\ref{thm:wn-obstruction}. Thus the sum is a
level-independent constant: $13$ for Virasoro and $250/3$ for
$\mathcal{W}_3$.

\item Conditional on the same-family DS/bar transport clause of
Theorem~\ref{thm:w-algebra-koszul-main}, the quantum complementarity theorem
\textup{(}Theorem~\textup{\ref{thm:quantum-complementarity-main})}
applied to the W-algebra same-family Koszul pair
\textup{(}Theorem~\textup{\ref{thm:w-algebra-koszul-main})} gives:
\[
Q_g(\mathcal{W}^k) \oplus Q_g(\mathcal{W}^{-k-2h^\vee})
= H^*(\overline{\mathcal{M}}_g,\, Z(\mathcal{W}^k)).
\]
This is the genus-$g$ complementarity statement for the principal
$\mathcal{W}$-pair.

\item For $\fg = \mathfrak{sl}_2$
\textup{(}Virasoro\textup{)}: the Virasoro at
$c$ and its same-family shadow partner at $c' = 26 - c$ have
curvature coefficients adding to~$13$, and their
genus-$g$ quantum corrections sum to
 $H^*(\overline{\mathcal{M}}_g, Z(\mathrm{Vir}_c))$. Here
 $\mathrm{Vir}_{26-c}$ is the proved M/S-level shadow partner used in
 the semi-infinite calculation; the stronger H-level
 infinite-generator realization belongs to MC4
 (MC4 closed; Vol~I Theorem~\ref{thm:completed-bar-cobar-strong}).
\end{enumerate}
\end{corollary}

\begin{proof}
Combine Theorem~\ref{thm:wn-obstruction} with
Theorem~\ref{thm:central-charge-complementarity}: for the principal
finite-type $\mathcal{W}$-algebra,
$\kappa(\mathcal{W}^k(\fg)) = c(\mathcal{W}^k(\fg))\cdot\varrho(\fg)$.
Part~(b) then follows from the quantum complementarity theorem.
Part~(c): specialize to $\fg = \mathfrak{sl}_2$ with $h^\vee = 2$,
for which $c + c' = 26$ and $\varrho(\mathfrak{sl}_2)=1/2$.
\end{proof}

\begin{remark}[Conditional semi-infinite interpretation]
\label{rem:w-semi-infinite-conditional}
If Theorem~\ref{thm:bar-semi-infinite-w} holds, then
Corollary~\ref{cor:anomaly-duality-w} acquires the intended
semi-infinite anomaly interpretation for principal
$\mathcal{W}$-algebra pairs. This interpretation remains conditional
on explicit DS/bar compatibility.
\end{remark}

\section{Holomorphic-topological field theories}

The physical content (Gaiotto's framework, boundary chiral
algebras, HT boundary conditions, W-algebras from Higgs branches)
is in Volume~II\@. The BV algebra structure on the bar complex
from configuration-space geometry is in
\S\ref{sec:complete-bv-structure}.

% Labels preserved for cross-reference compatibility.
\label{rem:ht-from-n4-sym}%
\label{rem:boundary-chiral-algebra-bv}%
\label{rem:bar-cobar-ht-boundary}%
% conj:holographic-bar-cobar defined in genus_complete.tex
\label{rem:w-algebra-bar-cobar}%


\section{The BV algebra structure}
\label{sec:complete-bv-structure}

\subsection{BV algebra definition}

\begin{definition}[BV algebra]
\label{def:bv-algebra-complete}
\index{BV algebra|textbf}
A \emph{Batalin--Vilkovisky algebra} is a graded commutative algebra
$(A, \cdot)$ equipped with a degree-$+1$ operator
$\Delta\colon A \to A$ satisfying:
\begin{enumerate}[label=\textup{(\roman*)}]
\item \emph{Nilpotency}: $\Delta^2 = 0$.
\item \emph{Second-order}: $\Delta$ is a second-order operator, meaning the
 \emph{BV bracket}
 \[
 \{a, b\} := (-1)^{|a|}\bigl[\Delta(ab) - \Delta(a)\,b
 - (-1)^{|a|} a\,\Delta(b)\bigr]
 \]
 (the failure of $\Delta$ to be a graded derivation) satisfies the
 graded Leibniz rule in each slot:
 \begin{gather*}
 \{a, bc\} = \{a, b\}\,c + (-1)^{(|a|+1)|b|} b\,\{a, c\},\\
 \{a, b\} = -(-1)^{(|a|+1)(|b|+1)}\{b, a\}.
 \end{gather*}
\end{enumerate}
The bracket automatically satisfies the graded Jacobi identity.
\end{definition}

\subsection{BV structure from configuration spaces}

\begin{theorem}[Conditional configuration-space BV package; \ClaimStatusConditional]
\label{thm:config-space-bv}
Assume that the diagonal-residue operator on the logarithmic bar
complex extends to a degree-$+1$ second-order operator
\[
\Delta\colon \bar{B}^{\mathrm{ch}}(\mathcal{A})
\longrightarrow
\bar{B}^{\mathrm{ch}}(\mathcal{A})[1]
\]
with $\Delta^2 = 0$, that the induced failure-of-derivation bracket is
the configuration-space residue pairing, and that these structures are
functorial in morphisms of chiral algebras. Then
$\bar{B}^{\mathrm{ch}}(\mathcal{A})$ carries a BV algebra structure
whose product is wedge product of logarithmic forms, whose bracket is
the residue pairing, and whose Laplacian is the diagonal-residue
operator.
\end{theorem}

\begin{proof}
Under the stated assumptions, the product is graded commutative by the
wedge-product sign rule, the bracket is generated from $\Delta$ by the
BV second-order identity, and the remaining BV axioms are exactly the
assumed nilpotence, Jacobi, and functorial compatibility conditions.
\end{proof}

\subsection{Quantum master equation}

\begin{remark}[Quantum master equation; \ClaimStatusHeuristic]
\label{rem:quantum-master-complete}%
The \emph{quantum master equation} is taken in the sign convention
of Appendix~\ref{app:signs}:
\[
\hbar \Delta_{\mathrm{BV}} S
+ \frac{1}{2}\{S,S\}_{\mathrm{BV}} = 0,
\]
equivalently $\Delta_{\mathrm{BV}} e^{S/\hbar}=0$ after the standard
exponential normalization. The bar side supplies the Maurer--Cartan
equation
$D\Theta+\frac{1}{2}[\Theta,\Theta]=0$. Identifying these equations
requires a BV Laplacian on the bar model and remains conditional on
Theorems~\ref{thm:config-space-bv} and~\ref{thm:bv-functor}.
\end{remark}

\begin{remark}[BV quantization and bar-cobar duality; \ClaimStatusHeuristic]
\label{rem:bv-equals-bar-cobar}%
Three operations must be kept separate. The BV or BRST complex computes
gauge-fixed physical cohomology. The bar complex computes the
collision coalgebra. The bar-cobar counit
\[
\Omega^{\mathrm{ch}}\barB^{\mathrm{ch}}(\mathcal{A})
\xrightarrow{\;\sim\;} \mathcal{A}
\]
is reconstruction on the Koszul locus; it does not produce
$\mathcal{A}^!$. The Koszul dual algebra enters through
Verdier duality applied to the bar coalgebra.

For affine Kac--Moody algebras, the proved statement is the filtered
comparison
\[
H^\bullet\!\bigl(\barB^{\mathrm{ch}}(\widehat{\fg}_k)\bigr)
\cong
H^{\infty/2+\bullet}(\widehat{\fg}_k,V_k)
\]
of Theorem~\ref{thm:bar-semi-infinite-km}, together with the separate
bar-cobar reconstruction of
Theorem~\ref{thm:bar-cobar-isomorphism-main}. For general chiral
algebras the BV reading remains heuristic unless a PBW filtration and
the conditional BV Laplacian package are supplied.
\end{remark}

\subsection{Summary: BV as functor}

\begin{theorem}[Conditional BV functor package; \ClaimStatusConditional]
\label{thm:bv-functor}
\index{BV functor}
Assume the conditional BV package of
Theorem~\ref{thm:config-space-bv} and a Verdier-duality comparison
identifying the dual of the bar coalgebra with a factorization algebra
$\cA^!_\infty$ compatibly with the bracket, coproduct/product, and
differential. Then the assignment
\[\mathrm{BV}: \mathrm{ChirAlg}_X \longrightarrow \mathrm{BV\text{-}Alg}\]
is a lax monoidal functor, and the Verdier dual
$\mathbb{D}_{\operatorname{Ran}}(\bar{B}(\mathcal{A}))$ identifies
with a factorization algebra $\cA^!_\infty$ whose underlying complex
is equivalent to $\bar{B}(\mathcal{A}^!)$
\textup{(}Theorem~\textup{\ref{thm:verdier-bar-cobar}}\textup{)}.
On the Koszul locus, $\cA^!_\infty \simeq \cA^!$.
\end{theorem}

\begin{proof}
Functoriality of the underlying bar construction is standard. The
extra BV functoriality and lax monoidal statements are exactly the
assumed functorial and tensor-compatibility properties of the
conditional BV package, while the Verdier-duality comparison is part
of the additional hypothesis in the theorem statement.
\end{proof}

\begin{remark}[Shifted symplectic structure from BV]\label{rem:bv-shifted-symplectic}
\index{shifted symplectic!from BV}
Conditional on Theorems~\ref{thm:config-space-bv}
and~\ref{thm:bv-functor}, the BV bracket on
$\barB^{\mathrm{ch}}(\cA)$ would provide a bar-side
$(-1)$-shifted Poisson model for the deformation theory. The
unconditional shifted-symplectic statements used later in the
manuscript instead come from the Verdier pairing on
$C_g = R\Gamma(\overline{\mathcal{M}}_g,\mathcal{Z}(\cA))$
\textup{(}Proposition~\ref{prop:ptvv-lagrangian}\textup{)} and from
the ambient cyclic deformation formal moduli problem
\textup{(}Theorem~\ref{thm:ambient-complementarity-fmp}\textup{)}.
\end{remark}

\section{Resolution for the Heisenberg algebra at all genera}
\label{sec:bv-bar-heisenberg-all-genera}
\index{Heisenberg!BV/bar identification|textbf}
\index{BV algebra!bar complex identification!Heisenberg}

The free-field case proves the scalar free-energy identity
predicted by Conjecture~\ref{conj:master-bv-brst}
for the Heisenberg family. The proof uses four independent
arguments that all produce the same identity.

\begin{theorem}[BV $=$ bar for the Heisenberg at all genera;
\ClaimStatusConditional]
\label{thm:heisenberg-bv-bar-all-genera-vol1}
\index{Heisenberg!determinant-line scalar}
\index{Faber--Pandharipande number!BV/bar identification}
For the Heisenberg vertex algebra $\cH_\kappa$ at level
$\kappa \neq 0$ and for every genus $g \geq 1$:
\begin{equation}\label{eq:bv-bar-heisenberg}
F_g^{\mathrm{BV}}(\cH_\kappa)
\;=\; F_g^{\mathrm{bar}}(\cH_\kappa)
\;=\; \kappa \cdot \lambda_g^{\mathrm{FP}},
\end{equation}
where $\lambda_g^{\mathrm{FP}}
= \frac{2^{2g-1}-1}{2^{2g-1}} \cdot
\frac{|B_{2g}|}{(2g)!}$
is the Faber--Pandharipande intersection number
$\int_{\overline{\mathcal{M}}_{g,1}} \psi^{2g-2}\lambda_g$.
\end{theorem}

\begin{proof}
We give four independent proofs, each arriving at the same
identity by a different route. The first two are the primary
arguments; the third and fourth provide independent confirmation.

\medskip
\noindent\textbf{Preliminary: the bar side.}
Under~$H_D^K$, Theorem~D gives
$\operatorname{ch}_g(\mathfrak O_g^K(\cA))
=(-1)^g\kappa(\cA)\lambda_g$.  Under~$H_D^{\mathrm{graph}}$ its
scalar-diagonal trace is $\kappa(\cA)\lambda_g^{\mathrm{FP}}$.
The Heisenberg $\cH_\kappa$ qualifies:
\begin{itemize}[nosep]
\item $\kappa(\cH_\kappa) = \kappa$ (the level);
\item single generator of conformal weight~$1$
 (uniform-weight, no cross-channel correction at any genus;
 Remark~\ref{rem:propagator-weight-universality});
\item class~G (Gaussian): the OPE $a(z)a(w) \sim \kappa/(z{-}w)^2$
 has only a double pole, so shadow depth $r_{\max} = 2$ and all
 higher shadow components $\mathfrak{C}, \mathfrak{Q}, \ldots$
 vanish.
\end{itemize}
Therefore
$F_g^{\mathrm{bar}}(\cH_\kappa) = \kappa \cdot \lambda_g^{\mathrm{FP}}$
at every genus $g \geq 1$. It remains to show that the BV free
energy agrees.

\medskip
\noindent\textbf{Path \textup{(a)}: Quillen anomaly + GRR.}
\index{Quillen metric!BV/bar proof|(}
The zero-mode-reduced gauge-fixed free-boson determinant-line scalar
on a compact Riemann surface~$\Sigma_g$ of genus~$g$ is defined by
the zeta-regularized determinant
\begin{equation}\label{eq:bv-partition-heisenberg}
Z_g^{\mathrm{det}}(\cH_\kappa;\,\Sigma_g)
= \bigl(\det\nolimits'_\zeta \bar\partial_{\Sigma_g}\bigr)^{-\kappa},
\end{equation}
where $\det'_\zeta$ denotes the zeta-regularized determinant
with the zero modes removed. This is the analytic determinant-line
model for the free BV Gaussian; no ordinary measure on the field space
is being asserted.
Apply GRR to the universal curve
$\pi\colon \overline{\mathcal{C}}_g
\to \overline{\mathcal{M}}_g$
for the structure sheaf~$\mathcal{O}$.

\emph{Step~1} \textup{(}Quillen anomaly\textup{)}.
The Quillen metric $h_Q$ on the determinant line bundle
$\det(R\pi_*\mathcal{O})$ over
$\overline{\mathcal{M}}_g$ has curvature
\[
\operatorname{curv}(h_Q)
= -2\pi i\, c_1(\mathbb{E}),
\]
where $\mathbb{E} = R^0\pi_*\omega_{\mathcal{C}/\mathcal{M}_g}$
is the Hodge bundle of rank~$g$ \cite{Quillen85}.

\emph{Step~2} \textup{(}Scaling\textup{)}.
For $\kappa$ copies of the free boson, the total
determinant line bundle is
$\det(R\pi_*\mathcal{O})^{\otimes\kappa}$
with curvature $\kappa \cdot c_1(\mathbb{E})$.

\emph{Step~3} \textup{(}GRR + family index\textup{)}.
By the Grothendieck--Riemann--Roch theorem:
\[
\operatorname{ch}(\mathbb{E})
= \pi_*\bigl(\operatorname{Td}(\omega_\pi)\bigr)
= g + \sum_{j \geq 1}
\frac{B_{2j}}{(2j)!}\,\kappa_{2j-1},
\]
where $\kappa_j = \pi_*(c_1(\omega_\pi)^{j+1})$ are the
Miller--Morita--Mumford classes. Under $H_D^K$, the signed Hodge
class is $\operatorname{obs}^{\mathrm{Hdg}}_g
=(-1)^g\kappa\lambda_g$.  The graph-evaluated family index is
$F_g^{\mathrm{sc}}=\kappa\lambda_g^{\mathrm{FP}}$, and its generating
function is
$\sum_{g \geq 1}F_g^{\mathrm{sc}}\,\hbar^{2g}
= \kappa\cdot(\hat{A}(i\hbar) - 1)$
\textup{(}Theorem~\ref{thm:family-index}\textup{)},
where $\hat{A}(x) = (x/2)/\sinh(x/2)$.

\emph{Step~4} \textup{(}Faber--Pandharipande\textup{)}.
The intersection number
\[
\int_{\overline{\mathcal{M}}_{g,1}}
\psi^{2g-2}\,\lambda_g
= \lambda_g^{\mathrm{FP}}
= \frac{2^{2g-1}-1}{2^{2g-1}}
\cdot \frac{|B_{2g}|}{(2g)!}
\]
is the Faber--Pandharipande formula~\cite{FP00}.
Combining Steps~1--4:
$F_g^{\mathrm{BV}}(\cH_\kappa)
= \kappa \cdot \lambda_g^{\mathrm{FP}}
= F_g^{\mathrm{bar}}(\cH_\kappa)$.
\index{Quillen metric!BV/bar proof|)}

\medskip
\noindent\textbf{Path \textup{(b)}: Selberg zeta function.}
\index{Selberg zeta function!BV/bar proof|(}
The D'Hoker--Phong determinant formula~\cite{DP86} expresses
the Laplacian determinant on a compact hyperbolic
surface~$\Sigma_g$ as
\begin{equation}\label{eq:dhoker-phong-det}
\det\nolimits'_\zeta(\Delta_{\Sigma_g})
= Z_{\mathrm{Sel}}(1;\Sigma_g)\cdot e^{c_g},
\end{equation}
where $Z_{\mathrm{Sel}}(s;\Sigma_g)
= \prod_{\gamma\, \text{prim.}} \prod_{n=0}^\infty
(1 - e^{-(s+n)\ell(\gamma)})$
is the Selberg zeta function, the product running
over primitive closed geodesics, and $c_g$ is a genus-dependent
spectral constant
\textup{(}Theorem~\ref{V2-thm:thqg-I-determinant-formula}\textup{)}.
The Heisenberg determinant-line scalar factorizes as
\begin{equation}\label{eq:heisenberg-selberg}
Z_g^{\mathrm{det}}(\cH_\kappa;\Sigma_g)
= \bigl(\det\operatorname{Im}\Omega\bigr)^{-\kappa/2}
\cdot Z_{\mathrm{Sel}}(1;\Sigma_g)^{-\kappa/2}
\cdot e^{-\kappa c_g/2},
\end{equation}
where $\Omega$ is the period matrix of~$\Sigma_g$
\textup{(}Corollary~\ref{V2-cor:thqg-I-heisenberg-selberg}\textup{)}.
The Mumford isomorphism
$\lambda = c_1(\det\mathbb{E})$ combined with the
Belavin--Knizhnik formula~\cite{BK86} gives
$\det\operatorname{Im}\Omega \sim |\vartheta|^2$ on
$\mathcal{M}_g$, so that the variation of~\eqref{eq:heisenberg-selberg}
over the moduli space produces a curvature form proportional
to $\kappa \cdot c_1(\mathbb{E})$. Extraction of
the top Chern class yields
$F_g^{\mathrm{BV}}(\cH_\kappa)
= \kappa \cdot \lambda_g^{\mathrm{FP}}$,
confirming~\eqref{eq:bv-bar-heisenberg} by a second route.
\index{Selberg zeta function!BV/bar proof|)}

\medskip
\noindent\textbf{Path \textup{(c)}: direct family index.}
\index{family index theorem!BV/bar proof|(}
Apply the Atiyah--Singer family index theorem
directly to $\bar\partial^{\oplus\kappa}$
on the universal curve $\pi \colon \overline{\mathcal{C}}_g
\to \overline{\mathcal{M}}_g$.
The family index of $\bar\partial$ in
$K^0(\overline{\mathcal{M}}_g)$ is $R\pi_*\mathcal{O}$,
and the Chern character of this virtual bundle is computed by GRR as
\[
\operatorname{ch}(R\pi_*\mathcal{O})
= 1 - g + \sum_{j \geq 1}
\frac{B_{2j}}{(2j)!}\,\kappa_{2j-1}.
\]
The $\hat{A}$-genus of the relative tangent bundle
$T_\pi = \omega_\pi^{-1}$ satisfies
$\hat{A}(T_\pi) = \operatorname{Td}(\omega_\pi)$
(since $\hat{A}$ and $\operatorname{Td}$ are related by
$\hat{A}(E) = e^{-c_1(E)/2}\operatorname{Td}(E)$, and for
a line bundle $L$,
$\hat{A}(L^{-1}) = \operatorname{Td}(L)$).
Therefore the generating function for the genus-$g$ contributions is
\[
\sum_{g \geq 1} F_g(\cH_\kappa)\,\hbar^{2g}
= \kappa \cdot \bigl(\hat{A}(i\hbar) - 1\bigr)
= \kappa \cdot \Bigl(\frac{\hbar/2}{\sin(\hbar/2)} - 1\Bigr)
\]
\textup{(}\S\ref{subsec:why-ahat}\textup{)}.
The Taylor coefficient at $\hbar^{2g}$ is
$\kappa \cdot \frac{(2^{2g-1}-1)\,|B_{2g}|}{2^{2g-1}\,(2g)!}
= \kappa \cdot \lambda_g^{\mathrm{FP}}$,
confirming~\eqref{eq:bv-bar-heisenberg} by a third route.
Note: Path~\textup{(a)} derives the same generating function
from the Quillen anomaly and then extracts coefficients; this path
arrives at the generating function from the family index theorem
without passing through the determinant line bundle.
\index{family index theorem!BV/bar proof|)}

\medskip
\noindent\textbf{Path \textup{(d)}: numerical verification.}
\index{multi-path verification!Heisenberg BV/bar}
At each genus $g = 1, \ldots, 15$, three independent
implementations produce the same rational number $\lambda_g^{\mathrm{FP}}$:
\begin{enumerate}[label=\textup{(\roman*)},nosep]
\item the Bernoulli formula
 $\lambda_g^{\mathrm{FP}}
 = \frac{(2^{2g-1}-1)\,|B_{2g}|}{2^{2g-1}\,(2g)!}$;
\item the Taylor extraction of the $\hat{A}$-series
 $(\hbar/2)/\sin(\hbar/2) - 1
 = \sum_{g \geq 1} \lambda_g^{\mathrm{FP}}\,\hbar^{2g}$;
\item a direct computation of the Faber--Pandharipande
 intersection integral on $\overline{\mathcal{M}}_{g,1}$
 via the string/dilaton equation and Witten's conjecture.
\end{enumerate}
All three agree at every genus tested. Sample values:
$\lambda_1^{\mathrm{FP}} = 1/24$,
$\lambda_2^{\mathrm{FP}} = 7/5760$,
$\lambda_3^{\mathrm{FP}} = 31/967680$.
\textup{(}Verified in the compute layer:
\texttt{heisenberg\_bv\_bar\_proof.py}, $60$ tests.\textup{)}
\end{proof}

\begin{remark}[Why the Heisenberg is special]
\label{rem:heisenberg-bv-bar-scope}
\index{BV algebra!bar complex identification!scope}
Theorem~\ref{thm:heisenberg-bv-bar-all-genera-vol1} proves the
\emph{scalar} \textup{(}partition function\slash free energy\textup{)}
consequence predicted by
Conjecture~\ref{conj:master-bv-brst} for the Heisenberg
family at all genera.
The Heisenberg succeeds because it is class~G:
a single Gaussian channel, no interaction vertices, and
all moduli dependence controlled by the Quillen anomaly.
For non-free theories, the scalar identification requires
the shadow obstruction tower machinery
\textup{(}Theorem~D for uniform-weight algebras;
the multi-weight genus expansion for $W_N$\textup{)}.

What remains open, even for
free fields, is:
\begin{enumerate}[label=\textup{(\roman*)},nosep]
\item the chain-level quasi-isomorphism between the BV complex and
 the bar complex at genus $g \geq 1$;
\item the identification of the BV Laplacian with the sewing
 operator on the bar complex;
\item the full quantum master equation as a chain-level identity
 in the modular deformation complex
 \textup{(}not just the scalar trace\textup{)}.
\end{enumerate}
The scalar identification is an index-theoretic statement
(the Euler characteristic of the BV complex equals the Euler
characteristic of the bar complex); the chain-level identification
is a quasi-isomorphism, which is strictly stronger.
\end{remark}

\begin{proposition}[Three chain-level obstructions and harmonic factorization;
\ClaimStatusProvedHere]
\label{prop:chain-level-three-obstructions}
\index{BV algebra!chain-level obstructions|textbf}
\index{chain-level BV/bar identification}
The proposed chain-level identification between the BV complex and
the bar complex at genus $g \geq 1$
faces three obstructions, controlled by the shadow depth together
with the separate harmonic higher-brace coherence data of the
algebra.
\begin{enumerate}[label=\textup{(\arabic*)},nosep]
\item \textbf{Propagator regularity.}
 The BV propagator $P(z,w) = \bar\partial^{-1}\delta(z,w)$
 is a distribution; the bar propagator
 $d\log E(z,w)$ is meromorphic.
 By Hodge decomposition, $P = d\log E + P_{\mathrm{ex}}
 + P_{\mathrm{harm}}$
 where $P_{\mathrm{ex}}$ is $\bar\partial$-exact
 and drops in Dolbeault cohomology.
 The harmonic part $P_{\mathrm{harm}}$ integrates to a
 moduli-dependent prefactor controlled by the Quillen anomaly.
 For the Heisenberg: resolved at all genera.
 For interacting theories: the smooth correction couples to
 the OPE, producing chain-level discrepancies.
\item \textbf{Moduli dependence.}
 Both the BV Green function and the bar prime form
 depend on the complex structure of~$\Sigma_g$.
 The Quillen anomaly formula relates
 $\operatorname{curv}(h_Q)$ to $c_1(\mathbb{E})$
 universally: this resolves moduli dependence for any algebra
 whose determinant-line scalar is determined by a zeta-regularized
 determinant.
\item \textbf{Higher-degree coupling through the harmonic propagator.}
 For non-Gaussian theories, the BV Laplacian contracts
 field-antifield pairs through interaction vertices.
 On the bar side, the sewing operator contracts bar
 generators through the Bergman kernel.
 When the OPE has cubic and higher poles, the BV contraction
 passes through interaction vertices, creating contributions
 that do not factor through the sewing kernel alone.
 \begin{itemize}[nosep]
 \item \emph{Class~G} \textup{(}Heisenberg\textup{)}: absent.
 No interaction vertices.
 \item \emph{Class~L} \textup{(}affine Kac--Moody\textup{)}:
 the cubic harmonic correction vanishes by the Jacobi identity
 $f^{ab[c}f^{d]be} = 0$.
 Shadow depth $r_{\max} = 3$ kills only the scalar
 $\Sigma$-coinvariant shadow invariants \(S_r\) for \(r\geq4\);
 it does not by itself kill all quartic or higher brace
 components before the shadow projection. Obstruction~\textup{(3)}
 is killed at every genus only under the corresponding harmonic
 higher-brace coherence condition.
 \item \emph{Class~C}
 \textup{(}$\beta\gamma$\textup{)}:
 conditional on harmonic decoupling.
 Although the quartic interaction vertex is nonzero
 \textup{(}shadow depth $r_{\max} = 4$\textup{)},
 the proved conclusion here is only the reduction: the first
 possible chain-level discrepancy is the quartic harmonic insertion,
 and its vanishing is exactly the harmonic-decoupling condition.
 The proposition does not prove that decoupling statement.
 \item \emph{Class~M}
 \textup{(}Virasoro, $\cW_N$\textup{)}:
 the naive ordinary chain-level comparison fails, but the
 quartic and higher interaction vertices produce a harmonic
 discrepancy that factors through the curvature:
 \[
\delta_r^{\mathrm{harm}}
\;\in\;
\operatorname{End}(H_g)\cdot m_0^{\lfloor r/2 \rfloor - 1}
\qquad\text{for every } r \geq 4.
 \]
 Hence the class~$\mathsf{M}$ obstruction is curvature-divisible
 even when it does not vanish chainwise.
 \end{itemize}
\end{enumerate}
\end{proposition}

\begin{proof}
Obstruction~(1): the Hodge decomposition
$P = P_{\mathrm{mer}} + P_{\mathrm{ex}} + P_{\mathrm{harm}}$
on $\Sigma_g$ is standard. The $\bar\partial$-exact part
$P_{\mathrm{ex}}$ vanishes in the Dolbeault cohomology
computing both the BV and bar homology.

Obstruction~(2): the Quillen anomaly formula
$\operatorname{curv}(h_Q) = -2\pi i\,c_1(\mathbb{E})$
\cite{Quillen85} is independent of the chiral algebra;
it depends only on the universal curve.

Obstruction~(3): let $H_g$ denote the harmonic summand of the
fiber model at genus~$g$. Hodge decomposition gives
\[
P_{\mathrm{BV}}
\;=\;
d\log E + [d_{\mathrm{fib}}, h] + P_{\mathrm{harm}},
\]
where $h$ is the Hodge homotopy and $P_{\mathrm{harm}}$ is the
harmonic projector. The commutator term contributes only a chain
homotopy, so the discrepancy between the BV differential and the
bar differential is concentrated in the harmonic part. Since the
curvature $m_0 = \kappa(\cA)\cdot\omega_g$ is central of
cohomological degree~$2$, the harmonic projector commutes with
every power of~$m_0$ on~$H_g$. A degree-$r$ harmonic discrepancy
must preserve the harmonic summand, lower the number of external
legs by pairs, and contribute even degree
$2\lfloor r/2 \rfloor - 2$. On~$H_g$ the only degree-compatible
central insertion with that degree is
$m_0^{\lfloor r/2 \rfloor - 1}$. Therefore
\[
\delta_r^{\mathrm{harm}}
\;\in\;
\operatorname{End}(H_g)\cdot m_0^{\lfloor r/2 \rfloor - 1}
\qquad\text{for every } r \geq 4.
\]
For class~G, there are no interaction vertices, so the harmonic
discrepancy vanishes. For class~L
\textup{(}affine Kac--Moody at non-critical level
$k \neq -h^\vee$\textup{)}, the interaction vertex is the
structure-constant tensor $f^{abc}$. The cubic harmonic
correction has the form
$\sum_c f^{abc}\cdot I_{\mathrm{harm}}(z_1,z_2)\cdot f^{cde}$,
where $I_{\mathrm{harm}}$ is the integral of the harmonic
propagator against the cubic vertex measure. Antisymmetry of
$f^{abc}$ and the Jacobi identity
$f^{abc}f^{cde} + \text{cyclic} = 0$ force this coefficient to
vanish. Shadow depth $r_{\max} = 3$ excludes scalar shadow
invariants \(S_r\) for \(r\geq4\), but it does not exclude
quartic or higher brace terms before the shadow projection; their
ordinary chain-level vanishing is the separate harmonic
higher-brace coherence input. For class~C
\textup{(}$\beta\gamma$\textup{)}, the quartic correction is the
first possible harmonic term. The ordinary chain-level coefficient
vanishes precisely under the additional harmonic-decoupling
condition, namely factorization of the harmonic insertion through the
composite free-field trace sector; this is a separate conditional
input and is not used in the proved classification above. For
class~M no cancellation identity forces
$c_r(\cA)$ to vanish, so the ordinary chain-level discrepancy
survives even though it is curvature-divisible.
\end{proof}

\begin{remark}[BV/bar identification for class~C:
conditional three-mechanism test; \ClaimStatusConditional]%
\label{rem:bv-bar-class-c-proof}%
\index{BV algebra!bar complex identification!class C|textbf}%
\index{$\beta\gamma$ system!BV/bar identification}%
\index{harmonic decoupling!role separation}%
The chain-level BV/bar identification
(Conjecture~\ref{conj:master-bv-brst})
for class~C algebras ($\beta\gamma$ systems) at genus~$1$ reduces to
the following harmonic-decoupling test. The test is not a proof of the
conjecture until the harmonic projector is shown to factor through the
free-field trace sector.
\begin{enumerate}[label=\textup{(\roman*)},nosep]
\item \emph{OPE structure.}
 The $\beta\gamma$ OPE $\beta(z)\gamma(w)
 \sim 1/(z{-}w)$ has a simple pole only.
 The quartic interaction vertex arises not from the
 binary OPE (which is abelian: $\beta\beta \sim 0$,
 $\gamma\gamma \sim 0$) but from the composite field
 $:\!\beta\partial\gamma\!:$ coupling to the stress tensor
 $T = :\!\beta\partial\gamma\!:$, which is the source of the
 shadow depth $r_{\max} = 4$.
 The harmonic-propagator correction at the quartic vertex
 therefore involves the composite source
 $:\!\beta\partial\gamma\!:$, not a fundamental cubic pole.
\item \emph{Hodge type.}
 The $\beta\gamma$ system has conformal weights
 $(h_\beta, h_\gamma) = (\lambda, 1{-}\lambda)$.
 The BV propagator's harmonic part $P_{\mathrm{harm}}$
 on $\Sigma_g$ is a $(0,1)$-form in
 $\bar\partial$-cohomology;
 the quartic vertex measure is built from holomorphic fields and their
 derivatives on the configuration space. This bidegree separation is an
 obstruction to a local scalar term, but it is not by itself a
 vanishing theorem on the compactified configuration space. The missing
 assertion is that the harmonic insertion has no boundary residue after
 the Fulton--MacPherson expansion.
\item \emph{Role separation.}
 The key structural feature distinguishing class~C from
 class~M: in the $\beta\gamma$ system, the field producing
 the quartic interaction ($T = :\!\beta\partial\gamma\!:$)
 is a \emph{composite} of the fundamental fields
 $\beta, \gamma$ that carry the propagator.
 The propagator sews $\beta$--$\gamma$ pairs;
 the quartic vertex is built from composites of those
 same pairs. The required decoupling statement is that, after expanding
 the harmonic correction in the fundamental fields, it factors through
 the genus-$1$ trace of the underlying free-field sector, whose scalar
 anomaly is already resolved.
 For class~M (Virasoro, $\cW_N$), the stress tensor $T$ is
 itself a fundamental generator with its own propagator channel;
 the quartic vertex involves $T$ coupling to $T$, and the
 harmonic correction does not factor through a free subsystem.
\end{enumerate}
Thus the genus-$1$ class~C statement is conditional on harmonic
decoupling: after the boundary-residue assertion is supplied, the
quartic harmonic-propagator correction vanishes despite the nonzero
quartic shadow $Q^{\mathrm{contact}}_{\beta\gamma}$. This is the local
model for the all-genera harmonic-decoupling hypothesis used in
Theorem~\ref{thm:bv-bar-coderived-vol1}: once every harmonic insertion
factors through the same composite free-field sector, the quartic and
higher chain-level corrections vanish. For class~$\mathsf{M}$, the
quartic harmonic discrepancy
$\delta_4^{\mathrm{harm}} \propto Q^{\mathrm{contact}} \cdot
\kappa / \mathrm{Im}(\tau)$ is not cancelled by the Fay trisecant
identity. Proposition~\ref{prop:chain-level-three-obstructions}
shows that the full harmonic discrepancy still factors through the
curvature. Conditional on the harmonic-factorization package of
Theorem~\ref{thm:bv-bar-coderived-vol1}, this identifies the
class~$\mathsf{M}$ BV and bar models in the coderived category;
the naive strict chain-level comparison still fails on the
uncompleted models.
\end{remark}

\begin{remark}[BV sewing at the chain level: class-by-class status;
 \ClaimStatusConditional]%
\label{rem:bv-sewing-chain-level-classes-vol1}%
\index{BV algebra!sewing operator identification|textbf}%
\index{sewing operator!BV Laplacian identification}%
\index{Bergman kernel!sewing contraction}%
The identification $\Delta_{\mathrm{BV}} = d_{\mathrm{sew}}$ at
the chain level asserts that the BV Laplacian and the bar sewing
operator agree as $(g,n) \to (g{+}1,n{-}2)$ operations on the
modular convolution algebra~$\gAmod$: both contract a pair of
inputs through the Bergman kernel $d\log E(z,w)$ along the
non-separating boundary divisor
$\delta^{\mathrm{ns}}\colon
\overline{\mathcal{M}}_{g,n+2} \to
\overline{\mathcal{M}}_{g+1,n}$.
Four complementary descriptions of this comparison
\textup{(}operator definition, spectral sequence, Heisenberg
extraction, modular operad\textup{)} lead to the following
class-by-class obstruction profile.
\begin{center}
\small
\renewcommand{\arraystretch}{1.15}
\begin{tabular}{lcl}
\toprule
Class & Chain-level status & Comment \\
\midrule
$\mathsf{G}$ \textup{(}Heisenberg\textup{)}
 & chain-level curved equivalence & no interaction vertices \\
$\mathsf{L}$ \textup{(}affine KM\textup{)}
 & chain-level curved equivalence & Jacobi kills the cubic harmonic term \\
$\mathsf{C}$ \textup{(}$\beta\gamma$\textup{)}
 & chain-level under harmonic decoupling & composite free-field factorization \\
$\mathsf{M}$ \textup{(}Virasoro, $\cW_N$\textup{)}
 & coderived equivalence; naive chain level false & harmonic correction survives but is curvature-divisible \\
\bottomrule
\end{tabular}
\end{center}
For class~$\mathsf{M}$, the quartic discrepancy
$\delta_4^{\mathrm{harm}} \propto Q^{\mathrm{contact}} \cdot
\kappa / \operatorname{Im}(\tau)$ is not a coboundary
in the ordinary chain complex. Proposition~\ref{prop:chain-level-three-obstructions}
shows that the full harmonic discrepancy is a sum of curvature
powers. Under the harmonic-factorization package in
Theorem~\ref{thm:bv-bar-coderived-vol1}, the resulting comparison
cone is coacyclic in the sense of
Definition~\ref{def:coacyclic-fact}. The remaining open questions are
the filtered-completed refinement for class~$\mathsf{M}$
\textup{(}Remark~\ref{rem:bv-bar-class-m-frontier}\textup{)} and the
all-genera harmonic-decoupling verification for class~$\mathsf{C}$.
\end{remark}

\begin{proposition}[Harmonic factorization of higher bar differentials;
\ClaimStatusConditional]%
\label{prop:harmonic-factorization}%
\index{harmonic factorization|textbf}%
\index{bar complex!harmonic projection}%
\index{curvature!harmonic factorization through}%
Let~$\cA$ be a chirally Koszul algebra with curvature
$m_0 = \kappa(\cA)\cdot\omega_g$ at genus $g \geq 1$
\textup{(}cohomological degree~$+2$, central\textup{)}.
Let $d_{\mathrm{fib}}$ denote the fiber bar differential,
$H_g \subset \ker(d_{\mathrm{fib}})$ the harmonic subspace
of the Hodge decomposition, and
$P_{\mathrm{harm}} \colon V_g \twoheadrightarrow H_g$
the harmonic projector. For each $r \geq 2$, let
$\delta_r$ denote the degree-$r$ component of the discrepancy
between the BV differential and the bar differential, arising
from insertion of the harmonic propagator at
$r$-valent interaction vertices. Then:
\begin{enumerate}[label=\textup{(\roman*)}]
\item The harmonic projection
$\delta_r^{\mathrm{harm}}
:= P_{\mathrm{harm}} \circ \delta_r \big|_{H_g}$
factors through powers of the curvature:
\[
\delta_r^{\mathrm{harm}}
\;=\;
\epsilon_r(\cA)\circ m_0^{\lfloor r/2 \rfloor - 1}
\qquad\text{for every } r \geq 2,
\]
for some central degree-zero endomorphism
$\epsilon_r(\cA)\in \operatorname{End}(H_g)$. The scalar
coefficient $c_r(\cA)$ denotes only the scalar/vacuum-channel
projection of~$\epsilon_r(\cA)$.
\item On the scalar/vacuum channel at low degrees:
$c_2(\cA) = 0$ \textup{(}the binary discrepancy is absorbed
by the Hodge homotopy\textup{)};
$c_3(\cA) = 0$ for classes~$\mathsf{G}$ and~$\mathsf{L}$
\textup{(}vanishes by the Jacobi identity for affine
Kac--Moody; absent for Heisenberg\textup{)};
$c_4(\cA) = S_4(\cA)$, the quartic shadow invariant.
\item For class~$\mathsf{G}$, $c_r(\cA) = 0$ for all $r$.
For class~$\mathsf{L}$, $c_r(\cA) = 0$ for all $r$
\textup{(}shadow depth $r_{\max} = 3$, so $c_r = 0$ for $r \geq 4$;
and $c_3 = 0$ by Jacobi\textup{)}.
For class~$\mathsf{M}$ \textup{(}Virasoro, $\cW_N$\textup{)},
the scalar/vacuum-channel coefficients $c_r(\cA)$ are generically
nonzero for $r \geq 4$, so the divisibility exponent cannot be
lowered on that channel.
\end{enumerate}
\end{proposition}

\begin{proof}
The proof proceeds through Hodge decomposition, degree
counting on the harmonic subspace, and centrality of the
curvature.

\emph{Step~1: Hodge decomposition isolates the harmonic
discrepancy.}
At genus $g \geq 1$, the Hodge decomposition of the BV
propagator is
\[
P_{\mathrm{BV}}
\;=\;
d\log E + [d_{\mathrm{fib}}, h] + P_{\mathrm{harm}},
\]
where $d\log E$ is the bar propagator \textup{(}the prime
form\textup{)}, $h$ is the Hodge homotopy, and
$P_{\mathrm{harm}}$ projects onto the harmonic
subspace~$H_g$. The commutator $[d_{\mathrm{fib}}, h]$
is $d_{\mathrm{fib}}$-exact and therefore contributes
only a chain homotopy: replacing $P_{\mathrm{BV}}$ by
$d\log E + P_{\mathrm{harm}}$ changes the BV differential
by a chain-homotopic deformation. The discrepancy between
the BV differential and the bar differential is therefore
concentrated in the harmonic part $P_{\mathrm{harm}}$.

\emph{Step~2: Centrality of $m_0$ on $H_g$.}
The curvature $m_0 = \kappa(\cA) \cdot \omega_g$
is central: it commutes with all operations in the fiber
model. In particular, $m_0$ commutes with the harmonic
projector $P_{\mathrm{harm}}$ and preserves the
decomposition
$V_g = \operatorname{im}(d_{\mathrm{fib}})
\oplus H_g \oplus \operatorname{im}(d_{\mathrm{fib}}^*)$.
The restriction of any power $m_0^k$ to~$H_g$ is again
central and has cohomological degree~$2k$.

\emph{Step~3: Degree counting.}
The degree-$r$ harmonic discrepancy
$\delta_r^{\mathrm{harm}}
= P_{\mathrm{harm}} \circ \delta_r |_{H_g}$
is a map from harmonics to harmonics. The $r$-valent
interaction vertex involves $r$ external legs from the fiber.
Insertion of the harmonic propagator contracts these legs
pairwise: each contraction through $P_{\mathrm{harm}}$
pairs two fiber directions and contributes cohomological
degree~$+2$. With $r$ legs, at most
$\lfloor r/2 \rfloor$ pairwise contractions are possible.
One contraction is already accounted for by the propagator
insertion that creates the discrepancy~$\delta_r$ itself.
The remaining $\lfloor r/2 \rfloor - 1$ contractions
contribute total cohomological degree
$2(\lfloor r/2 \rfloor - 1)$ to the harmonic-to-harmonic
map.

The degree count does \emph{not} imply that the residual
harmonic endomorphism is scalar. What it proves is the
divisibility statement: after the one propagator insertion
that creates~$\delta_r$, every remaining paired harmonic
contraction contributes a curvature factor. Thus there is a
central degree-zero endomorphism
$\epsilon_r(\cA)\in\operatorname{End}(H_g)$ such that
\[
\delta_r^{\mathrm{harm}}
\;=\;
\epsilon_r(\cA)\circ m_0^{\lfloor r/2 \rfloor - 1}.
\]
Identifying $\epsilon_r(\cA)$ with a scalar multiple of the
identity requires the separate vacuum-proportionality package
proved on the completed class-$\mathsf{M}$ surface in
Proposition~\ref{prop:central-m0-vacuum-proportionality}; it is not
assumed in this broad curvature-divisibility proposition.

\emph{Step~4: Low-degree evaluations.}
At $r = 2$: $\lfloor 2/2 \rfloor - 1 = 0$, so the scalar
channel of $\epsilon_2(\cA)$ is $c_2(\cA)$.
The binary discrepancy arises from replacing one bar
propagator by $P_{\mathrm{harm}}$. But the Hodge homotopy
$[d_{\mathrm{fib}}, h]$ already absorbs this replacement
at the binary level: the chain-homotopy equivalence between
$P_{\mathrm{BV}}$ and $d\log E$ at a single vertex
is the content of Step~1 applied at degree~$2$. Therefore
$c_2(\cA) = 0$.

At $r = 3$: $\lfloor 3/2 \rfloor - 1 = 0$, so the scalar
channel of $\epsilon_3(\cA)$ is $c_3(\cA)$.
For class~$\mathsf{G}$ \textup{(}Heisenberg\textup{)},
no cubic vertex exists, so $c_3 = 0$.
For class~$\mathsf{L}$ \textup{(}affine Kac--Moody at
$k \neq -h^\vee$\textup{)}, the cubic harmonic correction
takes the form
$\sum_c f^{abc} \cdot I_{\mathrm{harm}} \cdot f^{cde}$,
where $f^{abc}$ are the structure constants. By
antisymmetry of $f^{abc}$ and the Jacobi identity
$f^{ab[c}f^{d]be} = 0$, this coefficient vanishes:
$c_3(\cA) = 0$.

At $r = 4$: $\lfloor 4/2 \rfloor - 1 = 1$, so
$\delta_4^{\mathrm{harm}}$ is divisible by~$m_0$. The quartic
harmonic correction is the first term involving a nontrivial
curvature power. Its scalar/vacuum-channel coefficient
$c_4(\cA)$ is determined by the quartic shadow
invariant: $c_4(\cA) = S_4(\cA)$, the degree-$4$
projection of the Maurer--Cartan class
$\Theta_\cA$ in the convolution algebra.
For class~$\mathsf{G}$ and class~$\mathsf{L}$
\textup{(}shadow depth $r_{\max} \leq 3$\textup{)},
$S_4(\cA) = 0$ and thus $c_4(\cA) = 0$.
For class~$\mathsf{M}$ \textup{(}Virasoro\textup{)},
$S_4(\mathrm{Vir}_c) = 10/[c(5c+22)] \neq 0$
at generic~$c$, confirming that the quartic harmonic
discrepancy is nontrivial.

\emph{Step~5: Class-by-class vanishing.}
For class~$\mathsf{G}$: no interaction vertices at
any degree, so the harmonic discrepancy vanishes for all
$r \geq 2$.
For class~$\mathsf{L}$: shadow depth $r_{\max} = 3$
implies the scalar-channel coefficient $c_r(\cA) = 0$ for $r \geq 4$
\textup{(}no scalar shadow vertex survives at degree
$\geq 4$\textup{)}, and Jacobi kills $c_3$.
For class~$\mathsf{M}$: the shadow tower does not
terminate \textup{(}Corollary~\ref{cor:virasoro-postnikov-nontermination}\textup{)},
so the scalar-channel coefficients $c_r(\cA)$ are generically
nonzero for $r \geq 4$. The broad conclusion needed here is
curvature-divisibility: every harmonic discrepancy factors through
$m_0^{\lfloor r/2 \rfloor - 1}$, while scalar sharpness belongs to
the separate completed/vacuum-proportionality package.
\end{proof}

\begin{theorem}[BV$=$bar in the coderived category;
\ClaimStatusConditional]%
\label{thm:bv-bar-coderived-vol1}%
\index{BV algebra!bar complex identification!coderived category|textbf}%
\index{coderived category!BV/bar identification|textbf}%
Type signature: \textup{(}Open quadrant, BV/bar comparison
presentation, Beilinson level~$4$ compared with the level~$2$
bar coalgebra, hypothesis package: genus-$0$ geometric BV/bar
comparison, harmonic-factorization package of
Proposition~\textup{\ref{prop:harmonic-factorization}} in each genus,
three-obstruction chain-level analysis, curved weak-equivalence
ambient, and the coacyclic/coderived localization of
Definitions~\textup{\ref{def:coacyclic-fact}} and
\textup{\ref{def:coderived-fact}}\textup{)}.
Let~$\cA$ be a chirally Koszul algebra. The genus-$0$ BV/bar
comparison is the chain-level quasi-isomorphism of
Theorem~\ref{thm:bv-bar-geometric}. For each genus
$g \geq 1$, let
\[
 f_g \colon
 C^{\bullet}_{\mathrm{BV}}(\cA, \Sigma_g)
 \longrightarrow
 B^{(g)}(\cA)
\]
be the comparison morphism of filtered curved
factorization models
obtained by replacing the BV propagator with its Hodge
decomposition relative to the bar propagator. Then:
\begin{enumerate}[label=\textup{(\roman*)}]
\item for every genus $g \geq 0$,
\[
 C^{\bullet}_{\mathrm{BV}}(\cA, \Sigma_g)
 \;\simeq_{D^{\mathrm{co}}}\;
 B^{(g)}(\cA);
\]
\item if $\cA$ lies in class~$\mathsf{G}$ or~$\mathsf{L}$, then
 $f_g$ is a weak equivalence of filtered curved models in the sense
 of Definition~\ref{def:curved-weak-equiv};
\item if $\cA$ lies in class~$\mathsf{C}$ and harmonic decoupling
 holds, then $f_g$ is a weak equivalence of filtered curved models;
\item if $\cA$ lies in class~$\mathsf{M}$, then for every $r \geq 4$
 the harmonic discrepancy satisfies
\[
 \delta_r^{\mathrm{harm}}
 \;\in\;
 \operatorname{End}(H_g)\cdot m_0^{\lfloor r/2 \rfloor - 1}
\]
and the cone of $f_g$ is coacyclic.
\end{enumerate}
\end{theorem}

\begin{proof}
At genus~$0$, the comparison is the chain-level quasi-isomorphism
of Theorem~\ref{thm:bv-bar-geometric}.

For $g \geq 1$, Hodge decomposition writes the BV propagator as
\[
P_{\mathrm{BV}}
\;=\;
d\log E + [d_{\mathrm{fib}}, h] + P_{\mathrm{harm}}.
\]
The commutator term contributes only a chain homotopy, so the
failure of~$f_g$ to intertwine the two curved differentials is the
harmonic discrepancy
$\delta^{\mathrm{harm}} = \sum_{r \geq 4}\delta_r^{\mathrm{harm}}$.
Proposition~\ref{prop:harmonic-factorization} gives the
curvature-divisibility
$\delta_r^{\mathrm{harm}} \in \operatorname{End}(H_g)\cdot
m_0^{\lfloor r/2 \rfloor - 1}$, and
Proposition~\ref{prop:chain-level-three-obstructions} gives the
class-by-class description of the scalar-channel coefficients.

If $\cA$ lies in class~$\mathsf{G}$ or~$\mathsf{L}$, all harmonic
discrepancies vanish by
Proposition~\ref{prop:harmonic-factorization}(iii).
Hence $f_g$ intertwines the curved
differentials on the nose. On the PBW-associated graded, the map is
the genus-$0$ comparison on each weight piece, hence a
quasi-isomorphism. By Definition~\ref{def:curved-weak-equiv}, $f_g$
is a weak equivalence of filtered curved models. The same argument
applies to class~$\mathsf{C}$ under harmonic decoupling.

Now assume $\cA$ lies in class~$\mathsf{M}$, and let
$K_g := \operatorname{cone}(f_g)$. By
Proposition~\ref{prop:chain-level-three-obstructions}, every
filtration quotient of~$K_g$ is annihilated by a power of the
central curvature~$m_0$. If $m_0 = 0$, the same proposition gives
$\delta_r^{\mathrm{harm}} = 0$ for all $r \geq 4$, so there is no
class~$\mathsf{M}$ discrepancy. If $m_0 \neq 0$, each
$m_0$-power-torsion quotient is resolved by the one-variable
$m_0$-Koszul complex, whose totalization is exact; by
Definition~\ref{def:coacyclic-fact}, these totalizations are
coacyclic. Since the coacyclic subcategory is thick and closed
under extensions,~$K_g$ itself is coacyclic. Therefore~$f_g$
becomes an isomorphism in the coderived quotient
\[
D^{\mathrm{co}}(B^{(g)}(\cA)\text{-}\mathrm{CoFact})
\;=\;
\mathrm{Hot}(B^{(g)}(\cA)\text{-}\mathrm{CoFact})
\big/\mathrm{Acycl}^{\mathrm{co}}_{\mathrm{fact}}
\]
of Definition~\ref{def:coderived-fact}. This is the required
coderived comparison. The provisional localization is not needed
for the argument.
\end{proof}

\begin{remark}[Harmonic mechanism behind the coderived comparison]
\label{rem:bv-bar-coderived-higher-genus}%
\index{coderived category!higher-genus validity}%
The curvature-divisibility relation
$\delta_r^{\mathrm{harm}} \in \operatorname{End}(H_g)\cdot
m_0^{\lfloor r/2 \rfloor - 1}$
\textup{(}Proposition~\ref{prop:harmonic-factorization}\textup{)}
is the genus-independent mechanism behind
Theorem~\ref{thm:bv-bar-coderived-vol1}. The proof has two inputs:
\begin{enumerate}[label=\textup{(\roman*)},nosep]
\item Hodge decomposition on the fiber model isolates the harmonic
 part of the discrepancy, and the harmonic projector commutes with
 powers of the central curvature by degree counting on the harmonic
 summand.
\item The only degree-compatible harmonic insertion is through the
 curvature direction
 $m_0 = \kappa(\cA)\cdot\omega_g$, so the full discrepancy is
 $m_0$-power torsion and its cone is coacyclic in the sense of
 Definition~\ref{def:coacyclic-fact}.
\end{enumerate}
What remains open is the stronger vanishing of the scalar-channel
coefficients $c_r(\cA)$ in class~$\mathsf{M}$ and the all-genera
verification of harmonic decoupling in class~$\mathsf{C}$.
\end{remark}

\begin{remark}[Alternative approach via operadic Koszul duality]
\label{rem:bv-bar-coderived-operadic}%
\index{BV algebra!bar complex identification!operadic route}%
\index{Koszul duality!BV/bar comparison}%
There is a second route to
Theorem~\ref{thm:bv-bar-coderived-vol1}. By
Principle~\ref{princ:sc-two-incarnations}, the ordered bar complex
$\barB^{\mathrm{ord}}(\cA)$ is the open
$\Eone$-chiral coalgebra engine, while the derived center
$Z^{\mathrm{der}}_{\mathrm{ch}}(\cA)
= C^{\bullet}_{\mathrm{ch}}(\cA,\cA)$ is the closed object carrying
the $\mathsf{SC}^{\mathrm{ch,top}}$ structure. In this two-coloured
package, the BV operator $\Delta_{\mathrm{BV}}$ is the closed-colour
contraction induced by the Swiss-cheese pairing, whereas the bar
differential $d_{\barB}$ is the open-colour coderivation induced by
the chiral product. The two operators therefore come from the same
operadic bar-cobar datum, but they live in different colours. What is
compared is not ``the bar complex as a Swiss-cheese algebra'': the
Swiss-cheese object is the pair
\[
\bigl(Z^{\mathrm{der}}_{\mathrm{ch}}(\cA),\,\cA\bigr),
\]
and the ordered bar complex enters as the open-colour coalgebra model
used to resolve the boundary algebra.

There is also a conditional operadic derivation of the coderived
comparison. Its input is not classical Swiss-cheese Koszulity by
itself. One must import the full
$\mathsf{SC}^{\mathrm{ch,top}}$ homotopy-Koszulity package recorded in
Remark~\ref{rem:homotopy-koszulity-center}: classical Swiss-cheese
Koszulity, the fixed formality comparison from the chiral closed
colour to the classical $E_2$ closed colour, and transfer through
bar-cobar quasi-isomorphisms. Together with the explicit description
of the Koszul dual cooperad in
Proposition~\ref{prop:sc-koszul-dual-three-sectors} and the center
identification of Theorem~\ref{thm:operadic-center-hochschild}, this
package gives a second proof route: the general bar-cobar
correspondence for Koszul coloured operads identifies the
$\mathsf{SC}^{\mathrm{ch,top}}$-algebra pair
\[
\bigl(Z^{\mathrm{der}}_{\mathrm{ch}}(\cA),\,\cA\bigr)
\]
with its
$\mathsf{SC}^{\mathrm{ch,top},!}$-coalgebra resolution in the coderived
category. On the closed side this resolution computes the BV operator
through the Swiss-cheese contraction; on the open side it computes the
bar differential through the cofree ordered bar coalgebra. Thus the
BV/bar comparison in $D^{\mathrm{co}}$ is the Swiss-cheese instance of
operadic Koszul duality. From this viewpoint,
Theorem~\ref{thm:bv-bar-coderived-vol1} is, conditional on that
imported homotopy-Koszulity package, the Swiss-cheese instance of the
general curved Koszul principle: off the strict locus, curvature
forces passage to $D^{\mathrm{co}}$, but the comparison cone is
invisible there.

This also isolates the class~$\mathsf{M}$ gap in operadic terms. The
coderived equivalence survives because the operadic quasi-isomorphism
still exists after passing to the Koszul resolution, but a strict
chain-level quasi-isomorphism on the explicit models requires the
transferred $\mathsf{SC}^{\mathrm{ch,top}}$ structure to collapse to a
formal one. That collapse holds for classes~$\mathsf{G}$ and
$\mathsf{L}$, is the harmonic-decoupling input in class~$\mathsf{C}$,
and is not available for class~$\mathsf{M}$, where the
higher transferred operations remain visible. The class~$\mathsf{M}$
obstruction is therefore not a failure of coderived Koszul duality; it
is the failure of the operadic quasi-isomorphism to descend to a
strict chain-level quasi-isomorphism on a non-formal explicit model.
\end{remark}

\begin{remark}[Why the coderived category is inevitable]
\label{rem:bv-bar-coderived-why}
\index{coderived category!physical inevitability}
\index{BV algebra!bar complex identification!physical interpretation}
The passage from the ordinary derived category to the coderived
category is not a technical convenience but a physical necessity,
and its inevitability can be understood from three complementary
viewpoints.

\emph{The worldsheet viewpoint.}
On a genus-$g$ Riemann surface $\Sigma_g$, the worldsheet path
integral produces amplitudes via the sewing construction: one cuts
$\Sigma_g$ along a collar into two pieces and sums over a complete
set of intermediate states. The sewing sum introduces the
propagator $P = P_{\mathrm{bar}} + P_{\mathrm{exact}} +
P_{\mathrm{harm}}$, where $P_{\mathrm{harm}}$ is the harmonic
piece from $H^{0,1}(\Sigma_g)$. The BV formalism uses the
\emph{full} propagator; the bar complex uses only
$P_{\mathrm{bar}} = d\log E(z,w)$. The discrepancy at degree~$r$
is controlled by the harmonic piece. This shows why curvature
appears in higher genus: the harmonic sector is invisible to the
genus-$0$ bar differential but survives in the BV propagator.
Passing to the coderived setting is therefore natural once
$d^2 = m_0 \cdot \mathrm{id}$ enters the picture. Conditional on the
harmonic-factorization package, Theorem~\ref{thm:bv-bar-coderived-vol1}
supplies the missing step: curvature factorization forces the
comparison cone to be coacyclic.

\emph{The spacetime viewpoint.}
In the Costello--Li framework of twisted supergravity, the bulk
theory on $\mathrm{AdS}_{d+1}$ is computed by Witten diagrams
using the bulk propagator, while the boundary theory is computed
by the bar complex using the boundary OPE. The mismatch between
bulk and boundary propagators at one loop ($g = 1$) is the
holographic anomaly: the Weyl anomaly coefficient
$a - c$ in $4d$ becomes, on the algebraic side, a harmonic
correction to the BV/bar comparison. The coderived formalism
packages this anomaly as curvature data rather than forcing a
strict differential. Under the same harmonic-factorization package,
Theorem~\ref{thm:bv-bar-coderived-vol1} shows that the comparison map
becomes an isomorphism after localizing at coacyclic objects.

\emph{The categorical viewpoint.}
The coderived category $D^{\mathrm{co}}(\cA)$ is the natural
home for curved dg~algebras: it is the homotopy category of
dg~comodules over the bar coalgebra $B(\cA)$ in which
$d^2 = m_0 \cdot \mathrm{id}$ is permitted. The bar complex
$B(\cA)$ always satisfies $D_B^2 = 0$ (the total bar differential
squares to zero regardless of curvature), so $B(\cA)$ lives
natively in $D^{\mathrm{co}}$. The BV complex, with its
$d^2_{\mathrm{BV}} = \kappa \cdot \omega_g$ at genus~$g$,
satisfies the same curved relation. The coderived
category is the minimal ambient in which both objects are
well-defined. The ordinary derived
category, which requires $d^2 = 0$, is too restrictive: it
forces a choice of ``strictification'' (lifting the curvature
into the differential via a Maurer--Cartan twist), and different
choices produce inequivalent chain complexes. The coderived
category avoids this choice. By
Definitions~\ref{def:coacyclic-fact} and~\ref{def:coderived-fact},
curved objects are not set to zero in this localization; the
comparison map of Theorem~\ref{thm:bv-bar-coderived-vol1} becomes an
isomorphism precisely because its cone is coacyclic.

Thus the BV quantisation of the worldsheet sigma model and the
algebraic bar construction of the boundary vertex algebra live in the
same curved homological framework. The chain-level discrepancy for
class~$\mathsf{M}$ is therefore not a signal that the coderived
category is unnecessary; it is the reason the coderived category is
the correct ambient. What remains open is
the filtered-completed refinement in class~$\mathsf{M}$
\textup{(}Remark~\ref{rem:bv-bar-class-m-frontier}\textup{)} and the
all-genera harmonic-decoupling statement in class~$\mathsf{C}$.
\end{remark}

\begin{remark}[Class~$\mathsf{M}$ spectral and completed MC5]
\label{rem:bv-bar-class-m-frontier}
\index{BV algebra!class M frontier|textbf}%
\index{coderived category!filtered-completed refinement}%
Theorem~\ref{thm:bv-bar-coderived-vol1} and
Proposition~\ref{prop:chain-level-three-obstructions} isolate not
only the surviving coderived statement but also the exact reason the
strict chain-level form of MC5 fails for class~$\mathsf{M}$. For
Virasoro and principal $\cW_N$, the higher harmonic discrepancies
satisfy
\[
\delta_r^{\mathrm{harm}}
\;\in\;
\operatorname{End}(H_g)\cdot m_0^{\lfloor r/2 \rfloor - 1}
\qquad (r \geq 4).
\]
Corollary~\ref{cor:virasoro-postnikov-nontermination} shows that the
shadow obstruction tower does not terminate: the permanent cubic
source regenerates nonzero higher obstructions at every stage. In the
operadic language of
Remark~\ref{rem:bv-bar-coderived-operadic}, the transferred
class~$\mathsf{M}$ $A_\infty$-operations do not truncate. A strict
comparison on the raw direct-sum models would have to absorb
infinitely many nonvanishing higher corrections by a single ordinary
chain-homotopy package. This is what fails. The coderived theorem
survives because every defect is still divisible by the central
curvature.

A spectral replacement remains viable. Let
$K_g = \operatorname{cone}(f_g)$ and filter it by bar degree, total
conformal weight, or, more sharply for class~$\mathsf{M}$, by a
shadow-window filtration in which the quotient by $F^{R+1}$ keeps only
the discrepancy classes of shadow degree at most~$R$. Each finite
window then sees only finitely many nonzero operations, so the
associated-graded comparison becomes a finite problem.
Proposition~\ref{prop:coderived-bar-degree-spectral-sequence}
supplies the corresponding coderived spectral sequence, with strict
$E_1$ page because the curvature lies in positive filtration. The
research target is a spectral form of MC5 whose higher differentials
record the scalar-channel harmonic classes $c_r(\cA)$ and whose abutment is a
coacyclic cone. This replaces a false strict chain identity by a
filtered comparison that measures the class~$\mathsf{M}$ defect page
by page.

The strongest chain-level statement presently visible is completed
rather than strict. Proposition~\ref{prop:standard-strong-filtration}
and Theorem~\ref{thm:completed-bar-cobar-strong} already provide the
weight-completed bar side, and
Corollary~\ref{cor:virasoro-postnikov-nontermination} shows why
class~$\mathsf{M}$ also needs completion in shadow degree. Since the
full Maurer--Cartan class satisfies
$\Theta_\cA = \varprojlim_r \Theta_\cA^{\le r}$ in the completed
convolution algebra, the BV analogue belongs in the same completed
regime. The natural target is a continuous comparison of
bi-completed models
\[
\widehat f_g \colon
\widehat C^{\bullet}_{\mathrm{BV}}(\cA,\Sigma_g)
\longrightarrow
\widehat B^{(g)}(\cA),
\]
where the hats denote inverse limits over finite weight truncations
and finite shadow windows, equivalently a pro-object built from those
finite quotients. The conjectural statement is not that the raw
direct-sum complexes are quasi-isomorphic, but that each finite
quotient of $\widehat f_g$ is a strict quasi-isomorphism and that the
inverse-limit comparison is continuous and filtration-compatible. This
is stronger than Theorem~\ref{thm:bv-bar-coderived-vol1}, because it
remembers every finite stage, and weaker than a raw chain-level
quasi-isomorphism, because it allows the infinite class~$\mathsf{M}$
defect to die only after completion.

Between ordinary chain level and bare coderived equivalence lies an
intermediate strengthening. One may ask that $\widehat f_g$ be a weak
equivalence of filtered curved models in the sense of
Definition~\ref{def:curved-weak-equiv}, equivalently an isomorphism
already in the provisional coderived category of
Definition~\ref{def:provisional-coderived}. This retains the full
filtered-completed class~$\mathsf{M}$ tower while remaining compatible
with the proved coderived theorem. With the available hypotheses, this
filtered-completed spectral statement is the sharpest plausible
upgrade of MC5 for class~$\mathsf{M}$.
\end{remark}

\section{Non-Calabi--Yau local surfaces and the Burns datum}
\label{sec:non-cy-local-surfaces}
\index{Burns space|textbf}
\index{local surface!non-Calabi--Yau}

The holomorphic-topological boundary VOAs of twisted supergravity
on non-Calabi--Yau local surfaces furnish a family of class~$\mathsf{C}$
algebras whose Koszul datum is controlled by a single integer
\textup{(}the rank of the characteristic lattice\textup{)}. The Burns
space is the simplest member.

\begin{computation}[Burns space Koszul datum;
\ClaimStatusProvedElsewhere]
\label{comp:v1-burns-koszul-datum}
\index{Burns space!Koszul datum}
\index{$\beta\gamma$ system!Burns}
The boundary VOA of the Burns space consists of four $\beta\gamma$
pairs with an $SO(8)$ global symmetry, giving:
\begin{itemize}[nosep]
\item central charge $c_{\mathrm{Burns}}
 = 4 \cdot c_{\beta\gamma}(\lambda{=}1)
 = 4 \cdot 2 = 8$,
 using $c_{\beta\gamma}(1) = 2(6 - 6 + 1) = 2$;
\item modular characteristic
 $\kappa_{\mathrm{Burns}} = c_{\mathrm{Burns}}/2 = 4$;
\item shadow class~$\mathsf{C}$ globally \textup{(}shadow depth
 $r_{\max} = 4$, governed by the quartic contact term of the
 $\beta\gamma$ system\textup{)}, degenerating to class~$\mathsf{M}$
 \textup{(}$r_{\max} = \infty$\textup{)} along the $T$-line, where
 the stress tensor becomes a fundamental generator;
\item genus-$1$ obstruction
 $F_1^{\mathrm{Burns}} = \kappa_{\mathrm{Burns}}/24 = 4/24 = 1/6$;
\item genus-$2$ obstruction
 $F_2^{\mathrm{Burns}} = \kappa_{\mathrm{Burns}} \cdot
 \lambda_2^{\mathrm{FP}} = 4 \cdot 7/5760 = 7/1440$.
\end{itemize}
Both $F_1$ and $F_2$ are verified by multi-path computation in
\texttt{compute/lib/burns\_space\_koszul\_datum\_engine.py} and the
associated test file
\texttt{compute/tests/test\_theorem\_burns\_f2\_engine.py}
\textup{(}five paths for $\kappa$, six paths for $F_2$\textup{)}.
\end{computation}

\begin{remark}[Burns shadow class, global vs.\ $T$-line]
\label{rem:burns-shadow-class}
On the generic locus of the Burns parameter space, the only
propagating OPE channel is the $\beta\gamma$ binary OPE with a
simple pole, and the quartic contact term arises from the composite
stress tensor $T = {:}\!\beta\partial\gamma\!{:}$; this is the class
$\mathsf{C}$ mechanism of
Remark~\ref{rem:bv-bar-class-c-proof}. Along the $T$-line, $T$
couples to itself as a fundamental generator, so the shadow depth
jumps from $r_{\max} = 4$ to $r_{\max} = \infty$: the Burns datum
degenerates to class~$\mathsf{M}$. The genus expansion on the
$T$-line no longer closes at the scalar level, and the
multi-weight cross-channel correction $\delta F_g^{\mathrm{cross}}$
becomes nonzero at $g \geq 2$. The values
$F_1 = 1/6$ and $F_2 = 7/1440$ are computed on the class-$\mathsf{C}$
generic locus, not on the $T$-line.
\end{remark}

\begin{table}[htbp]
\caption{Non-Calabi--Yau local surfaces: Koszul datum
$(c, \kappa, \mathrm{class}, F_1, F_2)$. Burns is
\ClaimStatusProvedElsewhere{} \textup{(}verified in the compute
layer\textup{)}; the remaining rows are \ClaimStatusConjectured{}
until their \(c,\kappa,F_1,F_2\) values are computed directly.}
\label{tab:non-cy-local-surfaces}
\begin{center}
\small
\renewcommand{\arraystretch}{1.20}
\begin{tabular}{lccccc}
\toprule
Surface & $c$ & $\kappa$ & Class & $F_1$ & $F_2$ \\
\midrule
Burns space \textup{(}$4\beta\gamma$, $SO(8)$\textup{)}
 & $8$ & $4$ & $\mathsf{C}$
 & $1/6$ & $7/1440$ \\[1pt]
Hirzebruch $\mathbb{F}_n$ \textup{(}conj.\textup{)}
 & $2n$ & $n$ & $\mathsf{C}$
 & $n/24$ & $n \cdot 7/5760$ \\[1pt]
del Pezzo $\mathrm{dP}_k$ \textup{(}$k=0,\ldots,8$; conj.\textup{)}
 & $2(k{+}1)$ & $k{+}1$ & $\mathsf{C}$
 & $(k{+}1)/24$ & $(k{+}1) \cdot 7/5760$ \\[1pt]
Enriques surface \textup{(}conj.\textup{)}
 & $22$ & $11$ & $\mathsf{C}$
 & $11/24$ & $11 \cdot 7/5760$ \\
\bottomrule
\end{tabular}
\end{center}
\end{table}

\begin{conjecture}[Hirzebruch, del Pezzo, and Enriques Koszul data;
\ClaimStatusConjectured]
\label{conj:non-cy-local-surfaces}
\index{Hirzebruch surface!Koszul datum}
\index{del Pezzo surface!Koszul datum}
\index{Enriques surface!Koszul datum}
The boundary VOAs of twisted supergravity on the non-Calabi--Yau
local surfaces listed in Table~\textup{\ref{tab:non-cy-local-surfaces}}
have the Koszul data shown there. Concretely:
\begin{enumerate}[label=\textup{(\roman*)},nosep]
\item \emph{Hirzebruch $\mathbb{F}_n$:} the characteristic lattice
 $H^{2}(\mathbb{F}_n, \mathbb{Z})$ has rank~$2$, so the boundary
 VOA decomposes into $n$ isomorphic $\beta\gamma$ pairs
 \textup{(}one pair per self-intersection unit along the fibre
 direction\textup{)}, giving
 $\kappa(\mathbb{F}_n) = n$ linearly in the Hirzebruch parameter.
\item \emph{del Pezzo $\mathrm{dP}_k$:} the dimension
 $\dim H^{2}(\mathrm{dP}_k, \mathbb{R}) = k + 1$ controls the
 number of $\beta\gamma$ pairs in the boundary VOA, giving
 $\kappa(\mathrm{dP}_k) = k + 1$. At $k = 0$ \textup{(}$\mathbb{P}^2$,
 $\kappa = 1$\textup{)} and $k = 8$ \textup{(}maximal blow-up,
 $\kappa = 9$\textup{)} the numerics are predicted to match the
 characteristic-lattice dimension prediction.
\item \emph{Enriques surface:} a $\mathbb{Z}/2$ quotient of a K3 surface,
 and the K3 boundary VOA carries $\kappa(\mathrm{K3}) = 24$
 \textup{(}the $24$-dimensional Niemeier lattice\textup{)}. The
 quotient halves the characteristic lattice and loses a single unit
 from the fixed locus contribution, giving
 $\kappa(\mathrm{Enriques}) = 11$ rather than~$12$.
\end{enumerate}
In all three cases, the shadow class is conjectured to be
$\mathsf{C}$ on the generic locus \textup{(}four-pair $\beta\gamma$
structure\textup{)} and to degenerate to $\mathsf{M}$ along the
stress-tensor line, in direct analogy with the Burns space
\textup{(}Remark~\textup{\ref{rem:burns-shadow-class}}\textup{)}.
Direct verification via multi-path computation is open.
\end{conjecture}

\begin{remark}[Scope and status]
\label{rem:non-cy-scope}
The Burns row of Table~\ref{tab:non-cy-local-surfaces} is the only
entry that is presently computed from the engine; it is
\ClaimStatusProvedElsewhere{} by the five-path $\kappa$ verification
and the six-path $F_2$ verification in the compute layer. The
Hirzebruch, del~Pezzo, and Enriques rows are
\ClaimStatusConjectured{} predictions based on dimensional counting
of the characteristic lattice and the general class-$\mathsf{C}$
pattern $F_g = \kappa(\cA) \cdot \lambda_g^{\mathrm{FP}}$ on the
uniform-weight lane of Theorem~D.
No claim is made that these families are uniformly class~$\mathsf{C}$
on their full parameter spaces: the analogous jump to class
$\mathsf{M}$ along codimension-one loci for non-Calabi--Yau local
surfaces with composite stress tensor is a conjectural pattern, and it
requires a systematic discriminant and shadow-tower computation.
\end{remark}

\section{Affine WZW bar-BRST addendum}

\begin{proposition}[Genus-\texorpdfstring{$0$}{0} WZW BRST complex from the affine bar complex;
\ClaimStatusConditional]
\label{prop:wzw-brst-bar-genus0}
\index{WZW model!BRST complex as bar complex}
\index{BRST!WZW genus-0 bar comparison}
Let $\fg$ be a simple Lie algebra with dual Coxeter number~$h^\vee$,
and let $k \notin \Sigma(\fg) \cup \{-h^\vee\}$ be generic
\textup{(}in particular $k \neq -h^\vee$, so that the Sugawara
construction is well-defined\textup{)}. The affine vertex algebra
$V_k(\fg)$ carries the classical $r$-matrix
\begin{equation}\label{eq:wzw-r-matrix-brst}
r^{\mathrm{KM}}(z) = k\,\Omega_{\mathrm{tr}}/z,
\end{equation}
where $\Omega \in \fg \otimes \fg$ is the Casimir tensor, and the
curvature coefficient
\begin{equation}\label{eq:wzw-kappa-brst}
\kappa\bigl(V_k(\fg)\bigr)
= \frac{\dim(\fg)\,(k + h^\vee)}{2h^\vee}.
\end{equation}
Write
\[
 C_{\mathrm{BRST}}^{\bullet,\mathrm{WZW}}(\fg,k)
 \;:=\;
 \bigl(
 C^{\infty/2+\bullet}(\widehat{\fg}_k, V_k(\fg)),\,
 Q_{\mathrm{BRST}}^{\mathrm{WZW}}
 \bigr)
\]
for the standard semi-infinite BRST complex of the genus-$0$ WZW
model, and write
\[
 \barB^{\mathrm{ch}}_0\bigl(V_k(\fg)\bigr)
 \;=\;
 T^c\bigl(s^{-1}\overline{V_k(\fg)}\bigr),
 \qquad
 \overline{V_k(\fg)} := \ker(\epsilon),
\]
for the genus-$0$ chiral bar complex, with desuspension
$|s^{-1}v| = |v| - 1$. Then for every degree $n \geq 3$
\textup{(}the stability range $2g - 2 + n > 0$ at genus~$0$\textup{)},
the WZW BRST complex is computed by the affine bar complex: there
is a natural filtered quasi-isomorphism
\begin{equation}\label{eq:wzw-brst-bar-qi}
 C_{\mathrm{BRST}}^{\bullet,\mathrm{WZW}}(\fg,k)
 \xrightarrow{\;\sim\;}
 \bigl(
 \barB^{\mathrm{ch}}_0\bigl(V_k(\fg)\bigr),\,
 d_{\mathrm{bar}}
 \bigr).
\end{equation}
On the PBW-associated graded, the BRST differential
$Q_{\mathrm{BRST}}^{\mathrm{WZW}}$ equals the bar differential
$d_{\mathrm{bar}}$; both reduce to the Chevalley--Eilenberg
differential $d_{\mathrm{CE}}$ of the loop algebra subalgebra
$\widehat{\fg}_{k,-} := \fg \otimes t^{-1}\mathbb{C}[t^{-1}]$
acting on the vacuum module~$V_k(\fg)$.
\end{proposition}

\begin{proof}
The argument extends
Theorem~\ref{thm:brst-bar-genus0} \textup{(}the $c = 26$ string
case\textup{)} by replacing the Virasoro BRST complex with the
affine semi-infinite complex, and invokes
Theorem~\ref{thm:bar-semi-infinite-km} for the comparison.

\emph{Step~1: Bar--semi-infinite comparison.}
Theorem~\ref{thm:bar-semi-infinite-km} provides a filtered
quasi-isomorphism between the affine bar complex
$\barB^{\mathrm{ch}}(\widehat{\fg}_k)$ and the semi-infinite complex
$C^{\infty/2+\bullet}(\widehat{\fg}_k, V_k)$ for every
$k \neq -h^\vee$. The semi-infinite complex is the standard
genus-$0$ BRST presentation of the WZW model, so it remains to
identify the differentials on the PBW-associated graded.

\emph{Step~2: PBW filtration and Chevalley--Eilenberg reduction.}
At generic level, Theorem~\ref{thm:km-chiral-koszul} and
Proposition~\ref{prop:km-generic-acyclicity} place $V_k(\fg)$ on the
chirally Koszul locus. On the BRST side, the associated graded
differential is the Chevalley--Eilenberg differential~$d_{\mathrm{CE}}$
of the current Lie algebra $\widehat{\fg}_{k,-}$ acting on the
vacuum module. On the bar side, the associated graded differential
is extracted from the current OPE
$J^a(z) J^b(w) \sim k\,\Omega^{ab}/(z-w)^2 + f^{ab}{}_c J^c(w)/(z-w)$,
which is the same bracket that defines~$d_{\mathrm{CE}}$.
The filtered comparison therefore lifts the Chevalley--Eilenberg
identification to the full genus-$0$ complexes.

\emph{Step~3: Degree propagation.}
Since $V_k(\fg)$ is chirally Koszul at generic level, the bar spectral
sequence collapses on the Koszul line and each finite-length summand of
$T^c(s^{-1}\overline{V_k(\fg)})$ is generated by the binary residue
bracket and its iterates. The genus-$0$ BRST/bar identification
propagates from the binary comparison to every finite degree $n \geq 3$.
\end{proof}

\begin{conjecture}[All-genera WZW BRST/bar identification at generic level;
\ClaimStatusConjectured]
\label{conj:wzw-brst-bar-all-genera}
\index{WZW model!all-genera BRST/bar identification}
Let $\fg$ be simple and let
$k \notin \Sigma(\fg) \cup \{-h^\vee\}$ be generic. For every
genus $g \geq 0$ and every degree $n \geq 1$, the renormalized
BV-BRST complex of the genus-$g$ WZW model with $n$ marked points is
filtered quasi-isomorphic to the genus-$g$, degree-$n$ bar complex of
$V_k(\fg)$. At $g=0$ this is
Proposition~\ref{prop:wzw-brst-bar-genus0}; the remaining content is
the higher-genus MC5 comparison between the handle-gluing BRST
operator and the modular bar differential.
\end{conjecture}

\begin{remark}[Generic level, critical level, and the free-string limit]
\label{rem:wzw-brst-bar-generic-level}
The generic-level hypothesis is used twice. First, the proof of
Proposition~\ref{prop:wzw-brst-bar-genus0} uses the chirally Koszul
package of Theorem~\ref{thm:km-chiral-koszul} and
Proposition~\ref{prop:km-generic-acyclicity} to force the PBW
spectral sequence onto its Chevalley--Eilenberg line. Second, the
generic affine center is the small one used by the filtration
argument. At the critical level $k=-h^\vee$, the scalar class
$\kappa(V_k(\fg))$ vanishes, the Sugawara grading used in the generic
argument is replaced by the separate critical package, and the
Feigin--Frenkel center jumps
\textup{(}Theorem~\ref{thm:critical-level-structure}\textup{)}.
The proposition is therefore stated only on the generic lane.

The affine gauge ghosts contribute one fermionic $bc$ pair for each
basis vector of~$\fg$, hence
\[
 c_{\mathrm{ghost}}^{\mathrm{WZW}} = -2\dim(\fg).
\]
This is the same numerical contribution carried by the odd
Koszul-dual generators in the Chevalley--Eilenberg bar model.

For $\fg = \mathfrak{u}(1)$, the matter algebra is Heisenberg, so the
genus-$0$ comparison reduces to the abelian bar/BRST identification
already used in Theorem~\ref{thm:brst-bar-genus0}. Taking $26$
copies recovers the free matter sector $\mathcal{H}^{26}$ of the
bosonic string. Coupling that abelian matter to the reparametrizing
$bc$ ghosts gives the MC5 genus-$0$ case proved in
Theorem~\ref{thm:brst-bar-genus0} and
Corollary~\ref{cor:anomaly-physical-genus0-vol1}. The affine ghost count
$c_{\mathrm{ghost}}^{\mathrm{WZW}} = -2\dim(\fg)$ belongs to the WZW
gauge-theory sector; the string-theoretic ghost count
$c_{bc} = -26$ is a separate reparametrization contribution.
\end{remark}

\section{BV-BRST and the K3 1-loop anomaly}
\label{sec:bvbrst-supplement}

\begin{remark}[Holomorphic Chern--Simons and the five-dimensional
$\Omega$-background; \ClaimStatusProvedElsewhere]
\label{rem:bvbrst-6d-hcs-quartic}
\index{holomorphic Chern--Simons!anomaly complex|textbf}
Costello--Li construct the large-$N$ coupling of holomorphic
Chern--Simons theory to Kodaira--Spencer gravity and prove
perturbative anomaly cancellation for the coupled open--closed
topological B-model \cite{CL20}.  Their type-I models select gauge
group $\mathrm{SO}(8)$ in complex dimension~$3$ and
$\mathrm{SO}(32)$ in complex dimension~$5$.

Costello constructs a second theory by reducing eleven-dimensional
supergravity in a particular $\Omega$-background to
five-dimensional non-commutative gauge theory on
$\mathbb R\times X$, with $X$ a conical holomorphic symplectic
surface \cite{Costello2111}.  Theorems~9.0.2--9.0.3 of that paper give
the perturbative quantization and its deformation parameters.

A dimensional comparison is therefore a morphism
\[
 \rho_{\Omega}^{\mathrm{obs}}:
 H^1\operatorname{Obs}_{\mathrm{loc}}
   (\mathrm{hCS}_{Z_3}\oplus\mathrm{BCOV}_{Z_3})
 \longrightarrow
 H^1\operatorname{Obs}_{\mathrm{loc}}
   (\mathcal E^{\mathrm{nc}}_5(\mathbb R\times X))
\]
induced by a specified reduction of fields and interactions.  Its
effect on invariant polynomials determines the relation between the
open--closed anomaly class and the five-dimensional quantization
class.  Constructing $\rho_{\Omega}^{\mathrm{obs}}$ is the exact
cross-dimensional problem.
\end{remark}

\begin{remark}[Automorphic line and the proposed BV comparison;
\ClaimStatusConjectured]
\label{rem:bvbrst-paramodular-anomaly}
The Gritsenko--Nikulin form $\Delta _5$ is a section of
$\lambda^5\otimes\chi$ on $\mathcal A_2$ with divisor~$H_1$
\cite{GN96Igusa}.  Its weight is
\[
 \kappa_{\mathrm{BKM}}(\Delta _5)=\frac{f(0)}2=\frac{10}{2}=5.
\]
For a specified K3 background, a field-theoretic interpretation
requires a morphism from the one-loop local BV obstruction group to
$H^0(\mathcal A_2,\lambda^5\otimes\chi)$ carrying the obstruction
class to~$\Delta _5$.  The sourced twisted action, its quantization,
and this morphism constitute the conjectural comparison.  The
arithmetic line and its Kummer root classes are developed below.
\end{remark}

\begin{remark}[K3 conductor carriers: ghost, Mukai, and Borcherds;
\ClaimStatusHeuristic]
\label{rem:bvbrst-koszul-conductor}
The K3 geometry contains three conductors with different carriers. The
BRST ghost-resolution conductor is
$K_{\mathrm{univ}}^{\mathrm{BRST}}(A)
=-\tfrac12c_{\mathrm{ghost}}(\BRST(A))$ after choosing a ghost system.
The Mukai-enhanced Theorem~C conductor is
\[
 K^{\kappa_{\mathrm{ch}}}(\mathcal{H}_{\mathrm{Muk}})
 =2c_+(\mathrm{Mukai}(K3))=8,
\]
the $\mathcal B$-family value that extends the Theorem~C scalar set to
$\{0,8,13,250/3,25/3\}$. The Borcherds denominator conductor is
\[
 \kappa_{\mathrm{BKM}}(\mathfrak g_{\Delta_5})=f(0)/2=5,
\]
where $f(0)=10$ is the constant coefficient of the weight-$0$,
index-$1$ Jacobi form in the Gritsenko--Nikulin product
\cite{GN96Igusa}. Bruinier's Heegner--Chern reciprocity supplies an
automorphic divisor target~\cite[Proposition 5.1]{Bruinier2002}.
The three conductors have distinct carriers; an identification between
them is a separate comparison morphism. The
specialisation $\hbar_{\mathrm{qg}}^2K^{\kappa_{\mathrm{ch}}}=-1$ and
the Lusztig order $\ell=8$ require the separate monodromy, PBW
truncation, parity, and bulk-boundary comparison data recorded in the
K3 concordance. They are comparison targets for the corresponding
Vol.~I scalar morphisms.
\end{remark}

\section{The two-loop BV comparison}
\label{sec:bvbrst-supplement-ii}

Fix a global BV theory $\mathcal E_{11}^{\mathrm{tw}}(Y)$, classical
interaction~$I_0$, gauge fixing, and renormalization scheme.  The first
two obstruction cocycles are
\[
 \mathfrak o_1=\Delta_{\mathrm{BV}}I_0,\qquad
 \mathfrak o_2=\Delta_{\mathrm{BV}}I_1
 +\frac12\{I_1,I_1\}_{\mathrm{BV}}
 \in Z^1\operatorname{Obs}_{\mathrm{loc}}
        (\mathcal E_{11}^{\mathrm{tw}}(Y)).
\]

\begin{conjecture}[Two-loop BV/Fourier comparison;
\ClaimStatusConjectured]
\label{prop:bvbrst-2loop-obstruction}
For
\[
 \widetilde\Phi_{10}:=\frac{\Phi_{10}}{pqy},\qquad
 \Phi_{10}^{\circ}:=
 \frac{\widetilde\Phi_{10}(p,q,y)}
      {\widetilde\Phi_{10}(0,q,y)},
\]
construct a graded map
\[
 \Psi_2:\ H^1\operatorname{Obs}_{\mathrm{loc}}
        (\mathcal E_{11}^{\mathrm{tw}}(Y))
 \longrightarrow\bQ(y)[[q]]
\]
such that
\[
 \Psi_2([\mathfrak o_2])
 =[p^2]\log\Phi_{10}^{\circ}(p,q,y).
\]
This equation is the two-loop comparison statement.
\end{conjecture}

\begin{problem}[Theta-graph amplitude; \ClaimStatusOpen]
\label{rem:bvbrst-theta-graph}
Let $\Theta$ be a decorated theta graph of the chosen theory.  Compute
\[
 A_\Theta^{\mathrm{ren}}
 =\operatorname{FP}_{\epsilon\to0}
 \int_{\operatorname{Conf}_2(Y)}
 I_{0,v_1}I_{0,v_2}\,
 P_{\epsilon}^{\wedge3}
 +A_\Theta^{\mathrm{ct}}.
\]
The identity $\int_{K3}c_2(TK3)=24$ supplies a topological input to
this integral; the vertices, propagator, and counterterm determine its
field-theoretic coefficient.
\end{problem}

\begin{remark}[Geometric scope]
\label{rem:bvbrst-scope}
A compact K3 comparison may be posed on
$Y=\mathbb R\times K3\times E\times\mathbb C^2$, together with
specified M5 source worldvolumes.  The proof requires a sourced global
action, a solution of the QME through order~$2$, and the map~$\Psi_2$.
\end{remark}

\begin{remark}[Two deformation parameters]
\label{rem:bvbrst-two-hbar}
The BV parameter $\hbar_{\mathrm{BV}}$ records loop order in the
completed local-functional algebra.  A quantum-group parameter
$\hbar_{\mathrm{qg}}$ records deformation of a Hall or Drinfeld
algebra.  A bridge between them is a specified morphism of formal
deformation bases; a root-of-unity value is a further specialization
of the second base.
\end{remark}

\begin{remark}[Fourier--Jacobi target]
\label{rem:bvbrst-partition-match}
The arithmetic target at two loops is the complete series
\[
 C_2(q,y):=[p^2]\log\Phi_{10}^{\circ}(p,q,y).
\]
Expanding the Gritsenko--Nikulin product determines $C_2$ exactly.
The graph calculation determines $[\mathfrak o_2]$.  The map~$\Psi_2$
identifies their normalizations.
\end{remark}

\section{Roots of the automorphic line and BV counterterms}
\label{sec:bvbrst-supplement-iii}

Let $\lambda$ be the Hodge line on $\mathcal A_2$ and let $\chi$ be
the quadratic character of $\operatorname{Sp}_4(\mathbb Z)$.  The
form $\Delta _5$ is a section of $\lambda^5\otimes\chi$, with divisor
$H_1$ \cite{GN96Igusa}.  The Kummer sequence provides the intrinsic
obstruction to a chosen root of this automorphic line.

\begin{proposition}[Kummer cocycle for an eighth root;
\ClaimStatusProvedHere]
\label{prop:bvbrst-cech-cocycle}
Let $U\subset\mathcal A_2$ be an open substack and
$\{U_i\}$ an étale cover trivializing
$L:=(\lambda^5\otimes\chi)|_U$.  Write $g_{ij}\in
\mathcal O^\times(U_{ij})$ for its transition functions and choose
étale-local eighth roots $r_{ij}^8=g_{ij}$.  Then
\[
 F_{ijk}:=r_{ij}r_{jk}r_{ki}\in\mu_8(U_{ijk})
\]
is a \v Cech $2$-cocycle.  Its class
$[F]\in H^2(U,\mu_8)$ is the Kummer obstruction to an eighth root of
$L$, and $[F]=0$ exactly when such a root exists.
\end{proposition}

\begin{proof}
On quadruple overlaps, cancellation of the $r_{ij}$ gives
$\delta F=1$.  Replacing the local roots changes $F$ by a
$\mu_8$-valued coboundary.  The Kummer exact sequence
$1\to\mu_8\to\mathbb G_m\xrightarrow{(\cdot)^8}\mathbb G_m\to1$
identifies its cohomology class with the connecting image of~$[L]$.
\end{proof}

\begin{conjecture}[Banding-to-counterterm map;
\ClaimStatusConjectured]
\label{prop:bvbrst-banding-counterterm}
For a specified family of twisted BV theories over~$U$, construct a
map
\[
 \beta_{\mathrm{band}}:\ H^2(U,\mu_8)
 \longrightarrow H^0\operatorname{Obs}_{\mathrm{loc}}
      (\mathcal E_{11}^{\mathrm{tw}}/U)\otimes\mathbb Q/\mathbb Z .
\]
The counterterm associated with the automorphic root datum is
\[
 S_{\mathrm{ct}}^{\mathrm{band}}
 :=2\pi i\,\beta_{\mathrm{band}}([F]).
\]
Compatibility with the QME is the equation
$d_0S_{\mathrm{ct}}^{\mathrm{band}}=0$ together with its chosen
integral lift.
\end{conjecture}

\begin{problem}[Meridian realization of the Kummer class;
\ClaimStatusOpen]
\label{rem:bvbrst-linking-number}
Choose a real oriented $2$-cycle~$\Sigma$ in~$U$ and a $3$-chain
bounding it in a compactification.  Compute the intersection number
with the divisor of~$\Delta _5$ and compare
\[
 \exp\!\left(\frac{2\pi i}{8}
       \,\Sigma\cdot\operatorname{div}(\Delta _5)\right)
\]
with the holonomy of~$[F]$ on~$\Sigma$.  This gives a geometric
representative of the Kummer cocycle.
\end{problem}

\begin{remark}[Root-of-unity specialization]
\label{rem:bvbrst-two-hbar-banding}
The class $[F]$ takes values in~$\mu_8$.  A quantum-group
specialization may send its deformation parameter to a primitive
eighth root of unity.  The map $\beta_{\mathrm{band}}$ transports this
finite monodromy to the BV family; it is the precise bridge between
the two parameter spaces.
\end{remark}

\begin{remark}[Sixteenth-root refinement]
\label{rem:bvbrst-mu16-metaplectic-refinement}
The Kummer sequence with $\mu_{16}$ defines the obstruction
\[
 \delta_{16}([L])\in H^2(U,\mu_{16})
\]
to a sixteenth root of~$L$.  A metaplectic refinement consists of a
cover $\widetilde U\to U$ and a line $\widetilde L$ whose sixteenth
power is the pullback of~$L$.  Its square maps to the eighth-root
datum of Proposition~\ref{prop:bvbrst-cech-cocycle}.  Computing
$\delta_{16}([L])$ on meridians gives the refinement problem.
\end{remark}

\section{The Igusa product and the BV comparison problem}
\label{sec:bvbrst-heegner-all-order}

The arithmetic object and the field-theoretic obstruction live in
different categories.  We first state the automorphic theorem, then
formulate the morphism that would compare them.

\begin{theorem}[Gritsenko--Nikulin product for \texorpdfstring{$\Delta _5$}{Delta 5};
\ClaimStatusProvedElsewhere]
\label{thm:bvbrst-heegner-all-order}
\index{Igusa cusp form!Gritsenko--Nikulin product|textbf}
Let
\[
 \phi_{0,1}(\tau,z)=\sum_{n\geq0,\,r\in\bZ}f(4n-r^2)q^ny^r
 =y^{-1}+10+y+O(q)
\]
be the canonical weak Jacobi form of weight~$0$ and index~$1$ in the
Gritsenko--Nikulin normalization.  In the product chamber,
\begin{equation}\label{eq:bvbrst-heegner-identity}
 \Delta _5(p,q,y)
 =(pqy)^{1/2}\prod_{(m,n,r)>0}
       (1-p^mq^ny^r)^{f(4mn-r^2)},
 \qquad \Phi_{10}=\Delta _5^{\,2}.
\end{equation}
Here $(m,n,r)>0$ means $m,n\geq0$, with arbitrary $r$ when
$m+n>0$ and $r<0$ when $m=n=0$.
In particular, $f(0)=10$, $\operatorname{wt}(\Delta _5)=f(0)/2=5$,
and $\operatorname{wt}(\Phi_{10})=10$.
\end{theorem}

\begin{proof}
This is equation~(1.10), with the Jacobi form defined in
equations~(1.8)--(1.9), of Gritsenko--Nikulin
\cite[eqs.~(1.8)--(1.10)]{GN96Igusa}.  Their equation~(1.5) identifies
$\Delta _5^2$ with the Igusa cusp form~$\chi_{10}$.
\end{proof}

The weight-$(-2)$ generator belongs to a second calculation.  With
\[
 \phi_{-2,1}=-(y-2+y^{-1})
 \prod_{j\geq1}\frac{(1-q^jy)^2(1-q^j/y)^2}{(1-q^j)^4},
\]
direct expansion gives coefficients
$a(3)=-8$, $a(4)=12$, $a(7)=-39$, and $a(8)=56$.  The product
\eqref{eq:bvbrst-heegner-identity} uses the coefficients $f(D)$ of
$\phi_{0,1}$.  Thus the two sequences govern distinct Jacobi-form
constructions.

Let $\mathcal E_{11}^{\mathrm{tw}}(Y)$ denote a specified BV model of
the holomorphic--topological twist on an eleven-manifold~$Y$, and put
$d_0=Q+\{I_0,-\}_{\mathrm{BV}}$.  Once
$I_1,\ldots,I_{N-1}$ solve the QME through order $N-1$, its order-$N$
obstruction is the $d_0$-closed local functional
\[
 \mathfrak o_N=
 \Delta_{\mathrm{BV}}I_{N-1}
 +\frac12\sum_{i=1}^{N-1}\{I_i,I_{N-i}\}_{\mathrm{BV}}
 \in Z^1\operatorname{Obs}_{\mathrm{loc}}
       (\mathcal E_{11}^{\mathrm{tw}}(Y)).
\]
Write
\[
 \widetilde\Phi_{10}:=\frac{\Phi_{10}}{pqy},\qquad
 \Phi_{10}^{\circ}(p,q,y):=
 \frac{\widetilde\Phi_{10}(p,q,y)}
      {\widetilde\Phi_{10}(0,q,y)},\qquad
 C_N(q,y):=[p^N]\log\Phi_{10}^{\circ}
 \in\bQ(y)[[q]].
\]
The physics comparison asks for graded maps
\[
 \Psi_N:\ H^1\operatorname{Obs}_{\mathrm{loc}}
       (\mathcal E_{11}^{\mathrm{tw}}(Y))
 \longrightarrow \bQ(y)[[q]],
 \qquad \Psi_N([\mathfrak o_N])=C_N(q,y).
\]
Costello constructs a particular $\Omega$-background reduction of
eleven-dimensional supergravity to five-dimensional non-commutative
gauge theory and proves its perturbative quantization
\cite[Thms.~9.0.2--9.0.3]{Costello2111}.  Raghavendran--Saberi--Williams
construct an interacting twist on $\bR\times Z_5$ for a Calabi--Yau
fivefold~$Z_5$ and describe M2/M5 source backgrounds
\cite{RSW23}.  The maps $\Psi_N$ for a compact K3 geometry form the
remaining holographic construction.

\subsection{Four-loop graph integral}
\label{subsec:bvbrst-4loop-enumeration}

\begin{definition}[Decorated four-loop BV graph]
\label{def:bvbrst-4loop-diagram}
A decorated four-loop graph is a connected graph $\Gamma$ with
$h^1(\Gamma)=4$, vertices labelled by Taylor coefficients of the
chosen interaction~$I_0$, edges labelled by the chosen gauge-fixed
propagator, and divergent subgraphs labelled by the counterterms of
the chosen renormalization scheme.  Its amplitude is the pushforward
of the resulting density on compactified configuration space.
\end{definition}

\begin{table}[ht]
\centering
\begin{tabular}{@{}ll@{}}
\toprule
datum & required construction \\
\midrule
vertices & Taylor coefficients of the interacting twisted action \\
edges & gauge-fixed propagator with parity and orientation \\
counterterms & local extensions along collision strata \\
automorphic output & the map $\Psi_4$ \\
\bottomrule
\end{tabular}
\caption{Data determining the four-loop comparison.}
\label{tab:bvbrst-4loop-enum}
\end{table}

\begin{problem}[Graph-indexing comparison; \ClaimStatusOpen]
\label{prop:bvbrst-kreimer-count}
Construct a bijection between the decorated graphs of
Definition~\ref{def:bvbrst-4loop-diagram} and an external graph-Hopf
algebra in the same valence, tadpole, orientation, and automorphism
convention.  The bijection determines the symmetry factor
$|\operatorname{Aut}\Gamma|^{-1}$ for every field-theoretic graph.
\end{problem}

\begin{problem}[Renormalized K3 kernels; \ClaimStatusOpen]
\label{prop:bvbrst-k3-green-integrals}
Choose a gauge-fixed kinetic operator for
$\mathcal E_{11}^{\mathrm{tw}}(Y)$ and construct its heat-kernel
parametrix on the K3 factor.  For every decorated graph~$\Gamma$, form
the configuration-space integral
\begin{equation}\label{eq:bvbrst-k3-green-powers}
 A_\Gamma^{\mathrm{ren}}
 :=\operatorname{FP}_{\epsilon\to0}
 \int_{\operatorname{Conf}_{V(\Gamma)}(Y)}
 \Bigl(\bigwedge_{e\in E(\Gamma)}P_{\epsilon,e}\Bigr)
 \Bigl(\bigwedge_{v\in V(\Gamma)}I_{0,v}\Bigr)
 +A_\Gamma^{\mathrm{ct}}.
\end{equation}
This construction supplies the K3-dependent amplitude.
\end{problem}

\begin{conjecture}[Four-loop BV/Fourier comparison;
\ClaimStatusConjectured]
\label{thm:bvbrst-4loop-feynman-sum}
For a globally defined twisted theory and a comparison map~$\Psi_4$,
the renormalized order-four obstruction satisfies
\begin{equation}\label{eq:bvbrst-4loop-table}
 \mathfrak o_4=
 \Delta_{\mathrm{BV}}I_3+\{I_1,I_3\}_{\mathrm{BV}}
 +\frac12\{I_2,I_2\}_{\mathrm{BV}},
\end{equation}
and the desired arithmetic comparison is
\begin{equation}\label{eq:bvbrst-4loop-sum}
 \Psi_4([\mathfrak o_4])=C_4(q,y)
 =[p^4]\log\Phi_{10}^{\circ}(p,q,y).
\end{equation}
\end{conjecture}

\begin{remark}[Symmetry factors]
\label{rem:bvbrst-4loop-sym-kreimer}
The factor $|\operatorname{Aut}\Gamma|^{-1}$ is computed in the
groupoid of graphs decorated by the actual field species.  Erasing
the decorations defines a second automorphism group.
\end{remark}

\begin{problem}[Humbert-valued projection; \ClaimStatusOpen]
\label{rem:bvbrst-4loop-humbert-pair}
Construct a morphism from the automorphic target of~$\Psi_4$ to the
Picard group of the chosen compactification of~$\mathcal A_2$, and
compute the divisor of~$C_4$.  Gritsenko--Nikulin identify
$\operatorname{div}(\Delta _5)=H_1$ and
$\operatorname{div}(\chi_{35})=H_1+H_4$
\cite[pp.~3--4]{GN96Igusa}; these are the arithmetic divisor inputs.
\end{problem}

\begin{problem}[Four-loop extension across diagonals;
\ClaimStatusOpen]
\label{rem:bvbrst-4loop-divergence}
Extend every amplitude density from configuration space to its
Fulton--MacPherson compactification and solve the local QME on each
collision stratum.  This determines the finite counterterms and the
renormalized value of each tadpole subgraph.
\end{problem}

\begin{corollary}[Arithmetic coefficient and physical amplitude;
\ClaimStatusConditional]
\label{cor:bvbrst-4loop-cross-verify}
The Gritsenko--Nikulin product computes $C_4$ by formal expansion.
The equality in \eqref{eq:bvbrst-4loop-sum} follows from that arithmetic
identity together with the independently constructed map~$\Psi_4$.
\end{corollary}

\begin{remark}[Vol.~II comparison]
\label{rem:bvbrst-4loop-vol2-companion}
A cross-volume statement requires a morphism from the Vol.~I local
obstruction complex to the Vol.~II Wilson-line deformation complex,
followed by its scalar Fourier--Jacobi trace.  Compatibility with
$\Psi_4$ is a commutative-square condition on these maps.
\end{remark}

\begin{remark}[The coefficients at discriminants five and six]
\label{rem:bvbrst-heegner-c5-c6-vanishing}
The coefficients of $\phi_{-2,1}$ at discriminants five and six
equal zero by the index-one congruence.  The Igusa product is governed by
$\phi_{0,1}$ and its complete coefficients $f(4mn-r^2)$; its
Fourier--Jacobi coefficients $C_5$ and $C_6$ arise from the logarithm
of the full product.
\end{remark}

\begin{remark}[Independent mathematical inputs]
\label{rem:bvbrst-heegner-multipath}
The Jacobi expansion, the Siegel product, and the divisor theorem are
arithmetic statements.  The graph integrals compute classes in local
BV obstruction cohomology.  The maps~$\Psi_N$ compare these two
systems and therefore carry the holographic content.
\end{remark}

\subsection{Diagrammatic comparison maps}
\label{subsec:bvbrst-4loop-moves}

\begin{problem}[Wheel--ladder chain map; \ClaimStatusOpen]
\label{prop:bvbrst-4loop-pair-a}
Construct an orientation-preserving map between the compactified
configuration spaces of the proposed wheel and ladder graphs, pull
back every propagator and vertex density, and compute the determinant
line of the induced map.
\end{problem}

\begin{problem}[Theta--box period map; \ClaimStatusOpen]
\label{prop:bvbrst-4loop-pair-b}
Construct a correspondence between the two graph integrals and a
geometric correspondence between their period targets.  Its action on
$H_1$ and~$H_4$ determines the relative scalar after pushforward.
\end{problem}

\begin{problem}[Triangle--pretzel parity; \ClaimStatusOpen]
\label{prop:bvbrst-4loop-pair-c}
Compute the parity of the propagator determinant line for both graphs
from the BV field degrees and edge ordering.  A sign-reversing chain
map then gives the proposed cancellation.
\end{problem}

\begin{problem}[Renormalized tadpole; \ClaimStatusOpen]
\label{prop:bvbrst-4loop-tadpole-vanish}
For the selected kinetic operator and subtraction scheme, compute
\begin{equation}\label{eq:bvbrst-4loop-tadpole-reg}
 \operatorname{tad}_{\mathrm{ren}}(x)=
 \operatorname{FP}_{\epsilon\to0}
 \int_\epsilon^\infty K_t(x,x)\,dt+c_{\mathrm{ct}}(x).
\end{equation}
The result fixes every graph containing this local subgraph.
\end{problem}

\begin{remark}[Four-loop status]
\label{rem:bvbrst-4loop-summary}
The arithmetic coefficient~$C_4$ is explicit.  The field-theoretic
side consists of the interacting action, decorated graph list,
configuration-space extensions, determinant-line signs, and the map
$\Psi_4$.  These data constitute the four-loop proof problem.
\end{remark}

\subsection{Five-loop comparison}
\label{subsec:bvbrst-5loop-extension}

\begin{definition}[Decorated five-loop BV graph]
\label{def:bvbrst-5loop-diagram}
A decorated five-loop graph is defined as in
Definition~\ref{def:bvbrst-4loop-diagram}, with
$h^1(\Gamma)=5$.  Vertex valences are read from every Taylor
coefficient of the chosen interacting action.
\end{definition}

\begin{problem}[Five-loop graph enumeration; \ClaimStatusOpen]
\label{prop:bvbrst-5loop-kreimer}
Enumerate the decorated five-loop graphs of the twisted theory and
compute their automorphism groups in the same convention as the BV
Feynman rules.
\end{problem}

\begin{table}[ht]
\centering
\begin{tabular}{@{}ll@{}}
\toprule
stage & output \\
\midrule
interaction expansion & decorated vertex types \\
graph generation & isomorphism classes and automorphisms \\
renormalization & finite amplitudes $A_\Gamma^{\mathrm{ren}}$ \\
comparison & $\Psi_5([\mathfrak o_5])$ \\
\bottomrule
\end{tabular}
\caption{Five-loop construction stages.}
\label{tab:bvbrst-5loop-enum}
\end{table}

\begin{conjecture}[Five-loop BV/Fourier comparison;
\ClaimStatusConjectured]
\label{thm:bvbrst-5loop-sum}
For a globally defined twisted theory and comparison map~$\Psi_5$,
\begin{equation}\label{eq:bvbrst-5loop-sum}
 \Psi_5([\mathfrak o_5])=C_5(q,y)
 =[p^5]\log\Phi_{10}^{\circ}(p,q,y).
\end{equation}
\end{conjecture}

\begin{remark}[Five-loop arithmetic target]
\label{rem:bvbrst-5loop-cross-verify}
The product \eqref{eq:bvbrst-heegner-identity} determines $C_5$ as a
formal $q,y$-series.  This coefficient supplies the target of
\eqref{eq:bvbrst-5loop-sum}.
\end{remark}

\begin{remark}[Discriminant and loop order]
\label{rem:bvbrst-4loop-mod4-graph-level}
The discriminant $4mn-r^2$ indexes a Jacobi coefficient inside one
Borcherds factor.  The loop number indexes a graph and the power of
$\hbar_{\mathrm{BV}}$.  A comparison map~$\Psi_N$ provides the
relation between these gradings.
\end{remark}

\begin{remark}[Five-loop cross-volume map]
\label{rem:bvbrst-5loop-cross-volume}
The Vol.~II Wilson-line coefficient and the Vol.~III denominator
coefficient become comparable after constructing typed maps from the
same order-five BV obstruction class to both targets.
\end{remark}

\begin{remark}[Cross-volume arithmetic normalization]
\label{rem:bvbrst-heegner-cross-volume}
Every volume using the $\Delta _5$ denominator must employ the
weight-$0$, index-$1$ coefficients~$f(D)$ of
Theorem~\ref{thm:bvbrst-heegner-all-order}.  The weight-$(-2)$ form
may enter a derivative or insertion construction once that map is
written explicitly.
\end{remark}

\section{Instanton interpretation as a comparison problem}
\label{sec:bvbrst-nonperturbative-d3}

\begin{definition}[Euclidean D3 saddle datum]
\label{def:bvbrst-d3-instanton-amplitude}
A Euclidean D3 saddle consists of a real four-cycle~$W_4$, a classical
solution on~$W_4$, its gauge-fixed deformation complex, and an
orientation of its determinant line.  Its semiclassical contribution
has the form
\begin{equation}\label{eq:bvbrst-d3-amp}
 \mathcal A_{W_4}(\hbar)=
 e^{-S_{\mathrm{cl}}(W_4)/\hbar}
 \operatorname{Pfaff}_{\mathrm{ren}}(D_{W_4})\,
 \bigl(1+O(\hbar)\bigr).
\end{equation}
\end{definition}

\begin{remark}[Worldvolume dimension]
\label{rem:bvbrst-d3-worldvolume-humbert}
A holomorphic curve $C\subset K3$ has real dimension~$2$.  A D3
worldvolume therefore requires an additional real two-cycle or a
different four-cycle construction.  The duality frame must specify
$W_4$, its flux, and its behavior along the Humbert divisor.
\end{remark}

\begin{conjecture}[Borcherds factors from isolated saddles;
\ClaimStatusConjectured]
\label{thm:bvbrst-nonperturbative-completion}
Suppose a duality construction assigns to every positive product
index $(m,n,r)$ an isolated saddle~$W_{m,n,r}$ together with an
equality between its one-loop index and~$2f(4mn-r^2)$.  Then the
saddle product has the formal shape
\begin{equation}\label{eq:bvbrst-nonpert-completion}
 Z_{\mathrm{saddles}}=
 \prod_{(m,n,r)>0}
 \bigl(1-p^mq^ny^r\bigr)^{2f(4mn-r^2)}.
\end{equation}
The arithmetic target is
\begin{equation}\label{eq:bvbrst-phi10-product}
 \Phi_{10}(p,q,y)=pqy
 \prod_{(m,n,r)>0}
 \bigl(1-p^mq^ny^r\bigr)^{2f(4mn-r^2)}.
\end{equation}
\end{conjecture}

\begin{proof}
Equation~\eqref{eq:bvbrst-phi10-product} is the square of
\eqref{eq:bvbrst-heegner-identity}.  The physical statement is the
assignment $(m,n,r)\mapsto W_{m,n,r}$ and the determinant-line
identity in the hypothesis.
\end{proof}

\begin{problem}[Humbert-wall saddle analysis; \ClaimStatusOpen]
\label{rem:bvbrst-schwinger-humbert}
Construct the family of deformation operators~$D_{W_4}$ across a
Humbert wall and compute its spectral flow.  The resulting determinant
line monodromy is the precise wall-crossing datum.
\end{problem}

\begin{remark}[Independent instanton obligations]
\label{rem:bvbrst-d3-multipath}
The cycle assignment, classical action, fluctuation determinant,
and automorphic product are four separate objects.  Agreement follows
from explicit maps among them, evaluated one saddle at a time.
\end{remark}

\begin{remark}[Cross-volume instanton comparison]
\label{rem:bvbrst-nonpert-cross-volume}
The Vol.~II wall-crossing object and the Vol.~III BKM root object are
targets for the cycle assignment of
Theorem~\ref{thm:bvbrst-nonperturbative-completion}.  The cross-volume
statement is the commutativity of the corresponding determinant and
character maps.
\end{remark}

\section{Formal resummation of the arithmetic Fourier tower}
\label{sec:bvbrst-allloop-resummation}

\begin{theorem}[Fourier resummation of the Igusa product;
\ClaimStatusProvedHere]
\label{thm:bvbrst-allloop-resummation}
Let $C_N(q,y)=[p^N]\log\Phi_{10}^{\circ}$.  Then
\begin{equation}\label{eq:bvbrst-allloop-master}
 \exp\!\left(\sum_{N\geq1}
   \hbar_{\mathrm{BV}}^Np^N C_N(q,y)\right)
 =\Phi_{10}^{\circ}(\hbar_{\mathrm{BV}}p,q,y),
\end{equation}
and
\begin{equation}\label{eq:bvbrst-allloop-log}
 \sum_{N\geq1}p^NC_N(q,y)=\log\Phi_{10}^{\circ}(p,q,y).
\end{equation}
The comparison with the BV action is the family of equations
\begin{equation}\label{eq:bvbrst-allloop-effective-action}
 \Psi_N([\mathfrak o_N])=C_N(q,y),\qquad N\geq1.
\end{equation}
\end{theorem}

\begin{proof}
Taking the logarithm of \eqref{eq:bvbrst-phi10-product} gives
\begin{equation}\label{eq:bvbrst-allloop-log-expansion}
 \log\Phi_{10}^{\circ}
 =2\sum_{\substack{(m,n,r)>0\\m>0}}f(4mn-r^2)
      \log(1-p^mq^ny^r).
\end{equation}
Hence
\begin{equation}\label{eq:bvbrst-allloop-log-identity}
 C_N(q,y)=
 -2\!\sum_{\substack{(m,n,r)>0,\,m>0,\,k\geq1\\mk=N}}
 \frac{f(4mn-r^2)}{k}q^{nk}y^{rk}.
\end{equation}
Substitution $p\mapsto\hbar_{\mathrm{BV}}p$ proves
\eqref{eq:bvbrst-allloop-master}; the specialization
$\hbar_{\mathrm{BV}}=1$ proves \eqref{eq:bvbrst-allloop-log}.
\end{proof}

\begin{remark}[Analytic domain]
\label{rem:bvbrst-allloop-gevrey-scope}
Gritsenko--Nikulin prove absolute convergence of the product in a
Siegel product chamber.  Formal resummation uses the
$p$-adic topology; analytic continuation uses the automorphic form
$\Phi_{10}$.  A semiclassical expansion in $\hbar_{\mathrm{BV}}$
requires the maps~$\Psi_N$ and analytic estimates for the corresponding
graph integrals.
\end{remark}

\begin{remark}[Arithmetic exactness and physical interpretation]
\label{rem:bvbrst-allloop-nonpert-exact}
Equation~\eqref{eq:bvbrst-phi10-product} is an exact automorphic
identity.  Its instanton interpretation is the cycle-and-determinant
construction of
Theorem~\ref{thm:bvbrst-nonperturbative-completion}.
\end{remark}

\begin{conjecture}[Bulk and boundary character maps;
\ClaimStatusConjectured]
\label{cor:bvbrst-allloop-bulk-boundary}
Suppose a renormalized BV trace and a boundary character map satisfy
\begin{equation}\label{eq:bvbrst-allloop-effective}
 \operatorname{Tr}_{\mathrm{BV}}(I_{\mathrm{eff}})
 =\log\Phi_{10}^{\circ},
\end{equation}
\begin{equation}\label{eq:bvbrst-allloop-zbulk}
 Z_{\mathrm{bulk}}^{\mathrm{arith}}=\Phi_{10}^{\circ},
\end{equation}
and
\begin{equation}\label{eq:bvbrst-allloop-zboundary}
 \operatorname{ch}_{\partial}(\mathcal H_{\partial})
 =Z_{\mathrm{boundary}}^{\mathrm{arith}}
 =(\Phi_{10}^{\circ})^{-1}.
\end{equation}
Then
\begin{equation}\label{eq:bvbrst-allloop-bulk-boundary-product}
 Z_{\mathrm{bulk}}^{\mathrm{arith}}
 Z_{\mathrm{boundary}}^{\mathrm{arith}}=1,
\end{equation}
and in particular
\begin{equation}\label{eq:bvbrst-allloop-lusztig-duality}
 \left.Z_{\mathrm{bulk}}^{\mathrm{arith}}
 Z_{\mathrm{boundary}}^{\mathrm{arith}}
 \right|_{\hbar_{\mathrm{BV}}=1}=1.
\end{equation}
\end{conjecture}

\begin{proof}
Multiply equations~\eqref{eq:bvbrst-allloop-zbulk} and
\eqref{eq:bvbrst-allloop-zboundary}.  The field-theoretic content is
the construction of the two maps in the hypothesis.
\end{proof}

\begin{remark}[Arithmetic and field-theoretic routes]
\label{rem:bvbrst-allloop-multipath}
The product formula and its logarithmic expansion are two forms of one
arithmetic theorem.  The BV graph expansion is a second object.  The
comparison maps $\Psi_N$, $\operatorname{Tr}_{\mathrm{BV}}$, and
$\operatorname{ch}_{\partial}$ carry the relations between them.
\end{remark}

\begin{remark}[Cross-volume formulation]
\label{rem:bvbrst-allloop-cross-volume}
Vol.~I supplies local BV obstruction classes, Vol.~II supplies
boundary and line deformation objects, and Vol.~III supplies the
automorphic denominator.  A cross-volume theorem consists of typed
maps from the first two objects to the third, compatible with the
Fourier filtration and with
Theorem~\ref{thm:bvbrst-allloop-resummation}.
\end{remark}
